%=================================================================
%  Supplementary Materials for
%  "Composition-Level Constraints in Alternative Gravity:
%   A Recursive Constraint-Transport Criterion"
%  Supplementary file for Annales de la Fondation Louis de Broglie submission.
%=================================================================
\documentclass[11pt]{article}

\usepackage[utf8]{inputenc}
\usepackage[T1]{fontenc}
\usepackage{amsmath,amssymb,amsthm}
\usepackage{enumitem}
\usepackage{hyperref}
\hypersetup{hidelinks}
\usepackage[margin=1in]{geometry}

\title{Supplementary Material: Composition-Level Constraints in Alternative Gravity\\
--- A Second-Order Decision Criterion}

\author{Andy E. Williams\thanks{Correspondence: info@cc4ci.org}\\[4pt]
\small Caribbean Center for Collective Intelligence, St.~John's, Antigua and Barbuda\\
\small Nobeah Foundation, Nairobi, Kenya}

\date{}

%================== macros ==================
\newcommand{\Pspace}{\ensuremath{\mathcal{P}}}
\newcommand{\ARspace}{\ensuremath{\mathcal{A}_{R}}}
\newcommand{\Tspace}{\ensuremath{\mathcal{T}}}
\newcommand{\SOP}{\ensuremath{\mathrm{SOP}}}
\newcommand{\FOP}{\ensuremath{\mathrm{FOP}}}
\newcommand{\Reg}{\ensuremath{\mathsf{Reg}}}
\newcommand{\SEE}{\ensuremath{S_{\mathrm{EE}}}}
\newcommand{\Adm}{\operatorname{Adm}}
\newcommand{\Gobj}{\ensuremath{\mathbb{G}}}
\newcommand{\Audit}{\ensuremath{\mathsf{Audit}}}
\newcommand{\View}{\ensuremath{\mathsf{View}}}
\newcommand{\Emb}{\ensuremath{\mathsf{Emb}}}
\newcommand{\Tr}{\ensuremath{\mathsf{Tr}}}
\newcommand{\Inv}{\ensuremath{\mathsf{Inv}}}
\newcommand{\Register}{\ensuremath{\operatorname{Register}}}
\newcommand{\Recall}{\ensuremath{\operatorname{Recall}}}
\newcommand{\GOp}{\ensuremath{\mathsf{G}}}
\newcommand{\LOp}{\ensuremath{\mathsf{L}}}
\newcommand{\OpBasis}{\ensuremath{\mathsf{Op}}}
\newcommand{\CorrSig}{\ensuremath{\mathsf{Sig}_{\mathrm{Corr}}}}
\newcommand{\GenClos}{\ensuremath{\mathsf{Gen}_{\GOp}}}
\newcommand{\NovSig}{\ensuremath{\mathsf{Sig}^{\mathrm{Nov}}}}
\newcommand{\NovRes}{\ensuremath{\mathsf{Res}^{\mathrm{nov}}}}
\newcommand{\Sched}{\ensuremath{\mathsf{Sched}}}
\newcommand{\Supp}{\ensuremath{\mathsf{Supp}}}
\newcommand{\Top}{\ensuremath{\mathsf{Top}}}
\newcommand{\Role}{\ensuremath{\mathsf{Role}}}
\newcommand{\ObsMap}{\ensuremath{\mathsf{Obs}}}
\newcommand{\Bridge}{\ensuremath{\mathfrak{B}}}
\newcommand{\Comp}{\ensuremath{\mathsf{Comp}}}
\newcommand{\SyncRes}{\ensuremath{\rho_{\Sigma}}}
\newcommand{\Novel}{\ensuremath{\mathsf{Novel}}}
\newcommand{\GPair}{\ensuremath{\mathsf{G}^{P/A}}}
\newcommand{\ObsAdeq}{\ensuremath{\mathsf{ObsAdeq}}}
\newcommand{\ObsNeed}{\ensuremath{\mathsf{Need}^{\mathrm{obs}}}}
\newcommand{\ObsRes}{\ensuremath{\mathsf{Res}^{\mathrm{obs}}}}
\newcommand{\GCpass}{\textsc{[globally coherent]}}
\newcommand{\GCpending}{\textsc{[certificate pending]}}
\newcommand{\GCcert}{\ensuremath{\mathsf{GCert}}}
\newcommand{\GCverify}{\ensuremath{\mathsf{Verify}}}
\newcommand{\GainCoord}{\ensuremath{\mathsf{Gain}}}
\newcommand{\RCert}{\ensuremath{\mathsf{RCert}}}

% claim-type tags
\newcommand{\Defn}{\textsc{[def]}}
\newcommand{\Asm}{\textsc{[asm]}}
\newcommand{\Thm}{\textsc{[thm]}}
\newcommand{\Cnd}{\textsc{[cnd]}}
\newcommand{\Emp}{\textsc{[emp]}}
\newcommand{\Imp}{\textsc{[imported]}}
\newcommand{\Hyp}{\textsc{[hyp]}}
\newcommand{\Open}{\textsc{[open]}}

\theoremstyle{definition}
\newtheorem{sdefinition}{Definition}
\newtheorem{sproposition}{Proposition}
\newtheorem{sobligation}{Propagation Obligation}
\newtheorem{scard}{Imported-Result Card}
\newtheorem{compositioncard}{Source-to-Composition Card}
\newtheorem{shypothesis}[sdefinition]{Hypothesis}
\newtheorem{saxiom}[sdefinition]{Axiom}
\newtheorem{stheorem}[sdefinition]{Theorem}
\newtheorem{slemma}[sdefinition]{Lemma}
\newtheorem{scorollary}[sdefinition]{Corollary}

\begin{document}

\maketitle

\section*{Reading guide and claim-type discipline}
\addcontentsline{toc}{section}{Reading guide and claim-type discipline}

The main article derives $\GOp$ from that question. For a target $Q$, an
observation is adequate when $Q$ factors through a typed P-space to A(R)-space
observation transformation. The least non-erasing coupled closure that carries the
relations required by that factorization is the paired operator
$\GPair_{Q,\delta}$; its P-space and A(R)-space projections are the corresponding
single-space $G$ closures. Section~\ref{supp:observability-induced-closure} gives
the formal derivation before the gravitational path uses the operator.

The main article asks a simple question: can a material observation record erase a
relation needed to decide the physical target? The simplest grammar is recursive.
Every physical state is a potentially recursively composed physical-state node, and
every physical process is a potentially recursively composed edge or path. Every
observation is a potentially recursively composed registration node, and every
process of observation is a potentially recursively composed registration edge or
path. This supplement supplies the formal route through that picture. It first
states the two networks and the bridge between their interaction roles, then the
matter-observation composite and the residue that a local observation quotient can
erase. The later sections give the finite witness, the worked inverse-gravity
exhibit, and the recursive account of why a bounded local verdict and an unbounded
generated closure can both be valid assessments of different objects.

Theorem-level claims are bounded under $L$: they are assessed within a declared
target, finite record, and coherence budget. Global novelty claims are potentially
unbounded under a productive $G$-closure: their object is the recursively generated
set of theory, test, bridge, registration, and revision objects rather than a
single finite theorem record. Section~\ref{supp:assessment-relative} formalizes
when a limitation visible in a contracted record persists as a genuine obstruction
to this larger closure.

The paper makes the first finite portion of that closure visible rather than asking
an assessor to infer it from a single theorem. The path begins with existing
interaction, observation, and gravity hypotheses; passes through paired topology,
the four-role interaction hypothesis, a typed bridge, a synchronization residue,
and the blind-spot criterion; then generates a target-conditioned inverse-gravity
composition and a second-order theory of gravitation. Section~\ref{supp:generated-path}
states the stages, their inputs, their local projections, and the asymptotic
reachability condition that distinguishes a generated realization path from a list
of unrelated premises.

Readers seeking the shortest formal route may read Sections~\ref{supp:spaces},
\ref{supp:topological-bridge}, and~\ref{supp:matter-blindspot} before turning
to the calculus in Section~\ref{supp:calculus} and the strict-separation proof in
Section~\ref{supp:strict-separation}. Readers assessing the novelty claim may then
read Section~\ref{supp:novelty}. The remaining sections preserve the full
submission package: imported-result cards, the worked case, audit infrastructure,
and conditional temporal transport.

Every numbered statement is tagged by logical form: \Defn{} a definition;
\Asm{} a constitutive assumption of a declared model class; \Thm{} a theorem
relative to the stated definitions and assumptions; \Cnd{} a conditional
proposition; \Emp{} an empirical proposition; and \Imp{} a result imported from
the external literature and used as a cited source object. These form tags do not
replace the composition verdict. Each generated object carries the three-axis status
record
\begin{equation}
\mathsf{Status}_{B}(H)=
\bigl(\mathsf{Form}(H),\ \mathsf{GC}_{B}(H),\ \mathsf{Cont}_{B}(H)\bigr),
\label{eq:three-axis-status}
\end{equation}
where $\mathsf{Form}$ gives its logical form, $\mathsf{GC}_{B}$ gives its
budgeted global-coherence verdict, and $\mathsf{Cont}_{B}$ records the continuation,
realization, or typed defeat frontier. A source citation establishes its native
carrier, lawful relation, regime, observable, constraint, derivation, or empirical
registration. The manuscript records the explicit new bridge by which source
objects are composed. A hypothesis graph is globally coherent relative to a stated
assessment when its source objects, type assignments, bridges, premises, admissible
branches, and failure modes are explicit, a finite certificate records the checks
actually performed, and no typed contradiction has been established within that
certificate's budget. Thus a conditional relation or a hypothesis-valued coordinate
can be internal to a composition whose emitted verdict is \GCpass{}; its continuation
record specifies the next realization and quantitative-instantiation tasks internal
to that coherent structure. A form tag is never itself a global-coherence verdict.

\begin{sdefinition}[Assessment-bounded global-coherence verdict]
\label{def:assessment-budget}
\Defn{} Let $H$ be a typed hypothesis graph and $B$ a declared coherence-assessment
budget, specifying the branches, imports, bridge obligations, and finite tests
examined. Write
\begin{equation}
\mathsf{GC}_{B}(H)=\mathsf{pass}
\end{equation}
only when an exhibited certificate $\GCcert_B(H)$ records the tested claim scope,
imports, carrier assignments, bridge checks, regime checks, non-erasure checks,
defeat audit, and residue, and its stated verification returns \ensuremath{\checkmark}.
The verdict is a status-labeled in-budget result: its examined branches, imports,
bridge obligations, and finite tests are all recorded in $B$, and every unexamined
branch remains explicit residue. Before that verification, the appropriate
composition coordinate is \GCpending{}, not \GCpass{}.
\end{sdefinition}

\begin{sdefinition}[Finite global-coherence certificate and verification]
\label{def:global-coherence-certificate}
\Defn{} For a typed hypothesis graph $H$ and budget $B$, a finite
\emph{global-coherence certificate} is
\begin{equation}
\begin{aligned}
\GCcert_B(H)=\bigl(&\mathsf{ClaimScope},\mathsf{Imports},\mathsf{CarrierCheck},\\
&\mathsf{BridgeCheck},\mathsf{RegimeCheck},\mathsf{NonErasureCheck},\\
&\mathsf{DefeatAudit},\mathsf{Residue}\bigr)_B.
\end{aligned}
\label{eq:global-coherence-certificate}
\end{equation}
$\GCverify_B\bigl(\GCcert_B(H)\bigr)=\checkmark$ exactly when each listed
check is satisfied for the claims actually made by $H$ and the audit contains no
in-budget typed defeater. The verification neither instantiates a continuation
branch nor converts a conditional consequence into an empirical assertion; it
verifies that these remain correctly typed, non-erased records of $H$.
\end{sdefinition}

\begin{sdefinition}[Residue and typed defeater]
\label{def:residue-defeater}
\Defn{} The \emph{assessment residue} $\mathsf{Res}_{B}(H)$ is the set of
branches, bridge obligations, realization conditions, and possible countermodels
not adjudicated within $B$. Residue is retained with its declared status (for
example open, bridge-required, regime-limited, or empirical) and is not a defect
by default. A \emph{typed defeater} is a finite in-budget argument that, under the
declared semantics, supplies a contradiction, shows a witness ill-formed, or shows
that a conclusion already claimed requires an omitted obligation. Merely naming a
possible unassessed branch, alternative realization, or future certificate adds to
residue unless it establishes one of those finite failures.
\end{sdefinition}

\begin{sdefinition}[Source-to-composition record]
\label{def:source-to-composition-record}
\Defn{} A source-to-composition record for a generated stage $G_i:S_{i-1}\to S_i$
is
\begin{equation}
\mathsf{SC}_i=
\left(
\mathsf{Src}_i,\mathsf{Native}_i,\mathsf{Inv}_i,\mathsf{Bridge}_i,
\mathsf{Out}_i,\mathsf{Status}_{B,i},\mathsf{Defeat}_i
\right).
\label{eq:source-to-composition-record}
\end{equation}
Here $\mathsf{Src}_i$ is the cited set of pre-existing physical, mathematical, or
methodological objects; $\mathsf{Native}_i$ records their carriers, laws, regimes,
observables, derivations, and empirical registrations; $\mathsf{Inv}_i$ is the
target-relevant invariant retained from those objects; $\mathsf{Bridge}_i$ is the
new typed composition asserted in this manuscript; $\mathsf{Out}_i$ is the emitted
object; $\mathsf{Status}_{B,i}$ is the three-axis record of
Equation~\eqref{eq:three-axis-status}; and $\mathsf{Defeat}_i$ records a
contradiction, failed certificate, or other typed failure. When the source records and bridge obligations in the declared budget are
satisfied and the finite certificate of Definition~\ref{def:global-coherence-certificate}
verifies, their composition emits $\mathsf{GC}_{B}(H)=\mathsf{pass}$ for the
status-labeled hypothesis. Open entries then specify the next realization and
quantitative-instantiation tasks internal to that verdict. No entry substitutes a
source object for a bridge or a bridge for a source object.
\end{sdefinition}

\paragraph{Cited and re-theoremed antecedents.}
The standard parity/proper-support witness, quotient-factorization criterion, and
least-fixed-point theorem used below are cited as antecedents. They are also
re-stated and proved in the present typed setting. The re-theorems do not claim
original discovery of those mathematical facts; they identify the exact P-space,
A(R)-space, bridge, target, and residue relations through which each antecedent
functions in the composed physical assessment.

\begin{sdefinition}[Generated registration closure]
\label{def:generated-registration-closure}
\Defn{} Let $R$ be a recursive constraint registration carrying any typed,
registrable payload $p$, and let $\Theta$ specify a declared assessment target.
The \emph{$\GOp$-generated registration closure} is
\begin{equation}
\GenClos(R;\Theta)=
\left\{R'\;\middle|\;R'\text{ is emitted from }R\text{ by a finite or
 guarded-recursive typed }\GOp\text{-program directed to }\Theta\right\}.
\label{eq:generated-registration-closure}
\end{equation}
The closure includes $R$ and may include derived theorem registrations, typed
instances, target predicates, preservation signatures, witnesses, countermodels,
provenance records, and transfer obligations. There is no source-domain restriction
on $p$: a physical model, existing mathematical theorem, proof object, or prior
registration may be its payload. ``Emitted'' means introduced as a typed
registration with declared derivational and provenance relations. For a coupled
P-space/A(R)-space record, this registration closure is the A(R)-space projection
of the observability-induced paired closure derived in
Section~\ref{supp:observability-induced-closure}; that derivation supplies the
condition under which the registration path decides the physical target.
\end{sdefinition}

\begin{sdefinition}[Budget-relative novelty signature and residue]
\label{def:novelty-residue}
\Defn{} Fix a comparison registration $K_B$ containing the relations examined
under a declared coherence budget $B$. For target $\Theta$, the
\emph{novelty signature} of $R$ is the target-relevant tuple
\begin{equation}
\NovSig_{\Theta}(R)=
\left(
 p,\GenClos(R;\Theta),\Theta,
 \mathsf{Pred}_{\Theta},\mathsf{Witness}_{\Theta},
 \mathsf{Transfer}_{\Theta},\mathsf{Prov}
\right),
\label{eq:novelty-signature}
\end{equation}
where $p$ is the antecedent payload, $\mathsf{Pred}_{\Theta}$ the declared target
predicate, $\mathsf{Witness}_{\Theta}$ the generated witness or countermodel
relations, $\mathsf{Transfer}_{\Theta}$ the generated constraint-transfer
relations, and $\mathsf{Prov}$ the associated provenance. A licensed novelty
assessment is an $\LOp$-program with a certificate that its comparison preserves
$\NovSig_{\Theta}$ at budget $B$. It yields a status
\begin{equation}
\nu_{B,\Theta}(R)\in
\{\mathsf{accounted},\mathsf{nonzero\ residue},\mathsf{unresolved}\}
\quad\text{and}\quad
\NovRes_{B,\Theta}(R)
 =\GenClos(R;\Theta)\setminus\mathsf{Account}_{B,\Theta}(K_B).
\label{eq:novelty-residue}
\end{equation}
Here $\mathsf{Account}_{B,\Theta}(K_B)$ is the subset for which the comparison
record supplies a typed relation preserving the declared signature. A familiar
antecedent payload can therefore be accounted for while a nonempty generated
closure remains novelty residue. A finite $\LOp$ verdict is not an assertion that
novelty is exhausted outside its stated budget and target.
\end{sdefinition}

\begin{sproposition}[Familiar-form recognition is not a global novelty verdict]
\label{prop:familiar-form-not-novelty-verdict}
\Thm{} Let $R$ carry a theorem or other typed payload that a comparison record
$K_B$ recognizes as familiar. That recognition alone does not establish
$\NovRes_{B,\Theta}(R)=\varnothing$, nor does it establish an unbounded verdict
of no novelty. Such a conclusion is licensed only by an $\LOp$-program whose
certificate preserves $\NovSig_{\Theta}(R)$ over the generated registration
closure.
\end{sproposition}

\begin{proof}
Recognition of the antecedent payload establishes, at most, that one component of
$\NovSig_{\Theta}(R)$ is accounted for by $K_B$. By
Definition~\ref{def:novelty-residue}, all generated target, witness, transfer, and
provenance relations remain in $\NovRes_{B,\Theta}(R)$ unless an admissible
comparison preserves them. An $\LOp$ operation without that certificate identifies
records that may differ on the declared novelty signature and is therefore not a
licensed novelty assessment. The final clause follows because $B$ is finite and
all unassessed relations remain status-labeled residue.
\end{proof}

\paragraph{Assessment sequence.}
The formal separation result, the physical status of imported bridges, and the
empirical reference of the worked case must be assessed separately. A finite
defeater of the stated witness can defeat the separation claim. A missing bridge,
invalid regime, or unresolved safe witness instead routes the physical branch to
its declared status unless it also supplies such a defeater. This sequence prevents
an unresolved realization question from being substituted for a verdict on the
typed decision result.


The organizing object of this supplement is the \emph{decision record}. A
physical theory supplies a path of physical operations. The formal layer supplies
the registration through which that path, its bridges, regimes, support,
certificates, provenance, and status are assessed. The reader can therefore keep
one question in view throughout: after a proposed local simplification, does the
record still contain every relation on which the target depends?

The sections are arranged by function rather than by rhetorical priority. The
physical-state and registration-space definitions explain the objects; the
topological bridge and matter-observation sections explain the blind-spot model;
the calculus specifies admissible expansion and contraction; the strict witness
proves the finite separation result; and the later sections preserve the worked
case, audit, temporal, and successor-transport infrastructure. Every physical
realization claim remains attached to its bridge, regime, witness, and empirical
status rather than being silently promoted into a theorem about the decision
record.

\tableofcontents

\noindent\textit{Navigation.} The main article gives the physical statement first.
This supplement gives the compact formal answer: an observation is a recursively
composed registration node, an observation process is a recursively composed edge
or path, and a blind spot is a target-changing relation erased by the observer's
observation topology. The local record is adequate only when the target is constant
on the map's fibers. The other constructions expand, transport, test, or audit that
same decision condition.


%=================================================================
\section{The formal decision layer: what the record must retain}
\label{supp:calculus}

The finite corridor theorem specifies the assessment representation that
accompanies a physical theory. Let a candidate physical path be
\begin{equation}
\gamma=
\left(
 x_0\xrightarrow{u_1}x_1\xrightarrow{u_2}\cdots
 \xrightarrow{u_n}x_n
\right),
\qquad u_k\in\mathcal O_P.
\label{eq:physical-path-supp}
\end{equation}
The operators $u_k$ are the theory-relative physical operations of the path. The
formal operators below govern the record through which their admissibility, bridge
conditions, support, provenance, and status are transported.

\begin{sdefinition}[Recursive constraint registration]
\label{def:recursive-constraint-registration}
\Defn{} For a declared candidate composition $\mathfrak C$, a
\emph{recursive constraint registration} is a finite or guarded-recursive typed
record
\begin{equation}
R_{\mathfrak C}=
\left\langle
R_{\Pspace},R_{\ARspace},R_{\tau},R_{\mathcal B},R_{\gamma},
R_{\{\mathcal R_k\}},R_S,R_{\mathsf{SafeWit}},
R_{\mathsf{Cert}},R_{\mathsf{Status}},R_{\mathsf{Prov}}
\right\rangle .
\label{eq:recursive-constraint-registration}
\end{equation}
The components register, respectively, physical and registration representations,
typing, bridges, ordered physical path, regime records, support, safe-witness
state, certificates, current status, and provenance. Any component may carry a
prior registration and its provenance. Physical states and operations appear here
as declared payloads with their own types; they are not identified with the
registration operations in Definition~\ref{def:constraint-program}.
\end{sdefinition}

\begin{sdefinition}[Constraint-transport program]
\label{def:constraint-program}
\Defn{} Let
\begin{equation}
\OpBasis=\{\Register,\Recall,\GOp,\LOp\}.
\label{eq:constraint-operator-basis}
\end{equation}
A \emph{constraint-transport program} is a finite or guarded-recursive typed
program in $\langle\OpBasis\rangle_{\mathrm{rec}}$ acting on recursive constraint
registrations. $\Register$ introduces a typed datum, certificate, or status;
$\Recall$ recovers a datum together with the provenance and bridge conditions that
license its use; $\GOp$ non-erasively extends the active registration and may emit
new typed registrations derived from any registrable payload when a relation is not
determined by the current record; and $\LOp$ contracts the active registration
only under a target-preservation certificate as defined below.
\end{sdefinition}

\begin{sdefinition}[Corridor signature, licensed contraction, and required expansion]
\label{def:licensed-contraction}
\Defn{} For a corridor assessment, let
\begin{equation}
\CorrSig(\mathfrak C)=
\left(
\mathsf{Bridge}(\mathfrak C),
\mathsf{Regime}(\mathfrak C),
\mathsf{History}(\mathfrak C),
\mathsf{Support}(\mathfrak C),
\mathsf{SafeWit}(\mathfrak C),
\mathsf{Status}(\mathfrak C)
\right).
\label{eq:corridor-signature}
\end{equation}
A contraction $\LOp(R_{\mathfrak C};c)$ is \emph{licensed} for the declared
corridor target only when certificate $c$ establishes that every pair of active
records identified by the contraction has the same corridor value and the same
declared decision-relevant status. A $\GOp$-expansion is required when the current
record identifies a pair whose corridor value or target-relevant status can differ,
unless a new barrier, safe-witness, or other certificate changes the decision
branch without making that identification.
\end{sdefinition}

\begin{sproposition}[One closed assessment grammar]
\label{prop:one-assessment-grammar}
\Thm{} Every active second-order assessment object in this supplement---a typed
path record, bridge, regime certificate, support witness, audit quotient, residue,
trace record, or successor safety case---is a recursive constraint registration or
a declared payload carried by one. Every admissible second-order assessment
transition is a constraint-transport program.
\end{sproposition}

\begin{proof}
The claim is a closure declaration for the assessment representation. Each active
object is introduced as a typed component of a registration, a relation carried by
one, or a certificate referring to a preceding registration. Each subsequent
assessment action introduces, recovers, expands, or conditionally contracts such a
record and is therefore typed by Definition~\ref{def:constraint-program}. Since a
payload may be any typed registrable object, this is a schema-level closure across
branches of theory: an existing theorem, model, proof object, or prior registration
can be expanded by $\GOp$ into further registrations without adding a fifth
primitive assessment operator. The claim does not classify the physical $u_k$ as
$\Register$, $\Recall$, $\GOp$, or $\LOp$; their physical action is carried as
payload in $R_{\gamma}$.
\end{proof}

\begin{sproposition}[Local quotient as a candidate $\LOp$-program]
\label{prop:local-quotient-lop}
\Thm{} A local-observation quotient is a licensed $\LOp$-contraction for
corridor status if and only if corridor status is constant on every fiber of that
quotient. When this condition fails, an $\LOp$-only assessment cannot decide the
corridor predicate; the assessment must expand the record by $\GOp$ or retain the
branch as unresolved pending an adequate certificate.
\end{sproposition}

\begin{proof}
This is the program interpretation of the fiber criterion proved in
Proposition~\ref{prop:fiber-criterion}. The certificate condition in
Definition~\ref{def:licensed-contraction} is exactly constancy of the target on
contraction fibers. If it fails, a local program identifies instances that require
different decisions, so it cannot be sound and complete for that target.
\end{proof}

\begin{sdefinition}[Directional assessment]
\label{def:directional-assessment}
\Defn{} For a declared target $\Theta$, let $R^{(2)}_{\Theta}$ be the fully typed
second-order assessment record, retaining the physical trajectory, registration
trajectory, bridges, history, support, certificates, and target predicate. A
corresponding first-order assessment record has the form
\begin{equation}
R^{(1)}_{\Theta}=q_{\Theta}\bigl(R^{(2)}_{\Theta}\bigr),
\label{eq:directional-assessment}
\end{equation}
where $q_{\Theta}$ is a declared candidate $\LOp$-contraction. Assessment is
\emph{directional}: the second-order record is constructed first, and the
first-order record is derived from it. The contraction declaration states which
relations have been omitted from the derived record.
\end{sdefinition}

\begin{sproposition}[Verdict-changing contraction obligation]
\label{prop:verdict-changing-contraction}
\Thm{} Under Definition~\ref{def:directional-assessment}, a first-order verdict
may stand for the second-order target verdict only when $q_{\Theta}$ is licensed
for $\Theta$. If the two verdicts differ, the contraction must retain an explicit
record of an erased relation on which the target changes; without a certificate
showing that this relation is target-invariant, the first-order verdict cannot
close the second-order assessment.
\end{sproposition}

\begin{proof}
By Definition~\ref{def:licensed-contraction}, licensing requires constancy of the
target on every contraction fiber. A difference between the two verdicts exhibits a
fiber containing records with different target values. The contraction has therefore
erased a target-changing relation and is unlicensed unless a certificate changes
that condition. Proposition~\ref{prop:local-quotient-lop} gives the same result
for a local-observation quotient.
\end{proof}

%=================================================================
\section{Theorem-level bounds and global generated novelty}
\label{supp:novelty}

\paragraph{The simple distinction.}
Theorem-level claims are bounded under $L$: a finite contraction is assessed only
for the declared target it preserves. Global novelty claims are potentially
unbounded under $G$: a productive recursive closure can emit a non-equivalent
theory, test, bridge, registration, or revision object beyond every finite
comparison record. A representation can therefore leave every answer unchanged for
a chosen local task and still generate new objects when it is allowed to operate
recursively. This section separates those two assessments before stating their
formal conditions.

The preceding definitions make precise a distinction that is often lost when a
paper is compared with a familiar antecedent theorem. The antecedent theorem is a
registrable payload. Its recognition can be an accurate relation in $K_B$, but it
is not automatically a contraction certificate for the target-specific
registrations emitted when that theorem is deployed. A novelty verdict must
therefore follow the same discipline as any other adequacy verdict: first expand
what the target requires, then contract only when the declared signature is
preserved.

\begin{sproposition}[Novelty assessment is a constraint-transport program]
\label{prop:novelty-transport-program}
\Thm{} For a declared target $\Theta$, a valid novelty assessment of a typed
payload $p$ has the program form
\begin{equation}
R_p
\xrightarrow{\ \Register/\Recall\ }
\GenClos(R_p;\Theta)
\xrightarrow{\ \LOp_{B,\Theta}\ }
\left(\nu_{B,\Theta}(R_p),\NovRes_{B,\Theta}(R_p)\right),
\label{eq:novelty-transport-program}
\end{equation}
where the final contraction is licensed only if it preserves
$\NovSig_{\Theta}(R_p)$. Recognition of an antecedent theorem is one possible
input relation to this program, not a replacement for it.
\end{sproposition}

\begin{proof}
The first transition is a finite or guarded-recursive $\GOp$ expansion as in
Definition~\ref{def:generated-registration-closure}. The second transition is a
licensed $\LOp$ assessment precisely when it preserves the novelty signature of
Definition~\ref{def:novelty-residue}. Proposition~\ref{prop:familiar-form-not-novelty-verdict}
excludes treating antecedent recognition alone as the final verdict.
\end{proof}

\paragraph{Application to the present result.}
The familiar factorization form of the fiber criterion is an antecedent payload.
The generated deployment considered here includes the typed corridor predicate,
the proper-support parity witness, the four-coordinate family, the identification
of a local observation quotient as a candidate $\LOp$ contraction, and the
constraint-transport obligations that follow if that contraction is unlicensed. It
also includes the visible finite path of Section~\ref{supp:generated-path}: the
four-role interaction hypothesis, the typed bridge and synchronization residue, the
target-conditioned inverse-gravity composition, and the emitted second-order theory
of gravitation. The quotient step is one transition in that path, not its terminal
node. Whether a comparison record already contains all of those relations is a
budget-indexed question. The present article does not assign an unbounded novelty
verdict by familiar-form recognition alone. It supplies the distinct productivity
condition under which a $\GOp$-generated closure has unbounded generative capacity.
A reduction of this deployment to ``only a quotient fact'' is licensed only when
the reduction preserves the declared novelty signature and leaves no relevant
residue at the stated budget.


\subsection{A bounded local verdict and a generated global closure}
\label{supp:bookkeeping-generated-closure}

The phrase ``only bookkeeping'' has a precise, limited interpretation in this
calculus. It is a verdict about a declared finite task after a declared
contraction; it is not automatically a verdict about the recursively generated
closure of the representation.

\begin{sdefinition}[Bounded bookkeeping verdict]
\label{def:bounded-bookkeeping}
\Defn{} Let $\mathcal K$ be a typed physics record and let $D_B$ be a decision,
comparison, or explanatory task under coherence budget $B$. A \emph{bounded
bookkeeping verdict} is a licensed $\LOp$ contraction
\begin{equation}
\LOp_B:\mathcal K\longrightarrow\mathcal K_B
\label{eq:bounded-bookkeeping-contraction}
\end{equation}
for which there exists $\overline D_B$ satisfying
\begin{equation}
D_B=\overline D_B\circ\LOp_B.
\label{eq:bounded-bookkeeping-factorization}
\end{equation}
The verdict means that no distinction discarded by $\LOp_B$ changes $D_B$ within
$B$. It does not assert that all relations emitted by subsequent admissible
$\GOp$ programs factor through the same contraction.
\end{sdefinition}

\begin{sdefinition}[Tri-space generated physics closure]
\label{def:tri-space-generated-closure}
\Defn{} Let
\begin{equation}
\mathcal K^{(0)}=
\left(
\Pspace_{\mathcal E;R,\delta},
(\ARspace)_{\mathcal E;R,\delta},
\Tspace_{\mathcal E,\delta},B,\mathsf{Status},\mathsf{Prov}
\right)
\label{eq:tri-space-seed}
\end{equation}
be a typed physical, registration, and temporal-support record. Its one-step
generated extension is
\begin{equation}
\mathsf{Gen}(\mathcal K)=
\mathcal K\cup\GOp(\mathcal K)\cup\Comp_{P,A,T}(\mathcal K),
\qquad
\mathcal K^{(n+1)}=\mathsf{Gen}(\mathcal K^{(n)}),
\label{eq:tri-space-generation}
\end{equation}
where $\Comp_{P,A,T}$ contains the admitted compositions that connect a physical
path, its registration path, and its temporal-support record through their
declared bridges. The \emph{tri-space generated closure} is
\begin{equation}
\operatorname{cl}_{\GOp}^{P,A,T}(\mathcal K^{(0)})=
\bigcup_{n\geq 0}\mathcal K^{(n)}.
\label{eq:tri-space-generated-closure}
\end{equation}
The P-space/A(R)-space component of this construction is supplied by the
observability-induced paired closure of
Theorem~\ref{thm:observability-induced-paired-closure}; T-space extends that paired
closure with the temporal-support carrier of the executed program. The operation
$\GOp$ acts on typed registrations and programs while physical P-space operations
remain theory-relative dynamics.
\end{sdefinition}

\begin{sdefinition}[Theory-and-test program]
\label{def:theory-test-program}
\Defn{} A theory object is a typed registration
\begin{equation}
\tau=
\left(
\mathcal E_\tau,\mathcal D_\tau,\mathcal B_\tau,
\mathcal P_\tau,\mathcal M_\tau,\mathcal J_\tau
\right)
\in\ARspace,
\label{eq:theory-object-registration}
\end{equation}
where the entries specify its energy-functional composition, domain and regime,
typed bridges, predicted P-space paths, predicted A(R)-space registrations, and
comparison/revision conditions. Its execution is the typed cycle
\begin{equation}
\tau
\xrightarrow{\mathsf{design}}
\mathsf{Exp}_{\Tspace}(\tau)
\xrightarrow{\Pspace}
\pi^{P}
\xrightarrow{B}
\pi^{A}
\xrightarrow{\mathsf{compare}}
\tau'.
\label{eq:theory-test-cycle}
\end{equation}
T-space supplies the temporal-support conditions for the experiment and its
comparison path. Hypothesis, test, registration, and revision remain different
typed objects linked by the declared cycle.
\end{sdefinition}

\begin{sdefinition}[Productive generative capacity]
\label{def:productive-generative-capacity}
\Defn{} The tri-space generated closure is \emph{productive} when there is a
partially ordered signature space $\mathcal N$ and a signature map
\begin{equation}
\nu:\operatorname{cl}_{\GOp}^{P,A,T}(\mathcal K^{(0)})\longrightarrow\mathcal N
\label{eq:productive-signature-map}
\end{equation}
such that for every $n$ there is a generated object
$K_{n+1}\in\mathcal K^{(n+1)}$ with
\begin{equation}
\nu(K_{n+1})\not\preceq\nu(K)
\qquad
\text{for every }K\in\mathcal K^{(n)}.
\label{eq:productive-generative-capacity}
\end{equation}
The signature can contain support depth, bridge depth, registration depth,
temporal-support depth, theory depth, proof-obligation depth, or
experimental-composition depth. Productivity records semantic non-collapse, not
merely an unbounded set of operator strings.
\end{sdefinition}

\begin{sproposition}[Compatibility of bookkeeping and unbounded novelty]
\label{prop:bookkeeping-unbounded-novelty}
\Thm{} A bounded bookkeeping verdict for $D_B$ and productive unbounded
generative capacity of
$\operatorname{cl}_{\GOp}^{P,A,T}(\mathcal K^{(0)})$ are jointly coherent. The
first assesses $\mathcal K_B$ after a task-preserving $\LOp_B$ contraction. The
second assesses whether every finite signature-preserving contraction omits a
non-equivalent object generated at a later stage. In particular, if the closure is
productive, then for every finite comparison budget $B$ there exists a generated
object whose signature is not represented in $\mathcal K_B$.
\end{sproposition}

\begin{proof}
Definition~\ref{def:bounded-bookkeeping} licenses the first verdict only for the
fixed target $D_B$. Definition~\ref{def:productive-generative-capacity} supplies a
new non-dominated signature at every finite stage. A finite comparison record
contains only finitely many declared signatures and therefore cannot contain every
later non-dominated signature. The two assessments concern different objects: the
contracted record for the first verdict and the generated closure for the second.
\end{proof}

\paragraph{Computable physics.}
The novelty at issue is not a relabeling of known physical objects. Under the
tri-space closure, physical intervention in P-space, measurement and comparison in
A(R)-space, and temporal enactment in T-space are jointly executable components of
a theory-and-test program. The representation makes hypothesis formation,
experimental composition, registration, diagnosis of representation loss, and
revision into typed objects on which later programs can operate. When the closure
is productive, this capacity expands combinatorially through admitted bridge
compositions while the declared semantic signature prevents raw syntactic growth
from being misclassified as novelty.

%=================================================================
\section{Assessment-relative claims and persistent objections}
\label{supp:assessment-relative}

The preceding section distinguishes a bounded local theorem record from a
productive generated closure. This section makes the corresponding objection
discipline explicit. The distinction does not license a generated claim merely by
calling it global. It specifies what an objection must establish in order to
propagate from a contracted record to the full generated path.

\begin{sdefinition}[Assessment-relative record]
\label{def:assessment-relative-record}
\Defn{} Let $\mathcal K^{(0)}$ be a typed seed record and let
$\mathcal K_B=\LOp_{B,\Theta}(\mathcal K^{(0)})$ be a licensed finite contraction
for declared target $\Theta$. The associated $L$-assessment concerns only the
objects and target distinctions retained in $\mathcal K_B$. The associated
$G$-assessment concerns the observability-induced paired closure
and its guarded recursively generated continuation. Its registration projection is
\begin{equation}
\GenClos(\mathcal K^{(0)};\Theta)=
\bigcup_{n\geq 0}\mathcal K^{(n)},
\qquad
\mathcal K^{(n+1)}=\GOp_{\Theta}\!\left(\mathcal K^{(n)}\right),
\label{eq:assessment-relative-closure}
\end{equation}
with the declared typing, bridge, provenance, admissibility, and target-direction
conditions retained at each stage. The two assessments are non-interchangeable
because they evaluate different records.
\end{sdefinition}

\begin{sdefinition}[$G$-persistent objection]
\label{def:g-persistent-objection}
\Defn{} Let $D$ be an alleged defect in a generated path for target $\Theta$.
The defect is \emph{$G$-persistent} when either: (i) it is witnessed at a required
stage of every admissible continuation of the path; or (ii) there is a valid
certificate that every later generated object factors through the contracted record
on which $D$ was established. An absence statement confined to $\mathcal K_B$ is
not $G$-persistent merely because the contracted record contains no witness for a
later bridge, experiment, provenance relation, or target-directed continuation.
\end{sdefinition}

\begin{sproposition}[A local absence becomes a global defeater only by persistence]
\label{prop:local-absence-persistence}
\Thm{} Suppose an $L$-assessment establishes that a relation $b$ is absent from
$\mathcal K_B$. This establishes only
\begin{equation}
 b\notin \mathcal K_B.
\label{eq:local-absence-only}
\end{equation}
It does not establish $b\notin\GenClos(\mathcal K^{(0)};\Theta)$ unless the
absence is $G$-persistent in the sense of
Definition~\ref{def:g-persistent-objection}. Conversely, a contradiction, failed
bridge or regime condition, inadmissible transition, loss of target direction, or
nonproductive closure is a $G$-level defeater when it is $G$-persistent.
\end{sproposition}

\begin{proof}
Equation~\eqref{eq:local-absence-only} concerns the image of a contraction. A
relation absent from that image may appear only in a later extension, or may be
required but unresolved. It follows that a negative inference about the complete
closure requires either factorization of every later object through the contracted
image or failure of every admissible continuation. These are exactly the two
persistence conditions of Definition~\ref{def:g-persistent-objection}. The converse
holds because a required-stage failure or a certified exhaustive factorization
prevents the target-relevant generated object from being supplied by the path.
\end{proof}

\paragraph{Theorem-level claims and global novelty.}
A familiar parity or factorization form can be a mathematically correct
$L$-description of a finite transition. It does not, by itself, decide whether the
transition's role in a larger physical-theory path is novel. That further question
concerns the bridges, target conditions, test programs, and later generated objects
of the closure. The generated novelty claim succeeds only when these remain typed,
admissible, and non-equivalent to every finite signature-preserving contraction; it
fails if the requisite productivity or target-direction condition fails.

\paragraph{Derivation, testing, and contribution.}
Under a bounded $L$-only projection, the designated physics content remains
hypothesis-labeled because the projection retains selected local derivations, tested
relations, or source objects. Under the observability-induced paired closure, the
cited source objects, native regimes, typed bridges, preserved invariants,
successor relations, derivations, and empirical registrations compose into an
explicit, inspectable generated structure. At the declared assessment budget, the
source-to-composition record emits a global-coherence verdict for the
status-labeled hypothesis. Its open records identify the next realization,
measurement, and quantitative-instantiation tasks internal to that coherent
structure. In the present path, the meta-methodological contribution is visible
under $L$, while the additional physical-theory contribution
is the explicitly staged progression from paired topology to the second-order
gravitation object.

\paragraph{Title, abstract, and deliverable.}
A bounded projection may retain only the theorem-level or methodological component
and therefore make the paper's title or abstract appear to exceed the projected
output. The appropriate question is whether the displayed generated path supplies
the remaining components under its stated conditions. This question is decidable
stage by stage: a successful path has no such mismatch; an unfilled or failed
required stage marks the corresponding claim conditional or open rather than
settling it by local omission.

%=================================================================
\section{Physical states and their registrations}
\label{supp:spaces}

\paragraph{The objects before the notation.}
Every physical state is a potentially recursively defined node in a physical-state
network, and every physical process is a potentially recursively defined
composition of physical processes and therefore an edge or path in that network.
For a declared energy-functional composition, the full network represents the
physical states and processes available in its declared domain. Every observation
is a potentially recursively defined composition of observations and therefore a
node in a registration network, while every process of observation is a potentially
recursively defined composition of observation processes and therefore an edge or
path in that network. Under the same energy-functional composition and its declared
bridges, the registration network represents the corresponding possible
registrations. The definitions below give these two networks their compact P-space
and A(R)-space names without identifying one with the other.


\paragraph{Operational source objects.}
Operational measurement theory provides a source grammar in which a physical
system state, a measurement interaction, an instrument, and a registered outcome
are distinct but related objects~\cite{S_DaviesLewis1970,S_Ozawa1984}. That grammar
supplies the native physical meaning of a registration protocol: it does not
identify the record with the state, and it exposes the carrier and regime in which
state-to-record transfer is licensed. P-space and A(R)-space retain this distinction
while extending it to recursively composed physical and registration paths.

\begin{sdefinition}[Energy-functional composition and paired trajectories]
\label{def:energy-composition-paired-trajectories}
\Defn{} A declared source composition is written
\begin{equation}
\mathcal E=[E_1,\ldots,E_m]_{\mathrm{comp}},
\label{eq:energy-functional-composition}
\end{equation}
where the $E_i$ are the energy functionals used by the selected physical model,
together with their declared ordering, couplings, boundary conditions, scales, and
regime assumptions. The notation does not require that the $E_i$ collapse to one
scalar functional, Hamiltonian, or Lagrangian. If a source theory has no
energy-functional presentation, $\mathcal E$ denotes its declared equivalent
physical generating specification, with that difference retained as a type tag.

For an admissible source state $x_0$, $\mathcal E$ determines or predicts a
physical trajectory
\begin{equation}
\gamma^P_{\mathcal E}:
 x_0\xrightarrow{u_1}\cdots\xrightarrow{u_n}x_n,
 \qquad u_k\in\mathcal O_P.
\label{eq:paired-p-trajectory}
\end{equation}
Under a declared bridge/protocol $B$, the same composition also determines or
predicts the associated registration trajectory wherever the relevant
registrations are defined,
\begin{equation}
\gamma^A_{\mathcal E,B}:
 a_0\xrightarrow{r_1^B}\cdots\xrightarrow{r_n^B}a_n.
\label{eq:paired-a-trajectory}
\end{equation}
The first trajectory is physical; the second is a trajectory of registrations of
physical states or of predicted registrations of physical states. They are paired
by their declared source composition and bridge, not identified as the same object.
Neither totality nor one-to-one correspondence of the bridge is assumed.
\end{sdefinition}

\begin{sdefinition}[P-space]
\label{def:pspace}
\Defn{} A \emph{P-space} instance $\Pspace_{\mathcal E;R,\delta}$ at a declared
energy-functional composition $\mathcal E$, bounded region $R$, and finite
resolution $\delta$ is a finite strongly typed metric multigraph
\begin{equation}
\Pspace_{\mathcal E;R,\delta}=
\left(V_{R,\delta},E_{R,\delta},s,t,\tau,\mu,\partial,\mathcal O_{R,\delta}\right),
\label{eq:pspace-instance-v15}
\end{equation}
representing the observed physical states and/or the physical states predicted by
$\mathcal E$ in the bounded region. Its vertices are declared functional-state
classes, each status-tagged as observed, predicted, inferred, or otherwise
specified; its edges are admissible physical interactions and finite compositions;
its maps $s,t$ assign sources and targets; its typing map $\tau$ assigns carrier,
boundary, interface, operator, regime, support, and scale roles to each vertex and
edge; its metric $\mu$ records declared distinguishability or interaction data at
resolution $\delta$; its $\partial$ records boundary and port data; and its
$\mathcal O_{R,\delta}$ is the admitted operation family at the selected scale and
regime. An edge is admitted only when its source, target, operator, boundary,
interface, and regime types compose. A serialization of this finite graph
represents a bounded physical region in a three-dimensional geometry of arbitrary
curvature, to resolution $\delta$; the graph is not thereby identified with the
physical metric geometry. When $R$ is fixed by context, it may be suppressed from
the notation.
\end{sdefinition}

\begin{sdefinition}[Support signature]\label{def:support}
\Defn{} For a candidate operation $\Gamma$ at resolution $\delta$, the support
signature is $\mathsf{L}$ (local) when the retained registered effect is
recoverable from a bounded neighbourhood without loss of a declared retained
distinction, and $\mathsf{G}$ (global) when extended shape, topology, boundary
structure, or path composition remains irreducible at the retained quotient. Global
support is not acausality.
\end{sdefinition}

\begin{sproposition}[SOP interaction classification]
\label{prop:sop-interaction-classification}
\Asm{} The SOP interaction hypothesis assigns electromagnetic interactions to
$\mathsf{L}$ support and gravitational interactions to $\mathsf{G}$ support at the
declared observation quotient. The source interaction theories provide the native
carrier, law, and observational domains; the present proposition adds the
support-role assignment as a constitutive model-class hypothesis. Its assessment
therefore proceeds through a declared realization record, bridge certificate,
regime certificate, and target-sensitive defeater.
\end{sproposition}

\begin{sdefinition}[A(R)-space and registration non-substitution]
\label{def:arspace}
\Defn{} An \emph{A(R)-space} instance $(\ARspace)_{\mathcal E;R,\delta}$ is a
strongly typed graph representing the registrations generated under the same
declared composition $\mathcal E$ and under a declared registration protocol. Its
vertices are registration states---the states by which an observed physical state,
a predicted physical state, an observed registration, or a predicted registration
becomes coupled, transmitted, recorded, detected, displayed, compared, or made
available to another system---and its edges are admissible registration operations.
Each vertex is status-tagged at least as observed, predicted, inferred, or
otherwise specified. Thus A(R)-space contains, for example, detector records,
encoding states, displays, comparison records, predicted detector outcomes, and
predicted encodings. It does not contain a second copy of the physical state merely
because it registers that state.

A partial typed registration relation
\begin{equation}
\Reg^{B}_{\mathcal E;R,\delta}:
V^{P}_{\mathcal E;R,\delta}\rightharpoonup
V^{A}_{\mathcal E;R,\delta}
\label{eq:paired-registration-relation}
\end{equation}
associates a P-space state with a registration state under a declared bridge and
protocol. A predicted registration remains predicted unless the corresponding
protocol is executed and its observation conditions are met; an observed
registration remains a registration rather than automatically becoming the
physical state it represents. A claim about $p\in V^P$ may be used as a claim about
$a\in V^A$, or conversely, only under an explicit bridge record
\begin{equation}
B=\left(\Reg,\Pi,\mathcal U,\mathcal D,\mathcal K\right),
\label{eq:typed-bridge-v15}
\end{equation}
where $\Pi$ is the selected projection or display map, $\mathcal U$ an uncertainty
or distinguishability model, $\mathcal D$ the declared domain/regime, and
$\mathcal K$ the verification protocol. The bridge may be partial, lossy,
regime-limited, and many-to-one; the fact that both spaces are generated or
predicted from the same declared $\mathcal E$ supplies no unqualified transfer.
\end{sdefinition}

\begin{sdefinition}[Second-Order Physics record]\label{def:sop-record-v9}
\Defn{} An SOP record for a declared physical audit is
\begin{equation}
\SOP_{R,\delta}=\left(\Pspace_{\mathcal E;R,\delta},(\ARspace)_{\mathcal E;R,\delta},B,\Audit,\mathcal W\right),
\label{eq:sop-record-v9}
\end{equation}
where $B$ is the typed bridge, $\Audit$ the declared audit-question family, and
$\mathcal W$ a witness/certificate record. SOP is not a new physical substance. It
is a typed representation architecture that prevents a physical claim, a
registration claim, and a conceptual claim from being silently substituted for
one another.
\end{sdefinition}


%=================================================================
\section{Observability-induced paired closure}
\label{supp:observability-induced-closure}

The paired P-space/A(R)-space distinction supplies the input to a direct
construction of $G$. Observation is a recursively composed topological
transformation from admissible physical paths to admissible registration paths. A
target is observable exactly when the target predicate factors through that
transformation at the declared carrier, regime, support, history, and resolution.
The closure required for this factorization is therefore derived from the
observability condition rather than introduced as a free assessment convention.

\begin{sdefinition}[Declared recursive carrier and registration closures]
\label{def:declared-recursive-closures}
\Defn{} Let $B$ be a declared assessment budget and let $P_0$ and $A_0(R)$ be the
physical and registration seeds admitted for a target $Q$ at resolution $\delta$.
The recursive closures are
\begin{equation}
\mathcal P^{\ast}_{Q,\delta,B}=
\operatorname{Cl}_{\mathrm{adm},B}
\bigl(P_0;\ \mathsf{Comp}_{P},\mathsf{Trans}_{P}\bigr),
\qquad
\mathcal A^{\ast}_{Q,\delta,B}=
\operatorname{Cl}_{\mathrm{ne},B}
\bigl(A_0(R);\ \Reg,\Audit,\Reg\!\circ\!\Reg,\ldots\bigr).
\label{eq:declared-recursive-closures}
\end{equation}
The first retains declared admissible physical states, compositions of physical
states, and recursively generated physical transformations. The second retains
registrations of physical states, registrations of registrations, audit records,
and further recursively generated observation records. These are closures over a
declared carrier, transformation grammar, target, and budget; they do not claim an
enumeration of all physically possible states. A paired closure retains a typed
bridge between them rather than identifying their objects.
\end{sdefinition}

\begin{sdefinition}[Target-relative observation factorization]
\label{def:target-relative-observation-factorization}
\Defn{} Let $\Pi^P_{\mathcal E;R,\delta}$ be the declared admissible P-space path
family, let $\Pi^A_{\mathcal E;R,\delta}$ be the corresponding admissible
registration-path family, and let
\begin{equation}
\ObsMap_{B,\delta}:\Pi^P_{\mathcal E;R,\delta}\rightharpoonup
\Pi^A_{\mathcal E;R,\delta}
\label{eq:observation-transformation}
\end{equation}
be the recursively composed physical-to-registration transformation induced by the
typed bridge $B$ and its declared quotient. For a target predicate
$\Theta_Q:\Pi^P_{\mathcal E;R,\delta}\to\mathcal Q$, the coupled record carries the closures of
Definition~\ref{def:declared-recursive-closures}:
\begin{equation}
\mathcal C_{Q,\delta,B}=
\left(
\mathcal P^{\ast}_{Q,\delta,B},\mathcal A^{\ast}_{Q,\delta,B},B,
\Theta_Q,\mathsf{Regime},\mathsf{Status},\mathsf{Prov}
\right).
\label{eq:coupled-observation-record}
\end{equation}
The target is \emph{observably decidable} in this record when there is a typed
decoder $d_{Q,\delta}$ such that
\begin{equation}
\Theta_Q(\gamma^P)=
d_{Q,\delta}\!\left(\ObsMap_{B,\delta}(\gamma^P)\right)
\qquad
\text{for every }\gamma^P\text{ in the declared target domain.}
\label{eq:target-observation-factorization}
\end{equation}
Equivalently, $\Theta_Q$ is constant on each fiber of the observation
transformation within that domain. The equation preserves the physical and
registration types: $\Theta_Q$ is a predicate of a physical path and
$d_{Q,\delta}$ acts on its registration path through $B$.
\end{sdefinition}

\begin{sdefinition}[Observability obligation support]
\label{def:observability-obligation-support}
\Defn{} Let $\mathfrak U_{Q,\delta}(\mathcal C_0)$ be a declared universe of
typed, non-erasing coupled extensions of a seed record $\mathcal C_0$, ordered by
componentwise inclusion $\preceq$. The \emph{observability obligation support} is
\begin{equation}
\ObsNeed_{Q,\delta}(\mathcal C)=
\left\{r\ \middle|\ 
\begin{array}{l}
r\text{ is a physical relation, physical-composition relation, registration or}\\
r\text{registration-of-registration relation, bridge condition, regime condition,}\\
r\text{support relation, ordered-history relation, or audit record whose retention}\\
r\text{is required for Equation~\eqref{eq:target-observation-factorization} to remain}\\
r\text{typed, admissible, and target-preserving}
\end{array}
\right\}.
\label{eq:observability-obligation-support}
\end{equation}
The membership test is target-relative: a relation enters the support when its
erasure changes the target factorization, the declared physical admissibility of a
path, the legality of $B$, the declared regime, or the audit/provenance condition
through which the factorization is checked. This definition uses the observation
condition itself and introduces no prior $G$-closure.
\end{sdefinition}

\begin{sdefinition}[Observability extension generator]
\label{def:observability-extension-generator}
\Defn{} Let $\mathsf{Emit}_{\rm adm}$ map an observability obligation to its
admissibly realized typed relation or to its typed status record and provenance.
The \emph{observability extension generator} is
\begin{equation}
\mathsf E^{\mathrm{obs}}_{Q,\delta}(\mathcal C)=
\mathcal C\sqcup
\mathsf{Emit}_{\rm adm}\!\left(\ObsNeed_{Q,\delta}(\mathcal C)\right),
\label{eq:observability-extension-generator}
\end{equation}
where $\sqcup$ is non-erasing typed union. The declared closure universe is
\emph{observability complete} when it is a complete lattice under $\preceq$ and
$\mathsf E^{\mathrm{obs}}_{Q,\delta}$ is monotone on the upper interval
$[\mathcal C_0,\top]$. The generator then records the full target-relative support
required by observation while preserving all prior source, bridge, regime, and
status records.
\end{sdefinition}

\begin{stheorem}[Observability-induced paired closure (standard fixed-point construction, re-theoremed here)]
\label{thm:observability-induced-paired-closure}
\Imp{}/\Thm{} The least-fixed-point antecedent is standard~\cite{S_Tarski1955}.
It is re-theoremed here because the proof identifies the paired physical,
registration, bridge, and residue records retained by the present closure. Let
$\mathcal C_0$ be a coupled P-space/A(R)-space record for target $Q$ at resolution
$\delta$, and let its declared extension universe be observability complete in the
sense of Definition~\ref{def:observability-extension-generator}.
Then the least non-erasing paired closure
\begin{equation}
\GPair_{Q,\delta}(\mathcal C_0)=
\operatorname{lfp}_{\mathcal C\succeq\mathcal C_0}
\mathsf E^{\mathrm{obs}}_{Q,\delta}(\mathcal C)
\label{eq:observability-induced-paired-closure}
\end{equation}
exists. It is the least coupled record that contains every relation emitted by the
observability obligations for $Q$. Its derived component closures are
\begin{equation}
G^P_{Q,\delta}=\pi_P\!\left(\GPair_{Q,\delta}(\mathcal C_0)\right),
\qquad
G^{A(R)}_{Q,\delta}=\pi_A\!\left(\GPair_{Q,\delta}(\mathcal C_0)\right),
\label{eq:derived-component-g-closures}
\end{equation}
with the typed bridge retained in the paired record. The closure carries an
observability pass exactly when Equation~\eqref{eq:target-observation-factorization}
has a typed decoder on the declared target domain. Its remaining required
relations are the typed observability residue
\begin{equation}
\ObsRes_{Q,\delta}=
\ObsNeed_{Q,\delta}\!\left(\GPair_{Q,\delta}\right)
\setminus
\mathsf{Realized}\!\left(\GPair_{Q,\delta}\right),
\label{eq:observability-residue}
\end{equation}
which records the next bridge, realization, measurement, or quantitative task in
the same coupled closure.
\end{stheorem}

\begin{sdefinition}[Paired blind-spot locus]
\label{def:paired-blind-spot-locus}
\Defn{} Let $\mathsf{RegRec}(\mathcal A^{\ast}_{Q,\delta,B})$ denote the declared
registration, re-registration, and audit records in the A(R)-closure. The paired
blind-spot locus is
\begin{equation}
\mathsf{Blind}_{Q,\delta,B}(\GPair)=
\left\{(\gamma,\gamma')\ \middle|\ 
\Theta_Q(\gamma)\neq\Theta_Q(\gamma')\ \land\
\forall r\in\mathsf{RegRec}(\mathcal A^{\ast}_{Q,\delta,B}),\ r(\gamma)=r(\gamma')
\right\}.
\label{eq:paired-blind-spot-locus}
\end{equation}
For each such pair the record carries a typed reason field
$\mathsf{Why}\in\{\mathsf{carrier},\mathsf{support},\mathsf{history},
\mathsf{bridge},\mathsf{regime},\mathsf{audit}\}$. The locus is therefore not an
external failure label: it is an output of the same non-erasing paired closure that
records the available physical definitions, observation transformations, and
registration pathways.
\end{sdefinition}

\begin{sproposition}[Known quotient pattern; paired physical criterion]
\label{prop:known-pattern-paired-criterion}
\Thm{}/\Defn{} The fiber criterion and finite proper-support witnesses used below
are familiar mathematical patterns; threshold-information constructions provide one
standard witness family \cite{S_Shamir1979}. Physics also contains domain-specific
restricted-observable indistinguishability, including quantum data hiding
\cite{S_TerhalDiVincenzoLeung2001}, local topological order
\cite{S_BravyiHastingsMichalakis2010}, $N$-representability constraints
\cite{S_LiuChristandlVerstraete2007}, and nonlocal or relational
 gauge-invariant observables in gravity \cite{S_Giddings2005,S_Khavkine2015}.
To the author's knowledge, the present paired criterion is distinct: it computes,
for one declared carrier and budget, the blind-spot locus of recursively composed
physical states and recursively composed observation records jointly. This
proposition makes no claim that existing physics lacks observability blind spots;
it identifies the additional closure object used here to record their target-
relative location.
\end{sproposition}

\begin{proof}
The target-relative observation condition is Equation~\eqref{eq:target-observation-factorization}.
Definition~\ref{def:observability-obligation-support} identifies precisely the
relations that this condition requires to remain target-preserving and typed.
Definition~\ref{def:observability-extension-generator} adds those relations
non-erasively and is monotone on the declared complete lattice. The lattice
fixed-point theorem therefore supplies a least fixed point above $\mathcal C_0$
\cite{S_Tarski1955}, yielding Equation~\eqref{eq:observability-induced-paired-closure}.
Minimality follows because every fixed point above $\mathcal C_0$ contains the
non-erasing emissions required by the generator. The projections in
Equation~\eqref{eq:derived-component-g-closures} retain respectively the physical
and registration components; the bridge remains in the paired object and carries
the factorization relation between them. Equation~\eqref{eq:target-observation-factorization}
gives the pass condition directly. The set difference in
Equation~\eqref{eq:observability-residue} collects requirements emitted by the
same construction whose realization records remain open.
\end{proof}

\begin{sproposition}[The visible gravitational path is a finite prefix of the derived closure]
\label{prop:visible-path-prefix-derived-closure}
\Thm{} For the declared gravitational target, the finite sequence
$S_0\xrightarrow{\GOp_1}\cdots\xrightarrow{\GOp_6}S_6$ is assessed as a finite
source-to-composition prefix inside $\GPair_{Q,\delta}$. Each stage supplies a
typed source object, bridge, invariant, regime, realization record, or successor
relation required by a corresponding observability obligation.
\end{sproposition}

\begin{proof}
Each stage is already registered by a source-to-composition record containing
source objects, native conditions, an invariant, a bridge, an output, a status, and
a typed defeat condition. Such entries belong to the physical, registration,
bridge, regime, support, history, or audit classes in
Definition~\ref{def:observability-obligation-support}. Their non-erasing
composition is therefore an admissible finite prefix of the generator's fixed
point.
\end{proof}

\paragraph{Interpretive consequence.}
The derived object is paired at its origin. $G^P_{Q,\delta}$ carries the least
physical topology required for the target-relevant observation trace;
$G^{A(R)}_{Q,\delta}$ carries the least registration topology that retains,
decodes, audits, and propagates that trace. Their bridge records the physical-to-
registration transformation that makes the target's observability inspectable.
The resulting $G$ is therefore a closure generated by the requirements of
observation itself.

\begin{sdefinition}[Typed gravity-constraint record]
\label{def:typed-constraint-record}
\Defn{} A gravity-facing constraint is represented by the typed record
\begin{equation}
\mathsf{Con}=
\bigl(
\mathsf{State},\mathsf{Bridge},\mathsf{Regime},\mathsf{Support},
\mathsf{History},\mathsf{Register},\mathsf{Certificate}
\bigr).
\label{eq:typed-constraint-record}
\end{equation}
The entries respectively specify the state or transition to which a constraint
applies; the inter-theory or P-space/A(R)-space bridge by which it is transferred;
the regime in which the bridge and imported result are licensed; the support,
boundary, and scale conditions on which the result depends; the ordered preparation
and retention history; the registration conditions by which the result is encoded,
measured, or displayed; and the finite certificate through which a declared
verifier checks the conclusion. A constraint is transported soundly only when the
map to the next theory object, scale, verifier, or successor record preserves every
entry on which the declared audit family changes.
\end{sdefinition}

\begin{sproposition}[Constraint typing across the two spaces]\label{prop:typing}
\Defn{}/\Cnd{} A constraint proven about a P-space object (a physical
configuration) and a constraint proven about an A(R)-space object (a
registration/encoding structure) are different-typed constraints. Their transfer
requires an explicit bridge. The principal cross-theory constraint-transfer error
is the un-bridged substitution of one for the other: importing a bound on an
encoding structure as a bound on a physical configuration, or the reverse. The
worked case is constructed precisely on this boundary, which is why it is a clean
test of a simulation layer's type discipline.
\end{sproposition}


%=================================================================
\section{From interaction topology to registration topology}
\label{supp:topological-bridge}

This section states the role-level bridge required to use the physical interaction
classification in an A(R)-space observability result. It is a typed structural
correspondence. It neither identifies a physical force with a registration operator
nor introduces a new physical field equation.

\begin{sdefinition}[Typed topology of a state network]
\label{def:typed-topology}
\Defn{} At resolution $\delta$, the topology of a strongly typed network $X$ is
\begin{equation}
\Top_\delta(X)=
\left(V_X,E_X,\Pi_X^{\mathrm{adm}},\sim_X,\tau_X,\partial_X,
\Supp_X,\Sched_X\right),
\label{eq:typed-topology-supp}
\end{equation}
where $V_X$ and $E_X$ are typed nodes and edges,
$\Pi_X^{\mathrm{adm}}$ is the family of admissible paths, $\sim_X$ is the
declared path-equivalence relation, $\tau_X$ carries node and edge types,
$\partial_X$ records boundary attachment, and $\Supp_X,\Sched_X$ record support
and scheduling relations. An operation is \emph{topological in $X$} when it
transforms, preserves, composes, or compares this typed relational structure.
\end{sdefinition}

\paragraph{The generated hypothesis before its formal encoding.}

\paragraph{Source interaction objects.}
The strong, weak, and electromagnetic interactions enter this construction through
the established particle-physics interaction inventory~\cite{S_PDG2024}; gravitation
enters through the metric and constraint-bearing object of general relativity
and through its established observational comparison record~\cite{S_Einstein1916,
S_Will2014,S_Berti2015}. These citations anchor the native carriers, laws,
observables, and regimes. The four-role map below adds a new invariant-level
composition across those source objects. Its status is therefore \,\Hyp{}; its
four source-to-composition certificates are instantiated below, each with an
explicit witness, regime, bridge, and defeat field.

\paragraph{The generated hypothesis before its formal encoding.}
The observability-induced $\GOp$ closure emits globally coherent compositions of
existing hypotheses, registrations, and declared model relations. One such emitted
composition is the hypothesis that the weak, strong, electromagnetic, and
gravitational interactions occupy four distinct topological roles. The
support/scheduling signature below is a compact formal encoding of that generated
hypothesis. The bridge hypothesis asks whether the target-relevant invariants of
the two typed roles can be related while preserving their difference.

\begin{sdefinition}[Designated physical interaction-process record]
\label{def:designated-interaction-process-record}
\Defn{} A \emph{P-space interaction (designated physical interaction-process
record)} is a typed object
\begin{equation}
\Gamma_{i,\delta}^{\rho}=
\left(\mathcal C_i^{\rho},\mathcal V_i^{\rho},\Pi_{\delta}^{\mathrm{reg}},
\mathcal K_i^{\rho}\right),
\label{eq:designated-interaction-process-record}
\end{equation}
where $i$ identifies the native interaction sector, $\rho$ the role-bearing
physical process, $\mathcal C_i^{\rho}$ its carrier and declared regime,
$\mathcal V_i^{\rho}$ its interaction/evolution composition,
$\Pi_{\delta}^{\mathrm{reg}}$ its declared registration quotient, and
$\mathcal K_i^{\rho}$ the target-relevant support and scheduling invariant. The
role map therefore classifies neither an interaction name in the abstract nor the
spatial range of a mediator. It classifies the declared native-physics process
record through which that sector enters the P-space/A(R)-space composition.
\end{sdefinition}


\begin{compositioncard}[Strong binding realization]
\label{card:strong-binding-realization}
\Imp{}/\Defn{} Let a declared bounded interaction region $U$ carry a strong-sector
bound-state or confinement-sensitive placement record. The designated record is
\begin{equation}
\Gamma_{\mathrm s,\delta}^{\mathrm{bind}}=
\Pi_{\delta}^{\mathrm{bind}}\!\left[
\mathcal V_{\mathrm s}^{\mathrm{bind}}(U),
\mathcal C_{\mathrm s}^{\mathrm{bound}}(U),
\partial U
\right],
\label{eq:strong-binding-record}
\end{equation}
where $\mathcal V_{\mathrm s}^{\mathrm{bind}}$ is the declared strong-sector
interaction composition and $\mathcal C_{\mathrm s}^{\mathrm{bound}}$ retains the
bound or confinement-sensitive constraint at the selected support. The native
source object is the strong interaction sector in the standard particle-physics
inventory~\cite{S_PDG2024}. The record is \emph{local} at this quotient because the
role predicate is evaluated on the declared bounded support $U$ and its attachment
record. It is \emph{parallel} at a co-registered bound-state/constraint quotient
because the retained placement is evaluated as one jointly constrained
configuration rather than as an ordered observational history. This signature
classifies the designated binding record; scattering, hadronization, or time-resolved
strong-sector records can carry different schedule signatures at their own declared
quotients.
\end{compositioncard}

\begin{compositioncard}[Electromagnetic calibration realization]
\label{card:electromagnetic-calibration-realization}
\Imp{}/\Defn{} Let a declared electromagnetic source, local reference interface,
and detector/calibration readout be connected by a registered propagation path
$\gamma$. The designated record is
\begin{equation}
\Gamma_{\mathrm{em},\delta}^{\mathrm{cal}}=
\Pi_{\delta}^{\mathrm{cal}}\!\left[
\mathcal M_{\mathrm{det}}^{\mathrm{em}}\circ
\mathcal U_{\gamma}^{\mathrm{em}}[A]\circ
\mathcal C_{\mathrm{src}}^{\mathrm{em}}
\right],
\label{eq:electromagnetic-calibration-record}
\end{equation}
where $\mathcal C_{\mathrm{src}}^{\mathrm{em}}$ prepares the declared local
reference signal, $\mathcal U_{\gamma}^{\mathrm{em}}[A]$ propagates it in the
specified electromagnetic carrier, and $\mathcal M_{\mathrm{det}}^{\mathrm{em}}$
registers the calibration relation. The native source object is the electromagnetic
interaction sector in the standard particle-physics inventory~\cite{S_PDG2024}.
The record is \emph{local} at this quotient because its calibration predicate is
fixed by the declared source--interface--detector support and boundary relation. It
is \emph{sequential} because the source, propagation, reference comparison, and
readout order are retained in the calibration history. This signature classifies the
designated calibration/propagation record; extended distributed electromagnetic
records may be assigned another support signature at another declared quotient.
\end{compositioncard}

\begin{compositioncard}[Weak type-transformation realization]
\label{card:weak-transformation-realization}
\Imp{}/\Defn{} Let a weakly prepared flavor label $\alpha$ be weakly registered as
$\beta$ after coherent propagation along $\gamma$. The designated record is
\begin{equation}
\Gamma_{\mathrm w,\delta}^{\mathrm{trans}}=
\Pi_{\delta}^{\mathrm{reg}}\!\left[
\mathcal V_{\beta}^{\mathrm{CC}}\circ
\mathcal U_{\gamma}(\mathcal H_{\mathrm{flav}})\circ
\mathcal V_{\alpha}^{\mathrm{CC}}
\right],
\qquad
\mathcal U_{\gamma}=
\mathcal P\exp\!\left[-i\int_{\gamma}\mathcal H_{\mathrm{flav}}(x)\,dx\right],
\label{eq:weak-transformation-record}
\end{equation}
where $\mathcal V_{\alpha}^{\mathrm{CC}}$ and
$\mathcal V_{\beta}^{\mathrm{CC}}$ are charged-current production and detection
vertices and
\begin{equation}
\mathcal H_{\mathrm{flav}}(x)=
\frac{1}{2E}U M^2U^{\dagger}+V_{\mathrm{mat}}(x)
\label{eq:weak-flavor-hamiltonian}
\end{equation}
is the declared flavor evolution generator. The native source objects are weak
charged-current interaction, flavor/mass mixing, coherent propagation, and, where
present, matter-induced evolution~\cite{S_PDG2024,S_MakiNakagawaSakata1962,
S_Wolfenstein1978}. The record is \emph{global} at this quotient because the
registered $\alpha\!\to\!\beta$ predicate depends jointly on source type, path,
phase/matter profile, and detector type; no proper source-only or detector-only
record decides the registered transformation. It is \emph{sequential} because the
causal composition production $\prec$ propagation $\prec$ detection is retained,
and for declared position-dependent noncommuting flavor generators the propagator
is path ordered. The weak interaction remains short-ranged at its vertices. The
role signature belongs to the source--path--detector transformation record, not to
the range of a $W$ or $Z$ mediator.
\end{compositioncard}

\begin{compositioncard}[Gravitational geometry/constraint realization]
\label{card:gravity-geometry-realization}
\Imp{}/\Defn{} Let $\Sigma$ be a declared spatial slice or equivalent relational
geometry quotient. The designated gravitational record is
\begin{equation}
\Gamma_{\mathrm g,\Sigma,\delta}^{\mathrm{geom}}=
\Pi_{\delta}^{\mathrm{geom}}\!\left[
\Sigma,h_{ij},K_{ij},T_{\mu\nu}|_{\Sigma},
\mathcal H=0,\mathcal H_i=0,\partial\Sigma,\mathcal C_{\mathrm{rel}}
\right],
\label{eq:gravity-geometry-record}
\end{equation}
where $h_{ij}$ and $K_{ij}$ are the induced geometry and extrinsic-curvature data,
$\mathcal H=0$ and $\mathcal H_i=0$ are the Hamiltonian and momentum constraints,
$\partial\Sigma$ records boundary attachment, and $\mathcal C_{\mathrm{rel}}$
records the declared relational geometry registration. The native source object is
the constrained geometric formulation of metric gravitation~\cite{S_Einstein1916,
S_ArnowittDeserMisner1962}. The record is \emph{global} because the retained
geometry/constraint predicate depends on extended relational support, boundary data,
and the constraint-satisfying configuration; a bounded local metric record does not
generally decide that predicate. It is \emph{parallel} at the declared
co-registered slice/geometry quotient because the target is evaluated as one
jointly determined relational configuration, without an order among its contributing
local records. This signature classifies the geometry/constraint sector; radiative
or retarded gravitational process records may carry a sequential schedule signature
at another declared quotient.
\end{compositioncard}

\begin{shypothesis}[Four-role interaction/registration bridge]
\label{hyp:four-role-bridge}
\Hyp{} For each declared P-space interaction (designated physical
interaction-process record) $\Gamma_{i,\delta}^{\rho}$, let
\begin{equation}
\chi_\delta(\Gamma_{i,\delta}^{\rho})=
\left(\Supp_\delta(\Gamma_{i,\delta}^{\rho}),
\Sched_\delta(\Gamma_{i,\delta}^{\rho})\right)
\in\{\mathrm{local},\mathrm{global}\}\times
\{\mathrm{parallel},\mathrm{sequential}\}.
\label{eq:four-role-signature-supp}
\end{equation}
The four designated interaction-process roles and their A(R)-space role
counterparts are
\begin{equation}
\begin{array}{c|c|c}
\begin{gathered}\text{P-space interaction}\\
\text{(designated physical interaction-process record)}\end{gathered}&
\chi_\delta&\text{A(R)-space role}\\
\hline
\Gamma_{\mathrm{s},\delta}^{\mathrm{bind}}&
(\mathrm{local},\mathrm{parallel})&\Register\\
\Gamma_{\mathrm{w},\delta}^{\mathrm{trans}}&
(\mathrm{global},\mathrm{sequential})&\Recall\\
\Gamma_{\mathrm{em},\delta}^{\mathrm{cal}}&
(\mathrm{local},\mathrm{sequential})&\LOp\\
\Gamma_{\mathrm{g},\Sigma,\delta}^{\mathrm{geom}}&
(\mathrm{global},\mathrm{parallel})&\GOp
\end{array}
\label{eq:four-role-bridge-supp}
\end{equation}
The strong binding record supplies constrained local placement; the weak
transformation record supplies globally constrained typed reconfiguration along an
ordered path; the electromagnetic calibration/propagation record supplies local
ordered relational adjustment; and the gravitational geometry/constraint record
supplies extended-support, graph-shape, or geometry-sensitive coherence. The
hypothesis concerns invariant roles within the declared model class. It does not
assert raw operator identity, exhaust the physical content of any interaction, or
claim that one signature applies to every process or regime in a named interaction
sector.
\end{shypothesis}

\begin{sdefinition}[Role-map assessment record]
\label{def:role-map-assessment-record}
\Defn{} For each designated sector $i\in\{\mathrm s,\mathrm w,\mathrm{em},\mathrm g\}$,
a finite assessment of the four-role hypothesis carries
\begin{equation}
\mathsf{RoleCert}_{i,\delta}=
\left(
\mathsf{Src}_{i,\delta},\chi_{i,\delta},\Bridge_{i,\delta},
\mathsf{Regime}_{i,\delta},\mathsf{Witness}_{i,\delta},
\mathsf{Defeat}_{i,\delta}
\right),
\label{eq:role-map-assessment-record}
\end{equation}
and the four-role certificate family is
\begin{equation}
\mathsf{RoleCert}^{(4)}_\delta=
\left\{
\mathsf{RoleCert}_{\mathrm s,\delta},
\mathsf{RoleCert}_{\mathrm w,\delta},
\mathsf{RoleCert}_{\mathrm{em},\delta},
\mathsf{RoleCert}_{\mathrm g,\delta}
\right\}.
\label{eq:four-role-certificate-family}
\end{equation}
Here $\mathsf{Src}_{i,\delta}$ records the cited interaction object and designated
physical interaction-process record, $\chi_{i,\delta}$ the proposed invariant
signature, $\Bridge_{i,\delta}$ the cross-space map,
$\mathsf{Regime}_{i,\delta}$ the carrier and scale conditions,
$\mathsf{Witness}_{i,\delta}$ a realization or observational certificate, and
$\mathsf{Defeat}_{i,\delta}$ a typed contradiction or failed certificate. The table
in Hypothesis~\ref{hyp:four-role-bridge} is therefore assessed through an explicit
source-anchored certificate family carrying the cited objects, designated processes,
proposed invariants, bridges, regimes, realization witnesses, and defeat fields.
\end{sdefinition}

\begin{sdefinition}[Invariant interaction-registration bridge]
\label{def:invariant-interaction-registration-bridge}
\Defn{} Let
\begin{equation}
(C_{\mathrm{s}}^A,C_{\mathrm{w}}^A,C_{\mathrm{em}}^A,C_{\mathrm{g}}^A)
=(\Register,\Recall,\LOp,\GOp).
\label{eq:registration-role-basis}
\end{equation}
An \emph{invariant interaction-registration bridge} is a family of typed maps
\begin{equation}
\Bridge_{i,\delta}:
\Inv_\delta(\Gamma_i^P)\cong\Inv_\delta(C_i^A)
\label{eq:invariant-interaction-registration-bridge}
\end{equation}
that preserves the declared support, scheduling, boundary, admissible-path, and
role signatures. In particular, it preserves the topological information that can
change the declared target while retaining the type distinction
\begin{equation}
\Gamma_i^P\neq C_i^A.
\label{eq:nonidentity-force-registration}
\end{equation}
\end{sdefinition}


\begin{compositioncard}[Instantiated four-role certificates]
\label{card:instantiated-four-role-certificates}
\Hyp{} The certificate family of
Equation~\eqref{eq:four-role-certificate-family} is instantiated by the four
realization records already declared above:
\begin{align}
\mathsf{Witness}_{\mathrm s,\delta}&:=
\Gamma_{\mathrm s,\delta}^{\mathrm{bind}}
&&\text{(Card~\ref{card:strong-binding-realization})},\nonumber\\
\mathsf{Witness}_{\mathrm w,\delta}&:=
\Gamma_{\mathrm w,\delta}^{\mathrm{trans}}
&&\text{(Card~\ref{card:weak-transformation-realization})},\nonumber\\
\mathsf{Witness}_{\mathrm{em},\delta}&:=
\Gamma_{\mathrm{em},\delta}^{\mathrm{cal}}
&&\text{(Card~\ref{card:electromagnetic-calibration-realization})},\nonumber\\
\mathsf{Witness}_{\mathrm g,\Sigma,\delta}&:=
\Gamma_{\mathrm g,\Sigma,\delta}^{\mathrm{geom}}
&&\text{(Card~\ref{card:gravity-geometry-realization})}.
\label{eq:instantiated-role-witnesses}
\end{align}
Their bridges are the corresponding invariant maps
$\Bridge_{i,\delta}:\Inv_\delta(\Gamma_i^P)\cong\Inv_\delta(C_i^A)$ of
Definition~\ref{def:invariant-interaction-registration-bridge}. Their regimes are
respectively the declared bounded binding quotient, source--path--detector flavor
quotient, local electromagnetic calibration quotient, and co-registered
geometry/constraint quotient. \paragraph{Typed defeat fields.}
$\mathsf{Defeat}_{\mathrm s,\delta}$ is registered when the binding predicate
requires support or retained order outside its declared bounded quotient.
$\mathsf{Defeat}_{\mathrm w,\delta}$ is registered when the transformation
predicate fails to retain its source--path--detector support or declared causal
ordering. $\mathsf{Defeat}_{\mathrm{em},\delta}$ is registered when the calibration
predicate fails to retain its declared local support or ordered reference history.
$\mathsf{Defeat}_{\mathrm g,\Sigma,\delta}$ is registered when the
geometry/constraint predicate fails to retain extended relational support or
requires an order excluded by the declared co-registered quotient.
\label{eq:instantiated-role-defeats}
Thus the signatures are read first from independently specified P-space process
records and then transported through explicit invariant bridges to the A(R)-space
basis. At the declared assessment budget, the instantiated family emits
\begin{equation}
\mathsf{GC}_B\!\left(\mathsf{RoleCert}^{(4)}_\delta\right)=\mathsf{pass}
\label{eq:four-role-gc-pass}
\end{equation}
for the status-labeled four-role hypothesis: a globally coherent composition of
existing interaction hypotheses, native process records, and declared bridges.
\end{compositioncard}

\begin{sproposition}[Topological closure in the two spaces]
\label{prop:topological-closure-two-spaces}
\Thm{} Under Hypothesis~\ref{hyp:four-role-bridge}, every admitted designated
P-space interaction-process record has a declared topological interaction role,
and every admitted A(R)-space interaction is a typed program over
$\langle\Register,\Recall,\LOp,\GOp\rangle_{\rm rec}$. Larger-scale physical
interactions and larger-scale registrations are respectively typed compositions of
their P-space and A(R)-space bases. Their relation across spaces is licensed only
by an invariant interaction-registration bridge.
\end{sproposition}

\begin{proof}
The first clause follows from Definition~\ref{def:typed-topology} and the
support/scheduling role assignment of Hypothesis~\ref{hyp:four-role-bridge}. The
second is the recursive constraint-transport closure of
Proposition~\ref{prop:one-assessment-grammar}. Closure under typed composition gives
the larger-scale statement. Equation~\eqref{eq:nonidentity-force-registration}
excludes substitution of raw operations; Definition~\ref{def:invariant-interaction-registration-bridge}
supplies the only declared cross-space relation.
\end{proof}

\paragraph{Consequence for observability.}
Matter observes another material system through a topological composite: physical
access to a P-space subnetwork, a bridge that emits registrations, and a quotient
that retains only the observer's material record. A blind spot is therefore a
topological fiber condition: the composite map identifies physical paths or states
that remain distinct for the target. The next observability section gives this
condition in energy-uniform form and introduces the synchronization residue as a
specific topological coordinate that can vary inside such a fiber.

%=================================================================
\section{Interaction role: attractive, opposite, and inverse}
\label{supp:interaction-role}

The following distinction is definitional for the worked hypothesis. Ordered
admissible operations, attainable sets, and reachability furnish the relevant
control-theoretic source grammar. The classical formulation of that grammar is
given by Jurdjevic and Sussmann~\cite{S_JurdjevicSussmann1972}. The present
construction composes it with gravitational, bridge, regime, support, and
registration records. A conventional attractive-gravity description
propagates a current P-space state forward under a declared local interaction law.
An \emph{opposite} or repulsive-gravity proposal would alter a local interaction
term, for example by a sign reversal within a specified first-order theory. Neither
move is made here.

\begin{sdefinition}[Target-conditioned inverse-interaction schedule]
\label{def:target-conditioned-schedule}
\Defn{} Fix a declared target P-space configuration $x_\star$, an
energy-functional composition $\mathcal E$, and a bridge/registration discipline
$B$. A \emph{target-conditioned inverse-interaction schedule} is a typed ordered
record
\begin{equation}
\Sched_{\mathcal E,B}(x_\star)=
\left\langle
(u_k,\Delta t_k,\mathcal B_k,\mathcal R_k,S_k,B_k)
\right\rangle_{k=1}^{n},
\label{eq:target-conditioned-schedule}
\end{equation}
where $u_k\in\mathcal O_P$ is an ordinary forward-time admissible physical
operation, $\Delta t_k$ is its declared timing/duration, $\mathcal B_k$ its active
boundary/interface condition, $\mathcal R_k$ its regime record, $S_k$ its support,
and $B_k$ its physical-to-registration bridge. The schedule is selected relative
to $x_\star$ and is admissible only if each transition composes and its P-space and
A(R)-space records remain licensed. It realizes the inverse hypothesis only
conditionally, if its full physical composition reaches or appropriately
approximates $x_\star$ without violating the declared corridor predicate.
\end{sdefinition}

\begin{sproposition}[Scheduling is an implementation condition, not a sign flip]
\label{prop:schedule-not-signflip}
\Defn{}/\Cnd{} The adjective \emph{inverse} refers here to a target-conditioned
redirection of constraint propagation: a desired global configuration constrains
the selection of a forward-time sequence of admissible operations. It does not
refer to $u_k\mapsto-u_k$, to a reversal of physical temporal orientation, or to a
claim that a target-conditioned schedule exists. Timing is one coordinate of the
schedule, but scheduling also retains operation order, support, bridge,
boundary/interface, regime, and repeated basepoint/re-linearization conditions.
Consequently, a mere temporal pulse sequence without those typed relations is not
an inverse-interaction witness in the present sense. In the finite generated path of
Section~\ref{supp:generated-path}, this schedule is $S_5$: it is the composition
from which the second-order theory object $S_6$ is emitted.
\end{sproposition}


%=================================================================
\section{Visible generated path to a second-order theory of gravitation}
\label{supp:generated-path}

The observability factorization of
Section~\ref{supp:observability-induced-closure} derives the paired $\GPair$
closure before the physical-theory path begins. The finite witness and fiber
criterion occur at one auditable transition in a longer generated path inside that
closure. The point of making the path explicit is to expose a finite prefix whose
inputs, outputs, local projections, and failure conditions can each be inspected.
In the terminology of Section~\ref{supp:assessment-relative}, an objection to a
later stage is global when it persists through the corresponding admissible
continuations.

\begin{sdefinition}[Finite generated gravitation path]
\label{def:finite-generated-gravitation-path}
\Defn{} Let $S_0$ be a typed registration of existing hypotheses about physical
interaction, observation, gravitation, and their declared energy-functional domains.
Within the observability-induced paired closure for the declared gravitational
target, a \emph{finite generated gravitation path} is the sequence
\begin{equation}
S_0\xrightarrow{\GOp_1}S_1\xrightarrow{\GOp_2}S_2\xrightarrow{\GOp_3}S_3
\xrightarrow{\GOp_4}S_4\xrightarrow{\GOp_5}S_5\xrightarrow{\GOp_6}S_6,
\label{eq:finite-generated-gravitation-path}
\end{equation}
where
\begin{align}
S_0&:\ \text{existing interaction, observation, and gravitation hypotheses};\nonumber\\
S_1&:\ \text{paired recursively composed P-space and A(R)-space topologies};\nonumber\\
S_2&:\ \text{the generated four-role interaction hypothesis};\nonumber\\
S_3&:\ \text{the typed interaction-registration bridge and synchronization residue};\nonumber\\
S_4&:\ \text{the matter-observation blind-spot criterion};\nonumber\\
S_5&:\ \text{a target-conditioned inverse-gravity composition};\nonumber\\
S_6&:\ \text{a second-order theory of gravitation.}
\label{eq:finite-generated-gravitation-stages}
\end{align}
Each $\GOp_i$ is a typed composition program with declared inputs, bridge and
regime conditions, target, provenance, and possible failure predicates. The induced
paths in the three coupled spaces are successively composed as
\begin{equation}
P_a\rightarrow P_b\rightarrow P_c\rightarrow\cdots,\qquad
A_a\rightarrow A_b\rightarrow A_c\rightarrow\cdots,\qquad
T_a\rightarrow T_b\rightarrow T_c\rightarrow\cdots.
\label{eq:coupled-generated-paths}
\end{equation}
The symbols denote physical, registration, and theory/temporal-support paths;
they do not identify their nodes or operations.
\end{sdefinition}

\begin{compositioncard}[Visible source-to-composition ledger]
\label{card:visible-source-ledger}
\Defn{} The observability factorization first supplies the pre-path closure
condition: the coupled record retains the relations by which $\Theta_Q$ factors
through a declared registration. The path of
Definition~\ref{def:finite-generated-gravitation-path} is then assessed through the
following source-to-composition records inside that closure.
\begin{description}
\item[$\mathsf{SC}_{\mathrm{obs}}$---observability-induced closure.] \textbf{Source:}
the typed physical-state, measurement, and registration grammar of
Definitions~\ref{def:energy-composition-paired-trajectories}--\ref{def:arspace}.
\textbf{Invariant:} target factorization through the observation transformation.
\textbf{Bridge:} $B$ in Equation~\eqref{eq:paired-registration-relation}.
\textbf{Output:} $\GPair_{Q,\delta}$. \textbf{Status:} \Thm{} under the declared
complete-lattice and monotone-generator conditions of
Theorem~\ref{thm:observability-induced-paired-closure}.
\item[$\mathsf{SC}_0$---source registration.] \textbf{Sources:} standard-model
interaction objects~\cite{S_PDG2024}, metric gravity and its empirical constraint
practice~\cite{S_Einstein1916,S_Will2014,S_Berti2015}, and operational measurement
objects~\cite{S_DaviesLewis1970,S_Ozawa1984}. \textbf{Output:} $S_0$, a typed
registration of constituent source objects with native carriers and regimes.
\item[$\mathsf{SC}_1$---paired topology.] \textbf{Invariant:} the distinction
between physical-state evolution and state registration. \textbf{Bridge:} the
partial typed measurement/registration relation of
Definition~\ref{def:arspace}. \textbf{Output:} $S_1$. \textbf{Status:} \Defn{}.
\item[$\mathsf{SC}_2$.] \textbf{Interaction-role composition. Sources:}
strong binding and electromagnetic calibration/propagation records from the
particle-physics inventory~\cite{S_PDG2024}, weak charged-current flavor-
transformation records~\cite{S_MakiNakagawaSakata1962,S_Wolfenstein1978}, and
constrained gravitational geometry/constraint records~\cite{S_ArnowittDeserMisner1962}.
\textbf{Invariant:} designated interaction-process role at declared support and
scheduling resolution. \textbf{Bridge:} Hypothesis~\ref{hyp:four-role-bridge} and
the instantiated certificate family of Card~\ref{card:instantiated-four-role-certificates}.
\textbf{Output:} $S_2$ with $\mathsf{RoleCert}^{(4)}_\delta$.
\textbf{Status:} \Hyp{} with
$\mathsf{GC}_B(\mathsf{RoleCert}^{(4)}_\delta)=\mathsf{pass}$ at the declared budget.
\item[$\mathsf{SC}_3$---typed transfer.] \textbf{Invariant:} bridge legality,
regime, ordered history, support, and synchronization residue.
\textbf{Bridge:} Definition~\ref{def:invariant-interaction-registration-bridge}.
\textbf{Output:} $S_3$. \textbf{Status:} \Defn{}/\Hyp{}.
\item[$\mathsf{SC}_4$---blind-spot criterion.] \textbf{Source antecedent:}
proper-subset information separation, for which threshold-information constructions
provide a familiar source pattern~\cite{S_Shamir1979}. \textbf{Bridge:} the
corridor-specific quotient/fiber record of Sections~\ref{supp:strict-separation}
and~\ref{supp:matter-blindspot}. \textbf{Output:} $S_4$. \textbf{Status:} \Thm{}
relative to its witness class.
\item[$\mathsf{SC}_5$---target-conditioned schedule.] \textbf{Source object:}
ordered reachability under admissible generators~\cite{S_JurdjevicSussmann1972}.
\textbf{Bridge:} Definition~\ref{def:target-conditioned-schedule}. \textbf{Output:}
$S_5$. \textbf{Logical form:} \Cnd{} on a declared admissible schedule.
\textbf{Composition verdict:} schedule grammar retained as an audited component of
$\GCcert_{B_\Lambda}$, not a separately self-assigned terminal verdict.
\textbf{Continuation:} the same-carrier continuation certificate of
Definition~\ref{def:semiclassical-continuation-certificate} or its typed break
frontier.
\item[$\mathsf{SC}_{\Lambda}$.] \textbf{Gravitational leverage composition. Sources:}
extended encoding, entanglement--geometry coupling, path integration, metric-
gravity self-coupling, and ordered reachability objects of
Sections~\ref{supp:imported} and~\ref{supp:interaction-role}.
\textbf{Bridge:} Definition~\ref{def:gravitational-leverage-composition}.
\textbf{Output:} $\Lambda^{\mathrm{grav}}_\delta$.
\textbf{Logical form:} \Hyp{} with a named productivity coordinate and a
conditional target-depth consequence. \textbf{Composition verdict:}
\GCpass{} at $B_\Lambda$ by the verified certificate of
Definition~\ref{def:gravitational-leverage-gcert} and
Theorem~\ref{thm:verify-gravitational-leverage-gcert}.
\textbf{Continuation:} the certificate of
Definition~\ref{def:semiclassical-continuation-certificate}, its realization
record, and named break frontier.
\item[$\mathsf{SC}_{\mathrm{gain}}$---selected physical gain witnesses.]
\textbf{Sources:} the shock-wave/eikonal relation
\cite{S_DraytHooft1985,S_ShenkerStanford2014,S_ShenkerStanford2014Multiple}
and the Jeans Euler--Poisson relation
\cite{S_Jeans1902,S_BinneyTremaine2008,S_Chavanis2010}.
\textbf{Bridge:} Definition~\ref{def:recursive-interaction-network}.
\textbf{Output:} the selected shock gain family and the selected non-black-hole
Jeans feedback network. \textbf{Status:} \Thm{} with closed finite target
windows by Theorem~\ref{thm:nonempty-gravitational-gain-class}; no witness is
substituted for arbitrary metric control or the distinct holographic continuation
carrier.
\item[$\mathsf{SC}_6$---second-order gravitation object.] \textbf{Sources:}
gravitational constraint and comparison objects~\cite{S_Will2014,S_Berti2015},
the imported entanglement/encoding objects of Section~\ref{supp:imported}, and
the selected physical witness family of Section~\ref{supp:recursive-interaction-network-gain}.
\textbf{Bridge:} the joint theory/test/registration/leverage composition of
Definition~\ref{def:second-order-gravitation-theory}. \textbf{Output:} $S_6$.
\textbf{Logical form:} typed \Hyp{} with internal productivity and
realization coordinates. \textbf{Composition verdict:} \GCpass{} at $B_\Lambda$
through the verified conditional leverage certificate of
Theorem~\ref{thm:verify-gravitational-leverage-gcert}; no realization coordinate
is promoted to an empirical assertion by that verdict.
\textbf{Continuation:} carrier-indexed realization, empirical-instantiation,
and typed defeat fields; the selected shock and Jeans target windows are closed
physical witness records, while the generic holographic continuation remains a
distinct carrier-indexed record.
\end{description}
A source citation therefore establishes a constituent object in its native
domain, including its lawful derivational and empirical relations. The globally
coherent hypothesis is the explicit composition of those cited objects through the
recorded bridges, invariants, regimes, and defeaters. At the declared assessment
budget, that source-to-composition record emits the global-coherence verdict for the
status-labeled hypothesis; its open records identify the next realization and
quantitative-instantiation tasks.
\end{compositioncard}

\begin{sproposition}[A finite local projection does not exhaust the displayed path]
\label{prop:finite-path-not-exhausted-by-quotient}
\Thm{} The fiber criterion appears in the transition from the bridge/residue
structure $S_3$ to the blind-spot criterion $S_4$. Recognition that this transition
has a familiar quotient-theoretic form is an admissible local relation. It does not
by itself establish that $S_5$ or $S_6$ factor through the same local projection.
A global reduction of the displayed path requires a target-preservation certificate
for every later stage and its declared bridge, support, scheduling, and
experiment/revision relations.
\end{sproposition}

\begin{proof}
The fiber criterion decides only whether a declared target is constant on the fibers
of a selected quotient. The objects $S_5$ and $S_6$ add target-conditioned physical
schedules, their paired registrations, and the theory/test programs generated from
them. By Definition~\ref{def:bounded-bookkeeping}, a local contraction supplies a
global verdict only when it preserves the target-relevant signature of the objects
being assessed. The stage-specific relations are therefore obligations of any such
reduction.
\end{proof}

\begin{sdefinition}[Second-order theory of gravitation]
\label{def:second-order-gravitation-theory}
\Defn{} The output $S_6$ is a \emph{second-order theory of gravitation} when it is
a theory object
\begin{equation}
\tau_{\mathrm{grav}}^{(2)}=
\left(
\mathcal E,\mathcal D,\mathcal B,
\mathcal P_{\mathrm{grav}},\mathcal M_{\mathrm{grav}},
\mathcal J_{\mathrm{grav}},\mathcal S_{\star},
\Lambda^{\mathrm{grav}}_\delta,\mathcal W^{\mathrm{grav}}
\right),
\label{eq:second-order-gravitation-theory}
\end{equation}
whose entries specify: a gravitational energy-functional composition; domain and
regime conditions; typed physical-registration bridges; predicted gravitational
P-space paths; predicted A(R)-space registrations; comparison and revision
conditions; a target-conditioned schedule family $\mathcal S_{\star}$; the explicit
conditional leverage composition of
Definition~\ref{def:gravitational-leverage-composition}; and the selected physical
witness family $\mathcal W^{\mathrm{grav}}$ of
Theorem~\ref{thm:nonempty-gravitational-gain-class}. Its second-order content
is that these physical, registration, leverage, and theory/test objects are jointly
represented and recursively composable. It is not merely a first-order trajectory or
a force-sign alteration.
\end{sdefinition}

\paragraph{Gravitational source and empirical-instantiation objects.}
The first-order gravitational objects retained by $\tau_{\mathrm{grav}}^{(2)}$
include a metric/constraint-bearing gravitational theory and its established
weak-field, strong-field, and observational comparison practices
\cite{S_Einstein1916,S_Will2014,S_Berti2015}. Their role in $S_6$ is positive and
specific: they supply the physical carrier, reference observables, constraint
records, derivations, and empirical registrations through which the second-order
object is instantiated. Their typed composition with the declared bridges,
registrations, revisions, schedules, and
$\Lambda^{\mathrm{grav}}_\delta$ is assessed through the finite certificate of
Definition~\ref{def:gravitational-leverage-gcert}. Theorem~\ref{thm:verify-gravitational-leverage-gcert}
derives its in-budget global-coherence verdict for the conditional, status-labeled
$S_6$ hypothesis. The generic holographic leverage subrecord carries its
conditional productivity coordinate and target-depth consequence; the selected
witness family separately supplies finite physical gain recurrences and registered
target windows. The realization records preserve the remaining carrier-specific
quantitative and experimental tasks without conflating them.

\begin{sdefinition}[Asymptotically realizing schedule family]
\label{def:asymptotically-realizing-schedule-family}
\Defn{} Fix a target configuration $x_\star$ and declared target distance
$d_\star$. A family $\{\mathcal S_n\}_{n\geq 0}\subseteq\mathcal S_{\star}$ is
\emph{asymptotically realizing} when every $\mathcal S_n$ is an admissible
P-space/A(R)-space composition and
\begin{equation}
\lim_{n\rightarrow\infty}d_\star\!\left(
\mathsf{Run}(\mathcal S_n),x_\star\right)=0.
\label{eq:asymptotic-target-realization}
\end{equation}
Equivalently,
\begin{equation}
\forall\varepsilon>0\ \exists n\ \text{such that}\quad
 d_\star\!\left(\mathsf{Run}(\mathcal S_n),x_\star\right)<\varepsilon.
\label{eq:epsilon-asymptotic-target-realization}
\end{equation}
\end{sdefinition}

\begin{sproposition}[Generated asymptotic realization]
\label{prop:generated-asymptotic-realization}
\Cnd{} If the $S_5$ schedule family is emitted by productive $\GOp$ composition,
remains admissible under the declared bridges and regimes, and is asymptotically
realizing in the sense of Definition~\ref{def:asymptotically-realizing-schedule-family},
then $S_6$ contains an asymptotic realization guarantee within its declared model:
every finite target tolerance is reached by an emitted admissible schedule. This is
the relevant $G$-level physical-realization claim. A particular finite experiment
remains a separately status-labeled registration of that claim.
\end{sproposition}

%=================================================================
\section{Imported results used for coherence}
\label{supp:imported}

Each card records a cited external source object, marked \Imp{}, with its native
regime, the invariant it contributes, and its compositional role. The accompanying
extension status identifies the next bridge, realization condition, or empirical
registration required to extend that source object into the present hypothesis. No
card is re-derived.


\begin{scard}[Entanglement-built connectivity]\label{card:vr}
\Imp{} \textbf{Imported:} in the entanglement-first picture of emergent gravity,
spatial connectivity is associated with the entanglement pattern between
subsystems; reducing entanglement increases geometric separation and, in the limit,
disconnects regions \cite{S_VanRaamsdonk2010, S_MaldacenaSusskind2013}.
\textbf{Native regime:} anti-de~Sitter / conformal-field-theory duality and related
semiclassical settings. \textbf{Compositional role:} identifies an
entanglement/encoding structure of $\mathsf{G}$ support as the source registration
variable. \textbf{Extension status:} transfer to another gravitational carrier is
recorded as a separate bridge and empirical-registration condition.
\end{scard}

\begin{scard}[Local semiclassical entanglement--geometry response]
\label{card:jacobson}
\Imp{} \textbf{Imported:} under an entanglement-equilibrium hypothesis, a first-law
relation $\delta\SEE=\delta\langle H\rangle$ for a small region yields the linearized
Einstein equation \cite{S_Jacobson2016}. \textbf{Native regime:} local,
linearized, small-region, semiclassical. \textbf{Compositional role:} an
independent local-semiclassical realization family for a differential
entanglement--geometry relation. \textbf{Extension status:} it is not used as an
interchangeable all-regime kernel for the selected holographic worked path.
\end{scard}

\begin{scard}[Holographic linear entanglement--geometry response]
\label{card:faulkner-linear}
\Imp{} \textbf{Imported:} for small perturbations around the CFT vacuum with a
semiclassical holographic dual, first-law constraints for ball-shaped regions are
equivalent to gravitational equations linearized around pure AdS
\cite{S_FaulknerEtAl2014}. \textbf{Native carrier and regime:}
$\rho_{\mathrm{AdS/CFT}}$, linearized around AdS, holographic semiclassical code
subspace. \textbf{Compositional role:} supplies the carrier-indexed local response
$\delta g|_x=\mathcal K_{\rho}(x)\,\delta\SEE|_x$ used in the worked path.
\end{scard}

\begin{scard}[Second-order holographic continuation prefix]
\label{card:faulkner-nonlinear}
\Imp{} \textbf{Imported:} for a class of sourced CFT states, entanglement data for
ball-shaped regions admits an asymptotically AdS geometric representation satisfying
Einstein's equations perturbatively through second order in the sources
\cite{S_FaulknerEtAl2017}. \textbf{Native carrier and regime:}
$\rho_{\mathrm{AdS/CFT}}$, perturbative second order about the selected background.
\textbf{Compositional role:} supplies a nonlinear continuation prefix inside the
same carrier; it does not license an unrestricted strong-field kernel.
\end{scard}

\begin{scard}[Nearby-state bulk/boundary transport]
\label{card:jlms}
\Imp{} \textbf{Imported:} relative entropy of nearby boundary states is related to
bulk relative entropy to leading order in the bulk gravitational coupling, with a
corresponding modular-flow relation in the entanglement wedge
\cite{S_JafferisEtAl2016}. \textbf{Native carrier and regime:}
$\rho_{\mathrm{AdS/CFT}}$, nearby states, leading order in the bulk gravitational
coupling. \textbf{Compositional role:} supplies a same-carrier transport/stability
record for overlaps of successive local response charts.
\end{scard}

\begin{scard}[Holographic encoding and wedge-local reconstruction]\label{card:qec}
\Imp{} \textbf{Imported:} in holographic models, bulk geometric information is
encoded redundantly across an extended boundary as in a quantum error-correcting
code, and a bulk operator is reconstructible from a boundary region only within
that region's entanglement wedge, commuting with operations spacelike to it
\cite{S_AlmheiriDongHarlow2015, S_PastawskiEtAl2015, S_DongHarlowWall2016}.
\textbf{Native regime:} holographic (AdS/CFT-type) models. \textbf{Compositional
role:} supplies the extended-support encoding relation, protection of the
registration variable against single local operations, and the wedge-locality used
in Section~\ref{supp:nosignal}. \textbf{Extension status:} transfer to another
carrier is retained as an explicit reconstruction-locality bridge.
\end{scard}

\begin{scard}[No-signalling of quantum operations; nonlinearity and signalling]\label{card:nosig}
\Imp{}/\Thm{} \textbf{Imported:} (i)~any composition of local CPTP maps,
measurements, and classically communicated feedback is non-signalling: the reduced
state of an untouched subsystem is invariant under the choice of local operation
\cite{S_NielsenChuang2010}. (ii)~Deterministic nonlinear modifications of quantum
evolution generically permit superluminal signalling \cite{S_Gisin1990,
S_Polchinski1991}. \textbf{Use here:} (i)~makes a path of in-regime CPTP steps
non-signalling by composition; (ii)~is the obstruction that applies to a
\emph{direct} carrierless channel, motivating the actuator-composition framing,
which introduces no nonlinear state evolution.
\end{scard}

%=================================================================
\begin{compositioncard}[Entanglement--geometry composition junction]
\label{card:entanglement-junction}
\Defn{} Cards~\ref{card:vr}, \ref{card:faulkner-linear},
\ref{card:faulkner-nonlinear}, \ref{card:jlms}, and~\ref{card:qec} provide
extended-support entanglement/encoding, carrier-indexed local response,
second-order continuation-prefix, nearby-state transport, and distributed
reconstruction objects. The present worked case selects the common carrier
$\rho_{\mathrm{AdS/CFT}}$ and composes them by the typed record
\begin{equation}
\mathsf{Junc}_{\mathrm{EG}}=
\left(
\rho_{\mathrm{AdS/CFT}},\ \mathsf{Enc}_{\mathrm{ext}},\ 
\delta S_{\mathrm{EE}}\leftrightarrow\delta g,
\ \mathsf{Cont}^{\mathrm{sc}}_{\rho},\ \mathsf{Rec}_{\mathrm{wedge}},
\ \mathsf{Regime},\ \mathsf{Bridge},\ \mathsf{Status}
\right).
\end{equation}
The cited cards establish constituent objects in their native regimes. The new
hypothesis is that a same-carrier, in-regime, target-conditioned path can compose
them while retaining all bridge, support, history, error, and safety conditions.
Its continuation certificate makes the existence, obstruction, or break of that path
an explicit internal branch of the globally coherent composition.
\end{compositioncard}

%=================================================================
\section{The worked case as witness}
\label{supp:case}

The worked case exists to witness that the constraint class of
Proposition~\ref{prop:insufficiency-supp} is non-empty. It is the candidate
physical implementation of Definition~\ref{def:target-conditioned-schedule}: the
target condition selects an ordered, causal, in-regime preparation schedule, while
each physical interaction itself remains forward in time. The worked case selects the common holographic carrier
$\rho=\rho_{\mathrm{AdS/CFT}}$ of Cards~\ref{card:faulkner-linear},
\ref{card:faulkner-nonlinear}, \ref{card:jlms}, and~\ref{card:qec}. It treats the
linear response at a registered basepoint $x$ as
\begin{equation}
\delta g\big|_{x} \;=\; \mathcal K_{\rho}(x)\,\delta\SEE\big|_{x},
\qquad x \in \mathcal R_{\rho}(x),
\label{eq:coupling}
\end{equation}
with $\mathcal K_{\rho}(x)$ the imported carrier-indexed response kernel. The
second-order card supplies a perturbative continuation prefix; no fixed
$\mathcal K_{\rho}(x_0)$ is asserted outside its declared local regime.

\begin{sproposition}[Path integration]\label{prop:integration}
\Cnd{} Given a same-carrier semiclassical continuation certificate of
Definition~\ref{def:semiclassical-continuation-certificate} for a finite sequence
$x_0,\dots,x_n$, with $x_0$ ordinary matter and $x_n$ the target, the net change
$g(x_n)-g(x_0)$ is the path-ordered composition of the certified in-regime
increments \eqref{eq:coupling}. The endpoint may satisfy
$x_n\notin\mathcal R_{\rho}(x_0)$ while every successive pair lies in overlapping
local regimes certified at its actual basepoint. \textbf{Logical form:} \Cnd{} on
the existence of a certificate. \textbf{Composition verdict:} the continuation
grammar is a checked component of $\GCcert_{B_\Lambda}$; the full leverage verdict
is derived by Theorem~\ref{thm:verify-gravitational-leverage-gcert}.
\textbf{Continuation:} the certificate's existence, accumulated-error record,
semiclassical preservation conditions, and named break frontier. The construction uses the standard reachability grammar in
which ordered admissible generators build a finite transformation
\cite{S_JurdjevicSussmann1972}, with every step re-linearized at its registered
basepoint. It asserts integrability along a certified path, not a reversal of time
or force sign.
\end{sproposition}

\begin{sdefinition}[Same-carrier semiclassical continuation certificate]
\label{def:semiclassical-continuation-certificate}
\Defn{} For a declared carrier $\rho$, resolution $\delta$, and path
$\gamma=(x_0,\ldots,x_n)$, define
\begin{equation}
\mathsf{Cont}^{\mathrm{sc}}_{\rho,\delta}(\gamma)=
\left(
\rho,\gamma,\{\mathcal R_{\rho}(x_k)\},\{\mathcal O_{k,k+1}\},
\{\mathcal T_{k+1\leftarrow k}\},\{\epsilon_k\},
\mathcal M^{\mathrm{sc}}_{\rho},\mathsf{Break}_{\gamma}
\right).
\label{eq:semiclassical-continuation-certificate}
\end{equation}
Here $\mathcal O_{k,k+1}\subseteq\mathcal R_{\rho}(x_k)\cap
\mathcal R_{\rho}(x_{k+1})$ is an overlap relation, $\mathcal T_{k+1\leftarrow k}$
is a typed transport of the local response and registration records across that
overlap, $\epsilon_k$ is the recorded approximation error, and
$\mathcal M^{\mathrm{sc}}_{\rho}$ records preservation of the declared
semiclassical and code-subspace conditions. The break record includes loss of the
declared carrier, loss of overlap, loss of reconstruction, loss of semiclassicality,
divergent error, singularity, or any failed next response relation. The certificate
therefore distinguishes an endpoint outside the initial chart from a regime
violation: the former is compatible with a chain of overlapping certified charts;
the latter is a typed break recorded inside the composition.
\end{sdefinition}


\begin{sdefinition}[Globally coherent gravitational leverage composition]
\label{def:gravitational-leverage-composition}
\Hyp{} For declared target, resolution, and coupled P-space/A(R)-space regime,
define the gravitational leverage composition
\begin{equation}
\Lambda^{\mathrm{grav}}_\delta=
\left(
\mathsf{Enc}_{\mathrm{ext}},\
\mathsf{Lin}_{\rho,\delta},\
\mathsf{Cont}^{\mathrm{sc}}_{\rho,\delta}(\gamma),\
\mathsf{Sched}_{\star},\
\mathsf{Self}_{\mathrm{grav}},\
\mathsf{Adm}_{P/A,\delta},\
\mathsf{Haz}_{\delta},\
\mathsf{Prod}_{\delta}
\right).
\label{eq:gravitational-leverage-composition}
\end{equation}
Here $\mathsf{Enc}_{\mathrm{ext}}$ is the declared extended encoding support,
$\mathsf{Lin}_{\rho,\delta}$ is the carrier-indexed local
entanglement--geometry response relation, and
$\mathsf{Cont}^{\mathrm{sc}}_{\rho,\delta}(\gamma)$ is the same-carrier
continuation certificate replacing a bare re-linearization assertion. The remaining
coordinates are the target-conditioned admissible schedule,
gravitational self-coupling coordinate inherited from the metric/constraint carrier,
paired admissibility and bridge conditions, hazard/runaway record, and the declared
cumulative-growth coordinate
\begin{equation}
\mathsf{Prod}_{\delta}:\qquad \frac{G_n}{n}\longrightarrow\infty.
\label{eq:productivity-coordinate}
\end{equation}
The source objects are Cards~\ref{card:faulkner-linear},
\ref{card:faulkner-nonlinear}, \ref{card:jlms}, and~\ref{card:qec},
Definition~\ref{def:semiclassical-continuation-certificate},
Proposition~\ref{prop:integration}, the target-conditioned schedule of
Definition~\ref{def:target-conditioned-schedule}, and the gravity-native carrier
of Definition~\ref{def:second-order-gravitation-theory}. The definition does not
assign its own global-coherence verdict. Its productivity coordinate selects the
branch on which the cumulative response supports the target-depth consequence; the
admissibility and hazard coordinates retain the corresponding physical and safety
relations. The in-budget verdict is derived only after the explicit certificate
below is verified.
\end{sdefinition}

\begin{sdefinition}[Finite certificate for the gravitational leverage composition]
\label{def:gravitational-leverage-gcert}
\Defn{} Let $B_\Lambda$ be the finite assessment budget that contains exactly:
(i)~the claim scope that $\Lambda^{\mathrm{grav}}_\delta$ is a conditional
same-carrier leverage hypothesis and makes no assertion that a certified path,
productivity branch, or target realization has been empirically instantiated;
(ii)~the cited holographic source cards
\ref{card:faulkner-linear}--\ref{card:qec}, the gravity-native comparison carrier,
and the control grammar; (iii)~the common-carrier check
$\rho=\rho_{\mathrm{AdS/CFT}}$ for the holographic path; (iv)~the typed bridge
checks in Card~\ref{card:entanglement-junction}; (v)~the regime and continuation
checks of Definition~\ref{def:semiclassical-continuation-certificate};
(vi)~the non-erasure check that the certificate-existence, productivity, empirical,
and break branches remain recorded rather than substituted by a conclusion;
and (vii)~the named break and defeat fields. Define
\begin{equation}
\GCcert_{B_\Lambda}\!\left(\Lambda^{\mathrm{grav}}_\delta\right)=
\left(
\begin{array}{l}
\mathsf{ClaimScope}_{\mathrm{cnd}},\ 
\mathsf{Imports}_{\rho_{\mathrm{AdS/CFT}}},\ 
\mathsf{CarrierCheck}_{\rho},\\
\mathsf{BridgeCheck}_{\mathrm{EG}},\ 
\mathsf{RegimeCheck}_{\mathrm{sc}},\ 
\mathsf{NonErasureCheck}_{\mathrm{res}},\\
\mathsf{DefeatAudit}_{\mathrm{named}},\ 
\mathsf{Residue}_{\exists\gamma,\mathrm{Prod},\mathrm{Emp}}
\end{array}
\right).
\label{eq:gravitational-leverage-gcert}
\end{equation}
The continuation branch is present as a condition and residue record; it is not
entered as a realized path. Likewise, the productivity coordinate and empirical
instantiation remain explicit unadjudicated branches of the conditional hypothesis.
\end{sdefinition}

\begin{stheorem}[Verified in-budget global-coherence verdict for $\Lambda^{\mathrm{grav}}_\delta$]
\label{thm:verify-gravitational-leverage-gcert}
\Thm{} Under the source and bridge records named in
Definition~\ref{def:gravitational-leverage-gcert},
\begin{equation}
\GCverify_{B_\Lambda}\!\left(
\GCcert_{B_\Lambda}\!\left(\Lambda^{\mathrm{grav}}_\delta\right)
\right)=\checkmark
\quad\Longrightarrow\quad
\mathsf{GC}_{B_\Lambda}\!\left(\Lambda^{\mathrm{grav}}_\delta\right)=\mathsf{pass}.
\label{eq:leverage-gc-pass}
\end{equation}
\end{stheorem}

\begin{proof}
The claim-scope field contains only the conditional leverage hypothesis, not an
assertion that $\exists\gamma\,\mathsf{Cont}^{\mathrm{sc}}_{\rho,\delta}(\gamma)$,
$\mathsf{Prod}_\delta$, or empirical realization already holds. The import field
keeps the holographic response, continuation-prefix, transport, and reconstruction
objects on their declared common carrier; Jacobson's independent local
semiclassical family is not substituted for that carrier. The bridge field points
to the explicit junction record, and the regime field retains overlap, transport,
error, code-subspace, and break conditions instead of extending one local kernel
outside its regime. The non-erasure field retains each unresolved branch as
residue. Finally, the defeat audit contains the named carrier, overlap,
reconstruction, semiclassicality, error, singularity, and next-response failures;
within $B_\Lambda$ none contradicts the conditional claim actually made. Each
finite check in Definition~\ref{def:global-coherence-certificate} is therefore
satisfied, so its verification returns $\checkmark$. The conclusion follows from
Definition~\ref{def:assessment-budget}. This theorem is an in-budget coherence
verdict for the conditional composition; it neither realizes a path nor supplies a
quantitative gravitational prediction.
\end{proof}

\begin{sproposition}[Recursive leverage and target depth]\label{prop:amplify}
\Cnd{} Let the globally coherent leverage composition
$\Lambda^{\mathrm{grav}}_\delta$ of
Definition~\ref{def:gravitational-leverage-composition} have non-negative
admissible increments $\delta S_{\mathrm{EE},k}$ and cumulative response magnitude
\begin{equation}
G_n=\sum_{k=0}^{n-1}\mathcal K(x_k)\,\delta S_{\mathrm{EE},k}.
\label{eq:cumulative-response}
\end{equation}
On the productivity coordinate $\mathsf{Prod}_{\delta}$, namely
$G_n/n\to\infty$ along the declared direction of travel, the least step count
$N(\Delta)=\inf\{n:G_n\geq\Delta\}$ satisfies
\begin{equation}
\frac{N(\Delta)}{\Delta}\longrightarrow 0
\qquad (\Delta\to\infty),
\end{equation}
so target depth is sub-linear in the declared response size. The leverage
composition therefore contains a named conditional target-depth consequence of its
productivity coordinate. Its gravitational self-coupling coordinate,
$\mathsf{Self}_{\mathrm{grav}}$, carries the source relation through which the
response geometry enters the next registered step; the admissibility and hazard
coordinates retain the in-regime and runaway branches under
Obligation~\ref{ob:hazard}. \textbf{Type:} \Cnd{} within the \Hyp{} globally
coherent composition of Equation~\eqref{eq:gravitational-leverage-composition}.
\end{sproposition}


%=================================================================
\section{Recursive interaction-network gain: selected physical witness families}
\label{supp:recursive-interaction-network-gain}

\paragraph{Purpose and claim scope.}
The productivity coordinate of the generic holographic leverage composition is
not identified here with a freely chosen graph count or with an all-regime
entanglement--geometry kernel.  This section supplies selected physical witness
families in which a registered gain relation follows from declared governing
relations within a finite regime.  The witnesses serve a positive but bounded
role: they establish that recursive gravitational gain is physically nonempty,
that it is not confined to black-hole geometry, and that the gain relation can
be made explicit without silently transferring a result between carriers.

The section distinguishes three objects:
\begin{enumerate}[label=(R\arabic*)]
\item a finite GHZ-tree witness for target-relative paired-observability support;
\item a selected shock-wave \emph{gain family} with a registered eikonal shift;
\item a selected non-black-hole Jeans carrier with an actual forward-time
      self-gravitating feedback loop.
\end{enumerate}
Only the third is used below as a forward-time recursive physical-network
witness.  The shock family remains a quantitative gravitational gain witness,
and the GHZ tree remains an information-theoretic non-emptiness witness.  No
one of these objects is substituted for an arbitrary-target metric-control
claim or for the distinct generic holographic continuation branch.

\subsection{Recursive interaction-network grammar}
\label{supp:recursive-network-definition}

\begin{sdefinition}[Recursive interaction network]
\label{def:recursive-interaction-network}
\Defn{} For a declared carrier $\rho$, a recursive interaction network is
\begin{equation}
\mathfrak N_\rho=
\left(
X^P_\rho,\;X^A_\rho,\;\Reg_\rho,\;\Bridge_\rho,\;
\{\Phi_k\}_{k\ge0},\;\Theta_\rho,\;q_\rho,\;\mathcal D_\rho
\right),
\label{eq:recursive-interaction-network}
\end{equation}
where $X^P_\rho$ is a physical state space, $X^A_\rho$ a registration space,
$\Reg_\rho$ a typed physical-to-registration map, $\Bridge_\rho$ the
target-preserving bridge, $\Phi_k$ admissible forward-time physical updates,
$\Theta_\rho$ a declared finite target predicate, $q_\rho\ge0$ a registered
target coordinate, and $\mathcal D_\rho$ the declared carrier/regime domain.
The network is gain-bearing when, for every retained step in its domain,
\begin{equation}
q_\rho\!\left(\Reg_\rho(\Phi_k(x_k))\right)
\ge\lambda_\rho\,q_\rho\!\left(\Reg_\rho(x_k)\right),
\qquad \lambda_\rho>1.
\label{eq:physical-gain-recurrence}
\end{equation}
The factor must follow from a carrier relation or exact update law; it may not
be introduced only by selecting an arity, a counting convention, or a target
label.
\end{sdefinition}

\begin{sdefinition}[Selected realization certificate]
\label{def:selected-realization-certificate}
\Defn{} A selected realization certificate is
\begin{equation}
\RCert_\rho=
\left(
\rho,\;\mathcal D_\rho,\;x_0,\;\Theta_\rho,\;q_\rho,\;
\{\Phi_k\},\;\Reg_\rho,\;\Bridge_\rho,\;
\lambda_\rho,\;q_{\max},\;\mathsf{Defeat}_\rho
\right),
\label{eq:selected-realization-certificate}
\end{equation}
recording the carrier and governing equations, initial state, finite target,
in-domain route, target registration, derived gain relation, finite validity
bound, and named defeat conditions.
\end{sdefinition}

\begin{stheorem}[Finite-target gain theorem]
\label{thm:finite-target-gain}
\Thm{} Suppose a selected certificate verifies and
\begin{equation}
q_{k+1}\ge\lambda_\rho q_k,\qquad
\lambda_\rho>1,\qquad
0<q_0\le q_\star\le q_{\max}.
\label{eq:finite-target-gain-assumptions}
\end{equation}
Then the selected target threshold is reached within
\begin{equation}
N_\rho(q_\star)
\le
\left\lceil\frac{\log(q_\star/q_0)}{\log\lambda_\rho}\right\rceil
\label{eq:finite-target-depth}
\end{equation}
steps.  Corridor, continuation, productivity, and registration pass for this
selected finite target within the certificate's carrier and regime scope.
\end{stheorem}

\begin{proof}
Iteration gives $q_n\ge\lambda_\rho^nq_0$.  Solving
$\lambda_\rho^nq_0\ge q_\star$ gives
Equation~\eqref{eq:finite-target-depth}.  The certificate retains one
in-domain update at each step and a registration on which the same target
predicate is defined.
\end{proof}

\paragraph{Spiral interpretation.}
The recursion is a nonliteral spiral P-space traversal:
\begin{equation}
x_k\longrightarrow\Reg_\rho(x_k)
\longrightarrow\Bridge_\rho\!\left(\Reg_\rho(x_k)\right)
\longrightarrow\Phi_k(x_k)=x_{k+1}.
\label{eq:spiral-state-update}
\end{equation}
It is not a spiral in spacetime.  It denotes repeated typed passes through a
changed state/constraint topology: the present physical state changes a later
physical update or the registered conditions under which that update is
admissible.  It therefore differs from replaying one fixed operation on one
fixed topology.

\subsection{Discrete target-relative observability witness}
\label{supp:ghz-tree-witness}

\paragraph{Scope.}
The GHZ-tree record establishes that the paired target-observability schema has
an exact nonempty finite realization.  It is not a gravitational-amplification
claim.

\begin{sdefinition}[GHZ-tree target pair]
\label{def:ghz-tree-target-pair}
\Defn{} Fix $b\ge2$ and $N_n=b^n$.  Define
\begin{equation}
\left|\mathrm{GHZ}^{\pm}_n\right\rangle=
\frac{|0\rangle^{\otimes N_n}\pm|1\rangle^{\otimes N_n}}{\sqrt2}.
\label{eq:ghz-pm}
\end{equation}
Let $\Theta^{\rm GHZ}_n$ be the relative global phase.  The minimum
support that distinguishes the target is
\begin{equation}
d_{\Theta}(n)=
\min\left\{|S|:\exists O_S\text{ supported on }S\text{ with }
\langle O_S\rangle_+\ne\langle O_S\rangle_-\right\}.
\label{eq:ghz-target-support}
\end{equation}
\end{sdefinition}

\begin{sproposition}[Exact target-relative support recurrence]
\label{prop:ghz-support-recurrence}
\Thm{} For this target pair,
\begin{equation}
d_{\Theta}(n)=N_n=b^n,
\qquad d_{\Theta}(n+1)=b\,d_{\Theta}(n).
\label{eq:ghz-support-recurrence}
\end{equation}
Every proper reduced subsystem is identical for the two target states, whereas
$X^{\otimes N_n}$ distinguishes them.  The declared tree isometry maps
$|\mathrm{GHZ}^{\pm}_n\rangle$ to
$|\mathrm{GHZ}^{\pm}_{n+1}\rangle$ exactly.
\end{sproposition}

\begin{proof}
Tracing out at least one qubit removes the off-diagonal phase term; hence every
proper reduced state agrees between the two GHZ states.  The global Pauli
observable has expectation $+1$ and $-1$, respectively.  The tree isometry
preserves both computational-basis strings and their relative sign, yielding
the stated recurrence.
\end{proof}

\paragraph{Status.}
The recurrence is target-relative because it measures the minimum support
needed to distinguish a declared phase target rather than bare leaf count.  It
still establishes only a discrete information-theoretic witness; it does not
transfer a force, energy, or continuum response claim to gravity.

\subsection{Selected shock-wave gain family}
\label{supp:shock-realization}

\begin{scard}[Shock-wave gain source object]
\label{card:shock-gain-source}
\Imp{} In a declared nonextremal two-sided AdS black-hole background, an early
null perturbation in the standard shock-wave/eikonal regime produces a
near-horizon shift of the form
\begin{equation}
\alpha(\tau)=\alpha_0e^{\kappa\tau},
\qquad \alpha_0=C_{\rm sh}G_NE>0,
\label{eq:shock-shift}
\end{equation}
where $\tau$ is the positive early-time separation, $\kappa$ is the surface
gravity, and $C_{\rm sh}$ depends on the selected background and shell profile.
Shock-wave constructions and their holographic butterfly analysis provide this
relation in the stated regime~\cite{S_DraytHooft1985,S_ShenkerStanford2014,S_ShenkerStanford2014Multiple}.
\end{scard}

\begin{sdefinition}[Selected shock gain family]
\label{def:selected-shock-realization}
\Defn{} Let $\mathcal D_{\rm sh}$ be a finite shock/eikonal domain and fix
\begin{equation}
0<\alpha_0\le\alpha_\star\le\alpha_{\max}(\mathcal D_{\rm sh}).
\label{eq:shock-target-domain}
\end{equation}
The target is the registered relational null shift:
$\Theta_{\rm sh}(\alpha)=1$ iff $\alpha\ge\alpha_\star$.  Take
$\tau_k=k\Delta\tau$.  The sequence labels a declared family of early-time
preparations and their forward-time readouts; it is a quantitative gravitational
gain family, not an assertion that one black-hole experiment self-iterates an
arbitrary metric-control operation.
\end{sdefinition}

\begin{sproposition}[Quantitative shock-wave gain witness]
\label{prop:shock-wave-gain-realization}
\Thm{} Within $\mathcal D_{\rm sh}$,
\begin{equation}
\alpha_{k+1}=e^{\kappa\Delta\tau}\alpha_k,
\qquad \lambda_{\rm sh}=e^{\kappa\Delta\tau}>1,
\label{eq:shock-recurrence}
\end{equation}
and the selected finite threshold is reached by
\begin{equation}
N_{\rm sh}(\alpha_\star)=
\left\lceil\frac{\log(\alpha_\star/\alpha_0)}{\kappa\Delta\tau}\right\rceil.
\label{eq:shock-target-depth}
\end{equation}
The eikonal/horizon-shift record registers the target coordinate.  Thus the
selected shock family supplies a quantitative gravitational gain witness within
its declared domain.
\end{sproposition}

\begin{proof}
Equation~\eqref{eq:shock-recurrence} is immediate from
Equation~\eqref{eq:shock-shift}; solving the threshold inequality yields
Equation~\eqref{eq:shock-target-depth}.  The stated domain and the shift record
supply the carrier and registration checks for this selected family.
\end{proof}

\paragraph{Scope.}
The shock witness does not establish unrestricted metric control, a universal
gain law, or a silent transfer to the generic holographic continuation carrier.
It is defeated outside its specified shock/eikonal domain or if its selected
shift registration is lost.

\subsection{Selected non-black-hole self-gravitating-fluid network}
\label{supp:jeans-realization}

\begin{scard}[Jeans feedback source object]
\label{card:jeans-feedback-source}
\Imp{} For a homogeneous barotropic self-gravitating fluid with equilibrium
density $\rho_0$ and sound speed $c_s$, the linearized continuity, Euler, and
Poisson equations are
\begin{align}
\partial_t\delta\rho+\rho_0\nabla\!\cdot\delta\mathbf v&=0,
\label{eq:jeans-continuity}\\
\partial_t\delta\mathbf v&=-\frac{c_s^2}{\rho_0}\nabla\delta\rho
-\nabla\delta\Phi,
\label{eq:jeans-euler}\\
\nabla^2\delta\Phi&=4\pi G\,\delta\rho.
\label{eq:jeans-poisson}
\end{align}
For a plane mode of wavenumber $k$,
\begin{equation}
\omega^2=c_s^2k^2-4\pi G\rho_0.
\label{eq:jeans-dispersion}
\end{equation}
For $k<k_J:=\sqrt{4\pi G\rho_0}/c_s$, put
$\gamma_k:=\sqrt{4\pi G\rho_0-c_s^2k^2}>0$; the growing mode has
$A(t)=A_0e^{\gamma_kt}$.  These relations follow from the coupled
Euler--Poisson source system~\cite{S_Jeans1902,S_BinneyTremaine2008,S_Chavanis2010}.
\end{scard}

\begin{sdefinition}[Selected Jeans realization]
\label{def:selected-jeans-realization}
\Defn{} Let $\rho_{\rm J}^{\rm lin}$ be the declared linearized Euler--Poisson
carrier for one unstable Fourier mode $k<k_J$.  Choose
$0<A_0\le A_\star\le A_{\max}\ll\rho_0$ and define
\begin{equation}
\Theta_{\rm J}(A)=1\quad\Longleftrightarrow\quad A\ge A_\star.
\label{eq:jeans-target-predicate}
\end{equation}
The registration is the signed density-mode amplitude $A$ and the associated
potential mode
\begin{equation}
\delta\Phi_k=-\frac{4\pi G}{k^2}A.
\label{eq:jeans-registration}
\end{equation}
\end{sdefinition}

\begin{stheorem}[Non-black-hole gravitational feedback realization]
\label{thm:jeans-gain-realization}
\Thm{} Within the selected linear carrier, the sequence at
$t_k=k\Delta t$ obeys
\begin{equation}
A_{k+1}=e^{\gamma_k\Delta t}A_k,
\qquad \lambda_{\rm J}=e^{\gamma_k\Delta t}>1,
\label{eq:jeans-recurrence}
\end{equation}
and reaches the finite threshold after at most
\begin{equation}
N_{\rm J}(A_\star)=
\left\lceil\frac{\log(A_\star/A_0)}{\gamma_k\Delta t}\right\rceil.
\label{eq:jeans-target-depth}
\end{equation}
The physical feedback circuit is
\begin{equation}
\delta\rho_k
\xrightarrow{\;\nabla^2\delta\Phi_k=4\pi G\delta\rho_k\;}
\delta\Phi_k
\xrightarrow{\;-\nabla\delta\Phi_k\;}
\delta\mathbf v_{k+1}
\xrightarrow{\;-\rho_0\nabla\!\cdot\delta\mathbf v_{k+1}\;}
\delta\rho_{k+1}.
\label{eq:jeans-feedback-circuit}
\end{equation}
Therefore $\RCert_{\rm J}$ passes corridor, continuation, productivity, and
registration for the declared finite density-contrast target within
$\rho_{\rm J}^{\rm lin}$.
\end{stheorem}

\begin{proof}
The source equations combine to
\begin{equation}
\partial_t^2A+\left(c_s^2k^2-4\pi G\rho_0\right)A=0.
\label{eq:jeans-mode-equation}
\end{equation}
For $k<k_J$, the selected growing solution gives
Equation~\eqref{eq:jeans-recurrence}.  The target remains in the declared
linear carrier by $A_\star\le A_{\max}$, and the amplitude/potential pair
provides the target registration.  The depth estimate follows from
Theorem~\ref{thm:finite-target-gain}.
\end{proof}

\paragraph{Interpretation.}
This is a composition of fluid mechanics and Newtonian self-gravity, not a
black-hole construction.  A density perturbation changes potential; potential
changes velocity convergence; convergence changes the next density perturbation.
The gain factor is fixed by the physical dispersion relation rather than by a
chosen graph arity.  The witness is an instability realization, not a proof of
arbitrary target-directed gravitational control; its scope ends at $A_{\max}$.

\subsection{Witness-family theorem and terminal realization record}
\label{supp:witness-family-theorem}

\begin{stheorem}[Nonempty physical gravitational-gain class]
\label{thm:nonempty-gravitational-gain-class}
\Thm{} The class of selected physical gravitational gain objects is nonempty.
It contains the shock gain family of
Proposition~\ref{prop:shock-wave-gain-realization} and the non-black-hole
forward-time recursive feedback network of
Theorem~\ref{thm:jeans-gain-realization}.  The gain factors are
\begin{equation}
\lambda_{\rm sh}=e^{\kappa\Delta\tau}>1,
\qquad
\lambda_{\rm J}=e^{\gamma_k\Delta t}>1.
\label{eq:two-physical-gains}
\end{equation}
For the declared finite target windows,
\begin{equation}
\mathsf{GC}_{B_{\rm sh}}(\mathfrak F_{\rm sh})=\mathsf{pass},
\qquad
\mathsf{GC}_{B_{\rm J}}(\mathfrak N_{\rm J})=\mathsf{pass},
\label{eq:two-gc-passes}
\end{equation}
where each budget verifies carrier, governing relation, target, in-domain
physical family/path, registration, gain relation, and named defeat conditions.
\end{stheorem}

\begin{proof}
The cited shock relation gives the first selected gain family; the coupled
Euler--Poisson derivation gives the second selected forward-time feedback
network.  In both cases the gain relation is a consequence of stated physical
relations within a finite domain, rather than a free asymptotic premise.  The
listed certificate fields are thereby verified for their declared scope.
\end{proof}

\begin{compositioncard}[Selected gravitational realization record]
\label{card:terminal-realization-record}
\Defn{}
\textbf{Claim scope.} There exists a nonempty class of selected physical
gravitational systems carrying a quantitatively derived registered gain relation;
at least one selected realization is a non-black-hole forward-time recursive
self-gravitating feedback network.

\textbf{Established source objects.} The shock-wave source gives an exponential
near-horizon shift in a stated black-hole regime
\cite{S_DraytHooft1985,S_ShenkerStanford2014}; the Jeans source gives an
exponential unstable density mode in a stated self-gravitating-fluid regime
\cite{S_Jeans1902,S_BinneyTremaine2008,S_Chavanis2010}.

\textbf{Composed result.} The two families realize the shared
state--registration--gain grammar of
Definition~\ref{def:recursive-interaction-network}; the Jeans member realizes
the forward-time feedback circuit explicitly.

\textbf{Status.}
\begin{equation}
\left(
\textsc{[thm]},\;
\textsc{[globally coherent]},\;
\textsc{[closed selected physical realizations]}
\right).
\label{eq:terminal-realization-status}
\end{equation}

\textbf{Not entailed.} The record does not entail unrestricted control of an
arbitrary metric target, a universal gain law for all gravitating systems, or a
silent transfer to the generic holographic continuation carrier.
\end{compositioncard}

\paragraph{Relation to the holographic path.}
The entanglement/geometry cards remain a distinct carrier family.  Their local
and perturbative scope is not asked to bear the whole terminal physical burden.
The selected witnesses establish nonempty physical gain and a non-black-hole
feedback realization; a transfer to a selected holographic carrier remains
licensed only by its own bridge preserving target coordinate, regime, and error
conditions.  This separation prevents the physical witnesses from being mistaken
for an unproved universal AdS/CFT or metric-gravity claim.


%=================================================================
\section{Non-signalling of the composed channel}
\label{supp:nosignal}

\begin{sproposition}[Stepwise and composed non-signalling]\label{prop:nosignal}
\Thm{}/\Cnd{} Let each preparation step be a CPTP operation on a region of the
encoding support, coordinated by subluminal classical communication, with
reconstruction wedge-local in the sense of Card~\ref{card:qec}. Then (a)~each step
is non-signalling by Card~\ref{card:nosig}(i); (b)~the composition of such steps
is non-signalling; and (c)~no step introduces nonlinear state evolution, so the
Gisin--Polchinski route of Card~\ref{card:nosig}(ii) is not opened. Any inference
from entanglement-wedge reconstruction locality to confinement of a P-space effect
within a causal wedge is a separate bridge-required conditional under
Obligation~\ref{ob:wedge}; it is not supplied by Card~\ref{card:qec} alone.
\textbf{Type:} (a)--(c) \Thm{} relative to the imports; causal-wedge confinement
is \Cnd{} on the additional bridge. Non-signalling is therefore established within
the declared regime and model class, not unconditionally.
\end{sproposition}

\begin{sproposition}[Direct channel as degenerate limit]\label{prop:degenerate}
\Defn{}/\Thm{} A carrierless, nonlinear ``direct'' channel is the limit of the
composed channel in which the actuator composition collapses to a single
zero-support nonlinear operation. Card~\ref{card:nosig}(ii) applies to that limit;
hence the no-go obstruction constrains the degenerate limit, not the composed
channel, which is the coherent object.
\end{sproposition}

%=================================================================
\section{The two propositions of the main article, formally}
\label{supp:props}

\begin{sproposition}[Insufficiency of derive-and-test]\label{prop:insufficiency-supp}
\Thm{} Let a constraint $C$ be a predicate on the admissibility of a composition
$\gamma = (x_0,\dots,x_n)$---for example, ``a sequence $\gamma$ exists with all
$x_k$ admissible and all steps in-regime''---such that no declared local-observation
record decides $C$. Then $C$ cannot be evaluated by deriving an observable from a
first-order model of an individual step and testing it. Any sound and complete
evaluation requires a representation retaining the distinctions on which
admissibility depends. Proposition~\ref{prop:strict-separation-witness} supplies a
finite witness for the non-empty case in which the complete proper-support local
record is identical while $C$ differs. \textbf{Failure of the paradigm, not of
measurement:} the obstruction is that $C$ has no adequate local witness, so the
absence of a deciding measurement is structural.
\end{sproposition}

\begin{sproposition}[Validation--hazard coupling]\label{prop:coupling-supp}
\Cnd{} Suppose the worked-case path refers: a corridor exists
(Obligation~\ref{ob:composition}), the cumulative leverage-growth branch of
Proposition~\ref{prop:amplify} operates, and the current assessment record contains
no adequate safe non-traversal witness (Obligation~\ref{ob:safewitness}). Then
traversal is the currently licensed validation procedure, and it can light a
positive feedback on a globally supported geometric variable whose containment,
reversibility, and stopping conditions have not been established
(Obligation~\ref{ob:hazard}). This does not claim that no safe non-traversal route
exists in reality; a subsequently supplied adequate certificate changes the
assessment branch. Any causal-wedge confinement argument is separately
bridge-required (Obligation~\ref{ob:wedge}). Hence the current empirical
validation branch is coupled to the hazard the constraint describes. \textbf{Type:}
\Cnd{} on the antecedents. Global support alone does not establish
uncontrollability; the proposition applies only when the required containment and
safe-witness conditions remain open in the record. If the antecedent fails, there
is no established hazard; the proposition records the unresolved question without
converting it into an empirical claim.
\end{sproposition}


%=================================================================
\section{Strict separation and second-order non-reducibility}
\label{supp:strict-separation}

This section strengthens the necessity argument of the main article in two
ways. First, it gives a finite typed witness showing that two candidate
compositions can agree on every declared local observable while differing on
corridor existence. The result is therefore not merely that a presently known
local experiment is insufficient; it is that every assessment procedure which
factors through the specified local-observation quotient is insufficient.

Second, it proves a qualified non-reducibility result. For the declared class of
global-coherence constraints, any sound and complete decision procedure must
retain distinctions extensionally equivalent to: (i)~the typing and legality of
the P-space/A(R)-space bridge; (ii)~the regime record; (iii)~the ordered history
of the path; and (iv)~the support structure and global invariant on which the
composed constraint depends. This is the corridor-specific instance of the general
fiber necessity theorem (Theorem~\ref{arch:thm-fiber}): a representation cannot
identify two records when a required audit changes between them. It does not claim
that every physics problem requires one literal notation, that no safe global
witness can ever exist, or that a single view retains all physical information. It
states a narrower result: a representation which erases any of these distinctions
cannot decide this constraint class soundly and completely.


\paragraph{Source antecedent and typed deployment.}
The proper-support separation pattern used below has familiar information-theoretic
antecedents: threshold constructions exhibit target-relevant information that is
absent from designated proper subsets~\cite{S_Shamir1979}. The present deployment
uses that structural pattern as a source object and then fixes the distinct
corridor predicate, bridge, regime, history, support, and safe-witness coordinates
that define the physical constraint record. The theorem therefore establishes the
specified typed fiber condition relative to its witness class.

%-----------------------------------------------------------------
\paragraph{Representation-adequacy interface (orientation only).}
The following construction is a physical instance of a general
representation-adequacy schema. Let $X_{\mathrm{SOP}}$ be the typed
composition-instance space, let
$H_{\mathrm{loc}}:X_{\mathrm{SOP}}\to R_{\mathrm{loc}}$ be the declared
proper-support local-observation quotient, and let
$\Theta_{\mathrm{SOP}}$ be the corridor predicate together with its required
audit signature. A pair satisfying
\[
H_{\mathrm{loc}}(x)=H_{\mathrm{loc}}(y)
\qquad\text{and}\qquad
\Theta_{\mathrm{SOP}}(x)\neq\Theta_{\mathrm{SOP}}(y)
\]
certifies that the predicate does not factor through the quotient on the declared
witness class. The residue is not an untyped warning: an adequate corridor record
must retain the bridge, regime, ordered-history, support, and safe-witness
distinctions that the quotient erased. This is an interface annotation only; the
definitions and proofs below establish the SOP result directly.

\subsection{Typed composition instances and the corridor predicate}
\label{supp:typed-composition-instance}

\begin{sdefinition}[Typed composition instance]
\label{def:typed-composition-instance}
\Defn{} A \emph{typed composition instance} is a finite record
\begin{equation}
\mathfrak C
=
\left(
\Pspace,\ARspace,\tau,\mathcal B,
\{\mathcal R_k\}_{k=0}^{n},
\gamma,S,\mathcal U_{\mathrm{loc}},
\mathcal O_{\mathrm{loc}},
\mathsf{Adm},\mathsf{Amp},\mathsf{SafeWit}
\right),
\label{eq:typed-composition-instance}
\end{equation}
where:
\begin{enumerate}
\item $\Pspace$ is the declared physical-state representation and
$\ARspace$ is the declared registration-state representation;

\item $\tau$ is the role and type assignment, including the distinction
between P-space objects, A(R)-space objects, operators, boundaries,
registration records, and physical configurations;

\item $\mathcal B$ is the family of typed bridges permitted between declared
objects, including any bridge through which a registration preparation can
become a physical geometric update;

\item $\gamma$ is an ordered candidate path
\begin{equation}
\gamma
=
\left(
x_0
\xrightarrow{u_1}
x_1
\xrightarrow{u_2}
\cdots
\xrightarrow{u_n}
x_n
\right),
\label{eq:ordered-path}
\end{equation}
with $x_0$ the ordinary initial configuration and $x_n$ a declared target
configuration;

\item $\mathcal R_k$ is the regime record at step $k$, including the
conditions under which every imported coupling, reconstruction property,
non-signalling statement, and safety argument remains licensed;

\item $S$ is the declared support of the registration structure relevant to
the candidate bridge; $\mathcal U_{\mathrm{loc}}$ is the family of supports
treated as local by the first-order assessment procedure; and
$\mathcal O_{\mathrm{loc}}$ is the family of observations admissible on those
supports;

\item $\mathsf{Adm}$ is the composition-level admissibility predicate,
including the preservation of types, bridges, regime conditions, support
conditions, ordered retention conditions, and any gravity-native barrier
condition;

\item $\mathsf{Amp}$ records whether the declared path enters the candidate
self-amplifying sector; and $\mathsf{SafeWit}$ records whether an adequate safe
non-traversal witness for the relevant global predicate has been supplied to the
current assessment record. Its value $0$ records no positive certificate in that
record; it does not assert the global non-existence of a possible future witness.
\end{enumerate}
\end{sdefinition}

The point of the record is not that every physical theory must be written in
this notation. It is that a decision about a composed constraint requires a
representation in which the distinctions on which that decision depends have
not already been erased. In the calculus of Section~\ref{supp:calculus}, this
record is the declared physical and relational payload of $R_{\mathfrak C}$; it
does not expand the operator basis beyond $\{\Register,\Recall,\GOp,\LOp\}$.

\begin{sdefinition}[Corridor predicate]
\label{def:corridor-predicate}
\Defn{} For a typed composition instance $\mathfrak C$, define
\begin{equation}
\mathsf{Corr}(\mathfrak C)=1
\label{eq:corridor-predicate}
\end{equation}
if and only if there exists a declared path $\gamma$ from the ordinary
configuration $x_0$ to the target $x_n$ such that
\begin{equation}
\mathsf{Adm}(\gamma)=1,
\qquad
x_k,x_{k+1}\in\mathcal R_k
\quad
\text{for all }k,
\label{eq:corridor-admissibility}
\end{equation}
every required bridge in $\mathcal B$ is type-valid, every required retention
condition is met, and the required support condition on $S$ is attained.
Otherwise $\mathsf{Corr}(\mathfrak C)=0$.

The predicate states only whether the declared route is open within the
declared model class. It does not establish that the model class is realized in
nature, that the target is physically reachable, or that a geometric
modification has been empirically demonstrated.
\end{sdefinition}

For later use, define a conditional validation--hazard flag:
\begin{equation}
\mathsf{VHC}(\mathfrak C)
=
\mathsf{Corr}(\mathfrak C)
\wedge
\mathsf{Amp}(\mathfrak C)
\wedge
\neg\mathsf{SafeWit}(\mathfrak C).
\label{eq:validation-hazard-flag}
\end{equation}
This is a classification predicate, not an assertion that a real hazard has
been established. It records that, within the declared composition and assessment record, an open
amplifying route lacks a supplied safe non-traversal witness; it does not prove that
no such witness exists outside the record.

%-----------------------------------------------------------------
\subsection{Local-observation quotients}
\label{supp:local-quotient}

\begin{sdefinition}[Local-observation record]
\label{def:local-observation-record}
\Defn{} Let $S$ be a finite registration support and let
$\mathcal U_{\mathrm{loc}}$ be a family of proper supports
$U\subsetneq S$. For a registration distribution or state $\rho$ on $S$,
write $\pi_U$ for restriction to $U$. The \emph{local-observation record} of
a typed composition instance is
\begin{equation}
\Lambda_{\mathrm{loc}}(\mathfrak C)
=
\left(
\Pspace_{\mathrm{loc}},
(\ARspace)_{\mathrm{loc}},
\left\{
(\pi_U)_*\rho,\mathcal O_U
:
U\in\mathcal U_{\mathrm{loc}}
\right\},
\gamma_{\mathrm{loc}}
\right),
\label{eq:local-observation-record}
\end{equation}
where every $\mathcal O_U$ factors through the restriction $\pi_U$ and
$\gamma_{\mathrm{loc}}$ retains only the corresponding locally observable
step data. Two instances are \emph{locally equivalent}, written
\begin{equation}
\mathfrak C\equiv_{\mathrm{loc}}\mathfrak C',
\label{eq:local-equivalence}
\end{equation}
when their local-observation records are equal.
\end{sdefinition}

The definition allows the local procedure to collect and compare every declared
local record. It is therefore stronger than the claim that a single isolated
measurement is insufficient. What it excludes is evaluation of a full-support
invariant or a composition-level predicate that is not determined by the complete
family of proper-support records.

%-----------------------------------------------------------------
\subsection{A finite strict-separation witness}
\label{supp:finite-witness}

The following witness is intentionally finite and abstract. It is not a model of
actual gravitational microphysics, an assertion about holography, or a claim that
the worked-case path is physically realized. Its role is to prove the logical
point required here: a global composition predicate can vary while every declared
local observable remains exactly fixed.

Let
\begin{equation}
S=\{1,2,3\},
\qquad
\mathcal A_S=\{0,1\}^{3},
\label{eq:parity-support}
\end{equation}
and let the global parity functional be
\begin{equation}
\Pi(a_1,a_2,a_3)=a_1\oplus a_2\oplus a_3,
\label{eq:global-parity}
\end{equation}
where $\oplus$ denotes addition modulo two. Define two probability
distributions on $\mathcal A_S$:
\begin{equation}
\rho_b(a)=
\begin{cases}
\frac{1}{4}, & \Pi(a)=b,\\[4pt]
0, & \Pi(a)\neq b,
\end{cases}
\qquad b\in\{0,1\}.
\label{eq:parity-distributions}
\end{equation}
Thus $\rho_0$ is uniform on the even-parity states and $\rho_1$ is uniform on
the odd-parity states. The two distributions differ globally, but they have the
same marginal on every proper subset of $S$.

\begin{sproposition}[Proper-support indistinguishability]
\label{prop:proper-support-indistinguishability}
\Thm{} For every proper subset $U\subsetneq S$ and every assignment
$y\in\{0,1\}^{|U|}$,
\begin{equation}
(\pi_U)_*\rho_0(y)
=
(\pi_U)_*\rho_1(y)
=
2^{-|U|}.
\label{eq:proper-support-marginal}
\end{equation}
Consequently, every observable that factors through a proper-support restriction
has the same expectation value in $\rho_0$ and $\rho_1$.
\end{sproposition}

\begin{proof}
Fix a proper support $U$ and an assignment $y$ on that support. At least one
coordinate of $S$ remains unassigned. For either required global parity
$b\in\{0,1\}$, exactly
\begin{equation}
2^{|S|-|U|-1}
\end{equation}
completions of $y$ satisfy $\Pi(a)=b$. Each allowed completion has probability
$1/4$ under $\rho_b$. Hence
\begin{equation}
(\pi_U)_*\rho_b(y)
=
2^{|S|-|U|-1}\cdot\frac{1}{4}
=
2^{-|U|},
\end{equation}
independently of $b$. Any proper-support observable is a function of the
corresponding marginal and therefore has the same expectation under both
distributions.
\end{proof}

We now construct two finite typed candidate compositions. Let their P-space
records, local A(R)-space records, local operations, local bridge restrictions,
ordinary initial state, target declaration, and path prefix all be identical. Let
the declared full-support preparation terminate in $\rho_1$. The instances differ
only at the full-support composition boundary:
\begin{align}
\mathfrak C_{\mathrm{open}}:\quad&
\mathsf{Adm}(\gamma_S)=1,
\qquad
\mathcal B_S:\rho_1\longrightarrow p_\star
\quad\text{is type-valid and bridge-valid},
\label{eq:open-composition}
\\[4pt]
\mathfrak C_{\mathrm{blocked}}:\quad&
\mathsf{Adm}(\gamma_S)=0,
\qquad
\mathcal B_S
\quad\text{is blocked by a declared full-support constraint}.
\label{eq:blocked-composition}
\end{align}
The blocked instance does not deny that the same local preparation records can
be written. It denies that their full composition satisfies the global
admissibility condition needed for the bridge to the target. The distinction is
therefore not between two differently observed local systems. It is between two
globally typed compositions that induce the same complete proper-support
observation record.

\begin{sproposition}[Strict separation witness (standard parity pattern, instantiated here)]
\label{prop:strict-separation-witness}
\Imp{}/\Thm{} Proper-support parity separation is a familiar mathematical
antecedent, including in threshold-information constructions~\cite{S_Shamir1979}.
It is instantiated and re-proved here because the present witness attaches the
separation to typed bridge, regime, history, and support coordinates. The finite
candidate compositions
$\mathfrak C_{\mathrm{open}}$ and $\mathfrak C_{\mathrm{blocked}}$ satisfy
\begin{equation}
\mathfrak C_{\mathrm{open}}
\equiv_{\mathrm{loc}}
\mathfrak C_{\mathrm{blocked}},
\label{eq:strict-local-equivalence}
\end{equation}
while
\begin{equation}
\mathsf{Corr}(\mathfrak C_{\mathrm{open}})=1,
\qquad
\mathsf{Corr}(\mathfrak C_{\mathrm{blocked}})=0.
\label{eq:strict-corridor-separation}
\end{equation}
If the two instances share
\begin{equation}
\mathsf{Amp}=1,
\qquad
\mathsf{SafeWit}=0,
\end{equation}
then they also differ in conditional validation--hazard classification:
\begin{equation}
\mathsf{VHC}(\mathfrak C_{\mathrm{open}})=1,
\qquad
\mathsf{VHC}(\mathfrak C_{\mathrm{blocked}})=0.
\label{eq:strict-hazard-separation}
\end{equation}
\end{sproposition}

\begin{proof}
The physical and local registration records are identical by construction. The
equality of all proper-support registration marginals follows from
Proposition~\ref{prop:proper-support-indistinguishability}. Thus the complete
local-observation records coincide. The corridor predicate differs because the
full-support bridge is admissible in \eqref{eq:open-composition} and inadmissible
in \eqref{eq:blocked-composition}. The validation--hazard flag differs immediately
from its definition in \eqref{eq:validation-hazard-flag}.
\end{proof}

\begin{sproposition}[No local decider for the witness]
\label{prop:no-local-decider}
\Thm{} Let
\begin{equation}
D_{\mathrm{loc}}=F\circ\Lambda_{\mathrm{loc}}
\label{eq:local-decider}
\end{equation}
be any decision procedure whose input factors entirely through the declared
local-observation record. Then $D_{\mathrm{loc}}$ cannot be sound and complete
for $\mathsf{Corr}$ on a class containing both
$\mathfrak C_{\mathrm{open}}$ and $\mathfrak C_{\mathrm{blocked}}$.
\end{sproposition}

\begin{proof}
Equation~\eqref{eq:strict-local-equivalence} implies
\begin{equation}
D_{\mathrm{loc}}(\mathfrak C_{\mathrm{open}})
=
D_{\mathrm{loc}}(\mathfrak C_{\mathrm{blocked}}).
\end{equation}
Equation~\eqref{eq:strict-corridor-separation} requires a sound and complete
decider to return different values. The two requirements contradict one another.
\end{proof}

\begin{scorollary}[The strict witness is an unlicensed local contraction]
\label{cor:strict-witness-unlicensed-contraction}
\Thm{} Let $\LOp_{\mathrm{loc}}$ be the putative contraction that returns the
local-observation record $\Lambda_{\mathrm{loc}}$. On any class containing
$\mathfrak C_{\mathrm{open}}$ and $\mathfrak C_{\mathrm{blocked}}$ of
Proposition~\ref{prop:strict-separation-witness},
$\LOp_{\mathrm{loc}}$ is not licensed for the corridor target.
\end{scorollary}

\begin{proof}
The strict witness has
$\Lambda_{\mathrm{loc}}(\mathfrak C_{\mathrm{open}})=
\Lambda_{\mathrm{loc}}(\mathfrak C_{\mathrm{blocked}})$ while
$\mathsf{Corr}(\mathfrak C_{\mathrm{open}})\neq
\mathsf{Corr}(\mathfrak C_{\mathrm{blocked}})$. Thus corridor status is not
constant on a contraction fiber. Definition~\ref{def:licensed-contraction} and
Proposition~\ref{prop:local-quotient-lop} rule out a target-preservation
certificate for this contraction.
\end{proof}

\paragraph{What the witness proves and does not prove.}\label{rem:witness-scope}
The witness proves that no procedure restricted to the stated local quotient can
decide corridor existence. It does \emph{not} prove that no safe full-support
measurement, barrier certificate, topological invariant, or distributed
non-traversal calculation could ever decide it. The existence of such a safe
witness is exactly what $\mathsf{SafeWit}$ is intended to record. If one is
supplied and proven adequate, it narrows or removes the validation--hazard
coupling for that specific branch. It does not convert the composition predicate
into a local observable.

%-----------------------------------------------------------------
\subsection{Admissibility signatures and adequate reductions}
\label{supp:adequate-reductions}

The strict witness establishes that local observables are insufficient. The
stronger question is what a representation must retain in order to decide the
predicate without simply restating the full original record.

\begin{sdefinition}[Second-order admissibility signature]
\label{def:second-order-signature}
\Defn{} The \emph{second-order admissibility signature} of a typed composition
instance is
\begin{equation}
\Sigma_2(\mathfrak C)
=
\left(
\mathsf{Bridge}(\mathfrak C),
\mathsf{Regime}(\mathfrak C),
\mathsf{History}(\mathfrak C),
\mathsf{Support}(\mathfrak C)
\right),
\label{eq:second-order-signature}
\end{equation}
where
\begin{align}
\mathsf{Bridge}(\mathfrak C)
&=
\left(
\tau_{\Pspace/\ARspace},
\mathcal B,
\text{bridge-validity conditions}
\right),
\label{eq:signature-bridge}
\\
\mathsf{Regime}(\mathfrak C)
&=
\left(
\mathcal R_0,\ldots,\mathcal R_n,
\text{validity domains and exit conditions}
\right),
\label{eq:signature-regime}
\\
\mathsf{History}(\mathfrak C)
&=
\left(
u_1,\ldots,u_n,
\text{ordering},
\text{retention conditions},
\text{re-linearization basepoints}
\right),
\label{eq:signature-history}
\\
\mathsf{Support}(\mathfrak C)
&=
\left(
S,
\mathcal U_{\mathrm{loc}},
\text{incidence structure},
\text{full-support invariant}
\right).
\label{eq:signature-support}
\end{align}
Here $u_k$ denotes the step operator recorded with the temporal and
state-dependent information required for its composition; it is not merely an
unordered operation label. The signature is not asserted to be a unique data
structure or a minimal number of bits. It is a list of semantic distinctions that
the declared constraint class may require.
\end{sdefinition}

\begin{sdefinition}[Corridor-adequate reduction]
\label{def:corridor-adequate-reduction}
\Defn{} Let $\mathcal H_\star$ be a declared class of typed composition
instances. A reduction
\begin{equation}
q:\mathcal H_\star\longrightarrow Q
\label{eq:reduction-map}
\end{equation}
is \emph{corridor-adequate} when
\begin{equation}
q(\mathfrak C)=q(\mathfrak C')
\quad\Longrightarrow\quad
\mathsf{Corr}(\mathfrak C)=\mathsf{Corr}(\mathfrak C')
\qquad
\text{for all }\mathfrak C,\mathfrak C'\in\mathcal H_\star.
\label{eq:corridor-adequacy}
\end{equation}
Equivalently, $q$ is corridor-adequate if there exists a map
\begin{equation}
\overline{\mathsf{Corr}}:Q\longrightarrow\{0,1\}
\label{eq:reduced-corridor-map}
\end{equation}
such that
\begin{equation}
\mathsf{Corr}=\overline{\mathsf{Corr}}\circ q.
\label{eq:corridor-factorization}
\end{equation}
\end{sdefinition}

The definition separates two claims often conflated in informal discussion. A
representation may compress the original record substantially. But it may do so
only if its quotient classes never identify an open corridor with a blocked or
bridge-required corridor.

%-----------------------------------------------------------------
\subsection{The four-coordinate witness family}
\label{supp:four-coordinate-family}

To show why each component of $\Sigma_2$ is semantically necessary for this
constraint class, define a finite family of typed records
\begin{equation}
\mathfrak W
=
\left\{
\mathfrak C_{brhs}: b,r,h,s\in\{0,1\}
\right\},
\label{eq:four-coordinate-family}
\end{equation}
with the following interpretation.

\begin{enumerate}
\item $b=1$ means that the required bridge is legal: the P-space/A(R)-space
typing is preserved and the declared bridge is licensed. $b=0$ means that the
same untyped arrow exists only as an unlicensed substitution or an undeclared
bridge. Such a record is bridge-required rather than physically licensed.

\item $r=1$ means that every path increment lies inside the validity regime
of every imported relation used at that step. $r=0$ means that a step has left
the declared regime, even if its coarse numerical description appears continuous
with the preceding one.

\item $h=1$ means that the ordered path preserves the registration condition
needed by the next step: the required state is retained, the operations occur in
the required order, and re-linearization occurs at the actual current basepoint.
$h=0$ means that an otherwise identical unordered ledger of operations fails
because a prepared distinction relaxes, is overwritten, or is applied in an order
in which the required composition does not exist.

\item $s=1$ means that the full required support and invariant are attained.
$s=0$ means that only a proper-support proxy, local marginal, or incomplete
support pattern is attained. The parity construction above supplies a strict
finite realization of this distinction.
\end{enumerate}

For this witness family, define
\begin{equation}
\mathsf{Corr}(\mathfrak C_{brhs})=b\wedge r\wedge h\wedge s.
\label{eq:four-coordinate-corridor}
\end{equation}
Equation~\eqref{eq:four-coordinate-corridor} is not a claim that every physical
corridor has exactly four binary degrees of freedom. It is a finite typed witness
class in which each of the four semantic distinctions is independently decisive.
The physical worked case is more detailed: each coordinate expands into a
structured record of bridges, validity domains, history-dependent retention
conditions, and support topology.

The four decisive witness pairs are
\begin{align}
\mathfrak C_{1111}&\quad\text{versus}\quad\mathfrak C_{0111},
&&\text{bridge/type separation},
\label{eq:bridge-witness-pair}
\\
\mathfrak C_{1111}&\quad\text{versus}\quad\mathfrak C_{1011},
&&\text{regime separation},
\label{eq:regime-witness-pair}
\\
\mathfrak C_{1111}&\quad\text{versus}\quad\mathfrak C_{1101},
&&\text{history/retention separation},
\label{eq:history-witness-pair}
\\
\mathfrak C_{1111}&\quad\text{versus}\quad\mathfrak C_{1110},
&&\text{support/global-invariant separation}.
\label{eq:support-witness-pair}
\end{align}
Every pair differs in corridor status. A reduction that erases the relevant
coordinate identifies the pair and therefore cannot remain corridor-adequate.

%-----------------------------------------------------------------
\subsection{Second-order non-reducibility theorem}
\label{supp:nonreducibility}

\begin{sproposition}[Fiber criterion (standard quotient factorization)]
\label{prop:fiber-criterion}
\Imp{}/\Thm{} This is the standard quotient-factorization criterion.
It is re-theoremed in the present corridor setting because its proof fixes the
exact reduction fibers on which bridge, regime, history, and support distinctions
must be retained. A reduction $q:\mathcal H_\star\to Q$ admits a sound and
complete decider for $\mathsf{Corr}$ if and only if $\mathsf{Corr}$ is constant on
every fiber of $q$:
\begin{equation}
q^{-1}(z)=\left\{\mathfrak C\in\mathcal H_\star:q(\mathfrak C)=z\right\},
\qquad z\in Q.
\label{eq:reduction-fibers}
\end{equation}
\end{sproposition}

\begin{proof}
If a sound and complete decider
$\overline{\mathsf{Corr}}:Q\to\{0,1\}$ exists with
$\mathsf{Corr}=\overline{\mathsf{Corr}}\circ q$, then any two instances in the
same fiber receive the same corridor value. Conversely, if $\mathsf{Corr}$ is
constant on every fiber, define $\overline{\mathsf{Corr}}(z)$ as the common value
of $\mathsf{Corr}$ on $q^{-1}(z)$ whenever the fiber is non-empty, assigning an
arbitrary value to unused elements of $Q$. This yields the factorization in
Equation~\eqref{eq:corridor-factorization}.
\end{proof}

\begin{sproposition}[Second-order non-reducibility]
\label{prop:second-order-nonreducibility}
\Thm{} Let $\mathcal H_\star$ contain the four-coordinate witness family
\eqref{eq:four-coordinate-family}. Let
\begin{equation}
q:\mathcal H_\star\longrightarrow Q
\label{eq:putative-reduction}
\end{equation}
be any reduction through which a sound and complete corridor decision is
claimed to factor. Then $q$ must distinguish each of the four witness pairs
\eqref{eq:bridge-witness-pair}--\eqref{eq:support-witness-pair}. In particular,
\begin{align}
q(\mathfrak C_{1111})&\neq q(\mathfrak C_{0111}),
\label{eq:must-retain-bridge}
\\
q(\mathfrak C_{1111})&\neq q(\mathfrak C_{1011}),
\label{eq:must-retain-regime}
\\
q(\mathfrak C_{1111})&\neq q(\mathfrak C_{1101}),
\label{eq:must-retain-history}
\\
q(\mathfrak C_{1111})&\neq q(\mathfrak C_{1110}).
\label{eq:must-retain-support}
\end{align}
Therefore no sound and complete corridor decider can erase, respectively:
\begin{enumerate}
\item the distinction between a legal typed P-space/A(R)-space bridge and an
unlicensed substitution;

\item the distinction between remaining in the validity regime and leaving it;

\item the distinction between an ordered, retained, re-linearized path and an
unordered or relaxation-erased operation ledger; or

\item the distinction between attainment of the required full support/global
invariant and attainment of only locally indistinguishable proper-support records.
\end{enumerate}
Any representation which decides the predicate may encode these distinctions in
different symbols, data structures, equations, proof objects, or simulation
states. But it must preserve an extensionally equivalent distinction: it must not
identify any pair whose corridor values differ.
\end{sproposition}

\begin{proof}
By Equation~\eqref{eq:four-coordinate-corridor},
\begin{align}
\mathsf{Corr}(\mathfrak C_{1111})&=1,\\
\mathsf{Corr}(\mathfrak C_{0111})
&=\mathsf{Corr}(\mathfrak C_{1011})
=\mathsf{Corr}(\mathfrak C_{1101})
=\mathsf{Corr}(\mathfrak C_{1110})=0.
\end{align}
Suppose, for contradiction, that $q$ identifies $\mathfrak C_{1111}$ with any one
of the four records on the right. Then $\mathsf{Corr}$ is not constant on the
corresponding fiber of $q$. Proposition~\ref{prop:fiber-criterion} implies that no
sound and complete corridor decider can factor through $q$. The contradiction
applies independently to each witness pair.
\end{proof}

\begin{sproposition}[Minimality of the second-order layer]
\label{prop:second-order-minimality}
\Thm{} For the constraint class represented by $\mathcal H_\star$, a state-only,
local-observable-only, unordered-operation, or support-erasing representation is
not corridor-adequate. A decision procedure must instead retain a sufficient
statistic at least as discriminating as the second-order admissibility signature
$\Sigma_2$ on every output-relevant pair.
\end{sproposition}

\begin{proof}
A state-only or local-observable-only representation identifies the strict
separation witness of Proposition~\ref{prop:strict-separation-witness}. A
representation erasing bridge legality, regime state, history, or support
identifies at least one witness pair in
Equations~\eqref{eq:bridge-witness-pair}--\eqref{eq:support-witness-pair}. Each
identification violates the fiber criterion.
\end{proof}

\paragraph{The precise meaning of ``minimal''.}\label{rem:minimality-scope}
Proposition~\ref{prop:second-order-minimality} is a semantic minimality result,
not a claim that the present P-space/A(R)-space notation is uniquely privileged,
that the proposed simulation is computationally minimal, or that no shorter
mathematical derivation exists. A different framework defeats the result only by
supplying a reduction that remains corridor-adequate: it must prove that its
quotient never identifies an open corridor with a blocked, bridge-required,
regime-invalid, history-invalid, or support-incomplete one. Such a framework is
extensionally second-order with respect to this predicate, whether or not it uses
the terms P-space, A(R)-space, bridge, support, or global coherence.

The result is therefore not a claim of notation ownership. It is a claim about
the irreducible distinctions that any valid assessment must preserve. Recognition
that its fiber criterion has a familiar quotient-theoretic antecedent does not
negate this result: it accounts for one payload relation only. By
Section~\ref{supp:novelty}, a complete novelty verdict must also preserve the
typed registrations generated by its deployment or retain their difference as
budgeted novelty residue.

%-----------------------------------------------------------------
\subsection{Consequences for the main necessity claim}
\label{supp:consequences}

The two results sharpen the main article's argument. First,
Proposition~\ref{prop:strict-separation-witness} establishes a strict separation
between local derive-and-test and a composition-level corridor predicate. The
limitation is not merely that contemporary instruments may lack sensitivity, nor
merely that a proposed signal might be weak. The predicate can differ while the
entire proper-support local-observation record is unchanged.

Second, Proposition~\ref{prop:second-order-nonreducibility} specifies what the
second-order layer must preserve. The earlier obligation list is thereby no longer
only a prudent engineering checklist. For the declared witness class, the core
distinctions are necessary conditions of sound and complete assessment.

Third, the theorem identifies the exact burden for an attempted simplification. A
critic need not accept the present representation. But a valid replacement must
provide one of the following:
\begin{enumerate}
\item a corridor-adequate quotient that proves the omitted distinction cannot
alter the predicate in the declared class;

\item a barrier theorem showing that the full-support route is blocked before
the disputed distinction becomes relevant;

\item a safe full-support witness or certificate that decides the global
predicate without traversal, thereby changing the $\mathsf{SafeWit}$ branch in the assessment record; or

\item a typed counterexample showing that one of the witness pairs is ill-formed
under the declared semantics.
\end{enumerate}
Absent one of these, a reduction that discards bridge type, regime state, history,
or support structure has not decided the constraint. It has only changed the
representation until the deciding distinction is no longer visible. Questions not
settled by one of these finite routes remain assessment residue in the sense of
Definition~\ref{def:residue-defeater}; they do not become typed defects merely
because the assessment budget is finite.

%-----------------------------------------------------------------
\subsection{Scope}
\label{supp:strict-separation-scope}

Nothing in this section establishes that the worked-case corridor or its
target-conditioned schedule exists, that a physical registration variable can be
prepared as hypothesized, that entanglement--geometry relations transfer to the
intended target, that self-amplification occurs, that a force-sign reversal is
available, that time reversal occurs, or that a macroscopic geometric effect is
reachable.
Those remain conditional, bridge-required, or open exactly as in the rest of the
paper.

The contribution here is more limited and more formal. Given a declared class of
candidate compositions in which bridge legality, regime membership, ordered
retention, and full-support invariants can independently alter the corridor
predicate, no first-order local quotient can be a sound and complete assessment
representation. Any successful replacement must preserve those distinctions, or
prove that they are irrelevant for the class it claims to decide. A pass under a
finite coherence budget remains a status-labeled assessment result; unexamined
realization and empirical branches are residue unless a finite typed defeater is
supplied.



%=================================================================
\section{The matter-observation blind spot and hard problems}
\label{supp:matter-blindspot}

The finite corridor witness shows that a declared local quotient can erase a
physical decision distinction. This section specializes that result to
matter-sector observation and makes the resource question explicit. It begins with
the coupled assessment record and then tests the matter-sector quotient as a
candidate contraction. The relevant resource is not energy in isolation: it is the
availability of a fiber-separating carrier, interaction, or bridge.

\subsection{Matter-sector quotient and energy closure}

\begin{sdefinition}[Matter-sector assessment domain and protocol family]
\label{def:matter-domain}
\Defn{} Fix a physical target $Q$ and let $\mathcal C_Q$ be the coupled
P-space/A(R)-space domain of physical states, registration states, and typed bridge
records relevant to $Q$. For each available energy scale $E$, let
$\mathfrak M_E$ be the class of admissible matter-sector preparation, interaction,
detection, and registration protocols executable at or below $E$. A protocol
$\pi\in\mathfrak M_E$ has a physical-registration update
$\Phi_{\pi,E}:\mathcal C_Q\to\mathcal C_Q$ and an output registration map
$H_{\pi,E}:\mathcal C_Q\to\mathcal R_{\pi,E}$.

A \emph{matter-observation quotient} is a candidate first-order contraction of
this coupled assessment domain, written
\begin{equation}
q_{\mathrm M}:\mathcal C_Q\longrightarrow\overline{\mathcal C}_{\mathrm M}
\label{eq:matter-quotient-supp}
\end{equation}
whose fibers identify exactly the coupled states that the active matter-sector
assessment treats as observationally equivalent. The quotient is assessed from
$\mathcal C_Q$; it is not the starting object from which the coupled record is
reconstructed.
\end{sdefinition}

\begin{sdefinition}[Matter-sector energy closure]
\label{def:matter-energy-closure}
\Defn{} The matter-sector protocol family is \emph{energy closed} at
$q_{\mathrm M}$ when, for every available energy scale $E$ and every
$\pi\in\mathfrak M_E$, there are maps $\overline H_{\pi,E}$ and
$\overline\Phi_{\pi,E}$ such that
\begin{equation}
H_{\pi,E}=\overline H_{\pi,E}\circ q_{\mathrm M},
\qquad
q_{\mathrm M}\circ\Phi_{\pi,E}
=\overline\Phi_{\pi,E}\circ q_{\mathrm M}.
\label{eq:matter-energy-closure-supp}
\end{equation}
Energy closure means that the full available matter-sector program refines and
extends registrations inside one invariant quotient. It can generate new
matter-sector states, records, and measurements without separating states already
merged by $q_{\mathrm M}$.
\end{sdefinition}

\begin{sdefinition}[Matter-observability blind spot]
\label{def:matter-blindspot}
\Defn{} Let $\Theta_Q:\mathcal C_Q\to\mathcal Y_Q$ be the declared physical
predicate for target $Q$. An \emph{energy-independent matter-observability blind
spot} exists for $Q$ when the matter-sector protocol family is energy closed and
there exist $c_0,c_1\in\mathcal C_Q$ satisfying
\begin{equation}
q_{\mathrm M}(c_0)=q_{\mathrm M}(c_1),
\qquad
\Theta_Q(c_0)\neq\Theta_Q(c_1).
\label{eq:matter-target-nonfactor-supp}
\end{equation}
The quotient therefore erases a relation on which the physical target changes.
\end{sdefinition}

\begin{stheorem}[Energy-independent matter-observability blind spot]
\label{thm:matter-blindspot}
\Thm{} Under Definitions~\ref{def:matter-energy-closure} and
\ref{def:matter-blindspot}, no protocol in $\bigcup_E\mathfrak M_E$ decides
$\Theta_Q$. The blind spot is removed exactly by a change in observational
architecture that supplies a carrier, interaction, or bridge with a registration
map that separates a pair in a $q_{\mathrm M}$-fiber on which $\Theta_Q$ differs.
\end{stheorem}

\begin{proof}
Let $\pi\in\mathfrak M_E$. By Equation~\eqref{eq:matter-energy-closure-supp},
$H_{\pi,E}$ factors through $q_{\mathrm M}$. Every decision rule formed from this
output consequently assigns the same value to $c_0$ and $c_1$. By
Equation~\eqref{eq:matter-target-nonfactor-supp}, $\Theta_Q$ assigns them
different values. Thus no such rule decides $\Theta_Q$. The argument is uniform in
$E$, so enlarging energy leaves the conclusion unchanged while energy closure
holds. A new registration map that separates the pair changes the quotient and
therefore changes the available decision information.
\end{proof}

\begin{sproposition}[Blind-spot obligation for a physics program]
\label{prop:blindspot-obligation}
\Thm{} A physics program that uses a matter-observation quotient as exhaustive for
$Q$ requires a factorization certificate
\begin{equation}
\Theta_Q=\overline\Theta_Q\circ q_{\mathrm M}.
\label{eq:matter-factorization-certificate-supp}
\end{equation}
In the absence of that certificate, the program carries a matter-observability
blind-spot obligation: its matter-sector conclusions classify quotient-level
registrations while the target-relevant fiber distinction remains active in the
assessment record.
\end{sproposition}


\subsection{Synchronization residue as a topological fiber coordinate}
\label{supp:synchronization-residue}

The quotient criterion becomes physically inspectable when the omitted relation is
registered as a typed residue rather than named only as an abstract difference.

\begin{sdefinition}[Material observation composite]
\label{def:material-observation-composite}
\Defn{} Let $M$ be a material observer at time $t$. Let
$r_{M,t,\delta}$ restrict a P-space path to the subnetwork physically accessible
to $M$, let $\beta_{\mathcal E,\delta}$ be a typed P-space/A(R)-space bridge, and
let $q_{M,t,\delta}$ be the material registration quotient. The material
observation composite is
\begin{equation}
\ObsMap^{\mathcal E}_{M,t,\delta}=
q_{M,t,\delta}\circ\beta_{\mathcal E,\delta}\circ r_{M,t,\delta}.
\label{eq:material-observation-composite-supp}
\end{equation}
It maps a physical state or path to the equivalence class retained by the
observer's physical access topology, registration topology, bridge conditions, and
material quotient.
\end{sdefinition}

\begin{sdefinition}[Synchronization residue]
\label{def:synchronization-residue}
\Defn{} Let $a_M$ be an observation node generated by a material observer's
accessible recursively composed observation topology, and let $a_Q$ be a
registration node or observation process required to decide target $Q$. Each may
itself be a recursively composed object. The \emph{synchronization residue} is the
target-relative topological relation retained by $a_Q$ but not retained by $a_M$.
It is zero when a typed target-preserving admissible deformation identifies the two
observation topologies. It is nonzero when a support, scheduling, boundary, path,
or bridge relation remains distinct for $Q$.

For a path-level realization, choose a physical path $\pi^-_M(p;t)$ whose
registration is retained by $M$ and an admissible target-support path
$\pi^+_Q(p;t)$ required by $Q$. The residue is represented by
\begin{equation}
\SyncRes^{M,Q}(p;t,\delta)=
\left[
\beta_{\mathcal E,\delta}(\pi^-_M(p;t)),
\beta_{\mathcal E,\delta}(\pi^+_Q(p;t))
\right]_{\mathrm{rel},Q}.
\label{eq:synchronization-residue-supp}
\end{equation}
Ordered regimes are those in which a declared synchronizing composition reduces
the residue; disordered regimes are those in which the active composition
preserves, branches, or enlarges it.
\end{sdefinition}

\begin{sproposition}[Residue witness for a material blind spot]
\label{prop:residue-blindspot-witness}
\Thm{} Suppose $p_0,p_1$ satisfy
\begin{equation}
\ObsMap^{\mathcal E}_{M,t,\delta}(p_0)=
\ObsMap^{\mathcal E}_{M,t,\delta}(p_1),
\qquad
\SyncRes^{M,Q}(p_0;t,\delta)\neq
\SyncRes^{M,Q}(p_1;t,\delta),
\label{eq:residue-witness-equality}
\end{equation}
and the target varies with that residue:
\begin{equation}
\Theta_Q(p_0)\neq\Theta_Q(p_1).
\label{eq:residue-witness-target}
\end{equation}
Then the material observation composite has a target-distinct fiber and therefore
witnesses a matter-observability blind spot for $Q$.
\end{sproposition}

\begin{proof}
Equation~\eqref{eq:residue-witness-equality} places $p_0$ and $p_1$ in the same
fiber of the material observation composite. Equation~\eqref{eq:residue-witness-target}
shows that $\Theta_Q$ is not constant on that fiber. Thus no quotient-level
decision through the composite can determine $\Theta_Q$. The different residue
values identify a typed topological relation erased by the material quotient.
\end{proof}

\subsection{Productive dual branches and branch-selected conclusion attractors}

\begin{sdefinition}[Separation and quotient-closure branches]
\label{def:matter-dual-branches}
\Defn{} Fix an initial registration $R_0$ for target $Q$. A \emph{separation
branch} is a constraint-transport program $F_{\mathrm{sep}}$ that generates,
checks, or transports registrations bearing on a target-distinct pair within a
matter-observation fiber: typed witnesses, bridge requirements, contradiction
records, support certificates, or fiber-separating candidate measurements. A
\emph{quotient-closure branch} is a constraint-transport program
$F_{\mathrm{cl}}$ that generates, checks, or transports matter-sector
registrations and closure attempts supporting the conclusion label that the active
matter quotient is sufficient for $Q$.

The corresponding generated branch closures are
\begin{equation}
\mathcal B^{\infty}_{\mathrm{sep}}(R_0)=
\bigcup_{n\geq0}F_{\mathrm{sep}}^n(R_0),
\qquad
\mathcal B^{\infty}_{\mathrm{cl}}(R_0)=
\bigcup_{n\geq0}F_{\mathrm{cl}}^n(R_0).
\label{eq:matter-dual-branches}
\end{equation}
Each branch may contain theorems, proofs, empirical measurements, simulations,
registration records, comparison procedures, and other physical assessment
objects. Their common grammar is $\langle\Register,\Recall,\GOp,\LOp\rangle_{\rm rec}$.
\end{sdefinition}

\begin{sdefinition}[Productive unboundedness]
\label{def:productive-unboundedness}
\Defn{} A branch program $F$ is \emph{productive} when it carries a declared
provenance, support, or registration-depth measure
$\nu:\mathcal R\to\mathbb N$ such that
\begin{equation}
\nu(F^{n+1}(R_0))>\nu(F^n(R_0))
\qquad\text{for every }n\geq0.
\label{eq:productive-branch}
\end{equation}
A productive branch is unbounded in its generating, checking, and interaction
function: for every finite depth $N$, it contains a generated assessment object
with $\nu>N$. The condition distinguishes genuinely expanding physics-object
families from a finite cycle repeatedly presented as recursion.
\end{sdefinition}

\begin{sdefinition}[Branch-selected conclusion attractor]
\label{def:branch-attractor}
\Defn{} Let $\operatorname{cl}_{\rm adm}$ denote closure under the admitted
constraint-transport programs for a declared assessment. When a productive branch
preserves a conclusion label $\sigma_\star$ throughout its closure, its
\emph{branch-selected conclusion attractor} is
\begin{equation}
\mathsf{Attr}_{\star}(R_0)=
\operatorname{cl}_{\rm adm}\bigl(\mathcal B^{\infty}_{\star}(R_0)\bigr),
\qquad \star\in\{\mathrm{sep},\mathrm{cl}\}.
\label{eq:branch-attractor}
\end{equation}
The attractor is a stable conclusion object in registration space: its associated
label is preserved by the branch updates that define the closure. Its physical
interpretation is licensed only by the bridges retained in the same record.
\end{sdefinition}

\begin{stheorem}[Unbounded dual branches yield branch-selected attractors]
\label{thm:branch-attractors}
\Thm{} Suppose $F_{\mathrm{sep}}$ and $F_{\mathrm{cl}}$ are productive, preserve
distinct conclusion labels, and no active bridge-valid program resolves the
fiber distinction in Equation~\eqref{eq:matter-target-nonfactor-supp}. Then
$\mathsf{Attr}_{\mathrm{sep}}(R_0)$ and
$\mathsf{Attr}_{\mathrm{cl}}(R_0)$ are distinct unbounded conclusion attractors.
An assessment initialized in the basin of either branch reaches the conclusion
label preserved by that branch. A fiber-separating $\GOp$ expansion or a certified
factorization $\LOp$ contraction changes the branch dynamics by resolving the
active quotient relation.
\end{stheorem}

\begin{proof}
Productivity makes each branch closure unbounded by
Definition~\ref{def:productive-unboundedness}. Label preservation makes the label
stable under each admitted branch update and therefore under
$\operatorname{cl}_{\rm adm}$. The labels are distinct by assumption, so their
closures are distinct conclusion attractors. In the absence of a bridge-valid
separating program or a factorization certificate, no admitted update identifies
the two target values within the active quotient; the branch-selected label remains
stable. The indicated $\GOp$ or $\LOp$ operation supplies exactly the missing
resolution condition and thereby changes the available transition structure.
\end{proof}

The result classifies an assessment dynamic rather than a vote between physical
truths. The physical target has its declared value in $\Theta_Q$; the attractor
records the conclusion reached by an observation architecture whose generated
branch remains quotient-limited. The branch distinction therefore identifies a
structural source of persistent hard-problem behavior: unbounded evidence
production can reinforce incompatible conclusion labels while the
fiber-separating relation remains unregistered.

\subsection{Physical hard-problem metacategory}

\begin{sdefinition}[Physical hard-problem metacategory]
\label{def:physical-hard-metacategory}
\Defn{} Fix a coherence budget $\mathcal B$. The \emph{physical hard-problem
metacategory}
\begin{equation}
\mathsf{Hard}^{\mathrm{phys}}_{\GOp/\LOp,\mathcal B}
\label{eq:physical-hard-metacategory}
\end{equation}
contains an object for each target $Q$ with a record
\begin{equation}
\mathbf H_Q=
\left(
Q,\mathcal C_Q,q_{\mathrm M},\Theta_Q,
\mathcal B^{\infty}_{\mathrm{sep}},
\mathcal B^{\infty}_{\mathrm{cl}},
\mathsf{Attr}_{\mathrm{sep}},
\mathsf{Attr}_{\mathrm{cl}},
\mathcal B,\mathsf{Status}
\right)
\label{eq:hard-problem-object}
\end{equation}
when all of the following hold:
\begin{enumerate}[label=(\roman*)]
\item $q_{\mathrm M}$ identifies states on which $\Theta_Q$ differs;
\item the available matter-sector protocol family is energy closed at
      $q_{\mathrm M}$;
\item the separation and quotient-closure branches are productive;
\item the declared budget contains no bridge-valid fiber-separating program and
      no factorization certificate for $\Theta_Q$ through $q_{\mathrm M}$.
\end{enumerate}

Its arrows are admitted non-erasing constraint-transport programs that map one
such record to another while carrying the quotient, target predicate, branch
status, and bridge conditions needed by the destination audit. Arrow composition
is typed program composition when these signatures compose. The structure is a
metacategory because its objects are physics problems together with the
observation architectures that generate, assess, and transform their problem
records.
\end{sdefinition}

\begin{sproposition}[Hardness is an observability architecture]
\label{prop:hardness-observability}
\Thm{} Membership in
$\mathsf{Hard}^{\mathrm{phys}}_{\GOp/\LOp,\mathcal B}$ classifies a target by the
architecture of its observability rather than by its subject-matter label. A target
leaves the metacategory when a registered bridge, carrier, interaction, or
factorization certificate resolves the target-distinct matter-observation fiber.
\end{sproposition}

%=================================================================
\section{Supplementary audit infrastructure}
\label{supp:architecture}

This section gives optional audit infrastructure. The finite strict-separation
witness, fiber criterion, and non-reducibility theorem do not use its
generating-object, rendering, or human-verifier assumptions. When a selected
finite audit is required, the constructions below provide possible typed
implementations of $\GOp$-expansion and licensed $\LOp$-contraction: they retain
support, boundary, scale, bridge, and registration relations when a theory object
is restricted, refined, or compiled into a finite certificate. The resulting maps do not merely display a physical claim; they
specify which selected record continues to represent the same constraint. This is
consistent with the general need to state physical relations across scale explicitly
\cite{S_Wilson1971} and, in holographic settings, with support-sensitive
reconstruction through boundary regions and code subspaces
\cite{S_AlmheiriDongHarlow2015,S_DongHarlowWall2016}.

The audit component performs the next transport step. A physical question may
depend on several simultaneously active state, operation, boundary, regime,
support, history, and registration distinctions. The architecture compiles the
minimal exact set of active relations into a finite selected graph, a before/after
change record, a legend, and local failure checks. The resulting certificate is the
finite verification interface through which a bounded assessor can check, falsify,
transfer, or refuse a gravity claim. Its role is structurally analogous to a
proof-carrying certificate: the receiver checks a bounded object rather than
reconstructing the full production history \cite{S_Necula1997}. Provenance practice
likewise treats entities, activities, agents, and derivational relations as explicit
records rather than as an untyped residual trace \cite{S_MoreauEtAl2013}. The
result below formalizes exact semantic compression and a finite-budget certificate
condition. It does not require any empirical diagnosis of research activity.

\subsection{Typed generator and all-scale representation}

\begin{shypothesis}[Typed physical generating object]
\label{arch:hyp-generator}
\Hyp{} There exists a typed generating object
\begin{equation}
\Gobj=\left\langle
\mathsf{PCarrier},\Sigma,\mathcal O_P,\mathcal B,\mathcal I,\Adm,\mathsf{Comp}
\right\rangle,
\label{arch:eq-generator}
\end{equation}
where $\mathsf{PCarrier}$ is the carrier of physically admissible bounded-region
states; $\Sigma$ is a typed signature containing carrier, boundary, interface,
operator, regime, support, and scale roles; $\mathcal O_P$ is a family of primitive
physical operations; $\mathcal B$ and $\mathcal I$ assign boundary and interface
data; $\Adm$ determines admissible states, transitions, and finite compositions;
and $\mathsf{Comp}$ is partial composition for admissible transitions.
\end{shypothesis}

\begin{sproposition}[No scalar energy functional is presupposed]
\label{arch:prop-not-scalar}
\Defn{} The generating object of Hypothesis~\ref{arch:hyp-generator} is not assumed
to have scalar codomain or to be a Hamiltonian, Lagrangian, or variational
functional. Any such realization may be introduced only with its own field domain,
units, recovery map, and regime. Calling the object an energy functional before
those data are supplied would add structure not used below.
\end{sproposition}

\begin{saxiom}[Closure of admissible physical composition]
\label{arch:ax-closure}
\Cnd{} If $\gamma_1$ and $\gamma_2$ are admissible physical transitions and their
target/source, boundary, interface, and regime types compose, then
$\gamma_2\circ\gamma_1$ is defined and admissible. Every finite admissible path
therefore has a well-typed composite.
\end{saxiom}

\begin{sdefinition}[Bounded region and finite resolution]
\label{arch:def-region-resolution}
\Defn{} A \emph{bounded region} $R$ is a declared portion of the physical carrier
together with boundary and interface data. A \emph{finite resolution} $\delta$ is a
declared distinguishability threshold and finite active type vocabulary at which
two physical descriptions are identified only when their difference is below the
selected audit-relevant threshold. Neither $R$ nor $\delta$ is fixed across all
physical questions.
\end{sdefinition}

\begin{saxiom}[Finitary bounded-region serialization]
\label{arch:ax-finitary}
\Cnd{} For every bounded region $R$ and declared finite resolution $\delta$, the
admissible physical state and all transition distinctions retained at $(R,\delta)$
admit a finite typed serialization. Equivalently, a finite quotient of the
relevant state-transition structure has vertices, edges, labels, and metric data
that determine every distinction declared retained at that scale.
\end{saxiom}

\begin{saxiom}[Scale compatibility]
\label{arch:ax-scale}
\Cnd{} If $\delta'\preceq\delta$ denotes refinement to a finer resolution, then
there is a typed refinement map
\begin{equation}
\rho_{\delta'\to\delta}:\Pspace_{\mathcal E;R,\delta'}\longrightarrow
\Pspace_{\mathcal E;R,\delta}
\label{arch:eq-refinement}
\end{equation}
that preserves all distinctions retained at the coarser scale. If $R'\subseteq R$,
there is a typed restriction map
\begin{equation}
\operatorname{res}_{R\to R'}:\Pspace_{\mathcal E;R,\delta}\longrightarrow
\Pspace_{\mathcal E;R',\delta}
\label{arch:eq-restriction}
\end{equation}
that preserves the declared boundary/interface relation.
\end{saxiom}

\begin{sdefinition}[All-scale P-space schema]
\label{arch:def-allscale}
\Defn{} The \emph{all-scale P-space schema} of $\Gobj$ is the compatible family
\begin{equation}
\mathbf P(\Gobj)=
\left\{\Pspace_{\mathcal E;R,\delta};\rho_{\delta'\to\delta},
\operatorname{res}_{R\to R'}\right\}_{R,\delta}.
\label{arch:eq-allscale}
\end{equation}
It is not one finite graph. It is a refinement- and restriction-compatible family
of finite graphs, one for every bounded region and finite resolution in the
declared domain.
\end{sdefinition}

\begin{stheorem}[Universal finite-resolution representation]
\label{arch:thm-representation}
\Cnd{} Under Hypothesis~\ref{arch:hyp-generator} and
Axioms~\ref{arch:ax-closure}--\ref{arch:ax-scale}, every $\Gobj$-admissible
functional state of every bounded region $R$, and every admissible finite
transition path between such states, is represented by the all-scale P-space
schema at every resolution at which the relevant distinctions are retained.
\end{stheorem}

\begin{proof}
Fix $R$, $\delta$, and an admissible physical state $x$. By
Axiom~\ref{arch:ax-finitary}, the retained distinctions of $x$ at $(R,\delta)$
have a finite typed serialization. By Definition~\ref{def:pspace}, the resulting
state class is a vertex of $\Pspace_{\mathcal E;R,\delta}$. Now fix an admissible
transition $\gamma:x\to y$. Its source, target, operator, boundary, interface, and
regime data are admissible by Hypothesis~\ref{arch:hyp-generator}; hence
Definition~\ref{def:pspace} supplies a typed edge whenever those distinctions are
retained at $\delta$. A finite admissible path has a well-typed composite by
Axiom~\ref{arch:ax-closure}. Compatibility across resolutions and subregions is
supplied by Axiom~\ref{arch:ax-scale}. Therefore the state and every retained
finite transition path are represented at every relevant finite resolution.
\end{proof}

\begin{sproposition}[Universality is schema-level, not finite enumeration]
\label{arch:prop-universality}
\Defn{} Theorem~\ref{arch:thm-representation} does not claim that one finite
picture enumerates every possible microstate of an open-ended universe. It states
that every bounded region at every chosen finite resolution has a finite P-space
representation, with compatible refinement and restriction maps. The universal
object is the schema in Equation~\eqref{arch:eq-allscale}; the human-auditable
object is always a finite selected member of that schema.
\end{sproposition}

\subsection{Audit-minimal compression}

\begin{sdefinition}[Finite audit family]
\label{arch:def-audit-family}
\Defn{} Fix $(R,\delta)$ and let $X_{R,\delta}$ be the set of finite P-space/SOP
records in the declared domain. A finite \emph{audit family} is
\begin{equation}
\Audit=(a_1,\ldots,a_m),\qquad a_j:X_{R,\delta}\to Y_j,
\label{arch:eq-audit-family}
\end{equation}
where each $a_j$ is a finite question, for example whether an operation is
type-admissible, a bridge legal, a regime retained, a selected interaction changes
a declared invariant, or a before/after path closes under composition.
\end{sdefinition}

\begin{sdefinition}[Audit equivalence, quotient, and adequacy]
\label{arch:def-audit-quotient}
\Defn{} Two records $x,y\in X_{R,\delta}$ are audit-equivalent for $\Audit$ when
\begin{equation}
x\sim_{\Audit}y\quad\Longleftrightarrow\quad
\forall j\in\{1,\ldots,m\},\ a_j(x)=a_j(y).
\label{arch:eq-audit-equivalence}
\end{equation}
The audit quotient is $Q_{\Audit}(X_{R,\delta})=X_{R,\delta}/\sim_{\Audit}$. A
representation $r:X_{R,\delta}\to Z$ is sound and complete for $\Audit$ if every
$a_j$ factors through $r$, i.e. $a_j=\bar a_j\circ r$ for some
$\bar a_j:Z\to Y_j$.
\end{sdefinition}

\begin{stheorem}[Audit quotient is maximally compressed exact representation]
\label{arch:thm-compression}
\Thm{} The quotient map $q_{\Audit}:X_{R,\delta}\to Q_{\Audit}(X_{R,\delta})$ is
sound and complete for $\Audit$. If $r:X_{R,\delta}\to Z$ is any sound-and-complete
representation for $\Audit$, then
\begin{equation}
\ker(r)\subseteq\ker(q_{\Audit}).
\label{arch:eq-kernel-inclusion}
\end{equation}
Thus $r$ may retain more distinctions than the audit quotient, but may not identify
any pair that the audit quotient distinguishes. The quotient is the coarsest, hence
maximally compressed, exact semantic representation for the declared audit family.
\end{stheorem}

\begin{proof}
For each $j$, define $\widetilde a_j([x]_{\sim_{\Audit}})=a_j(x)$. It is
well-defined because equivalent records agree on every audit. Hence
$a_j=\widetilde a_j\circ q_{\Audit}$. Now let $r$ be sound and complete. If
$r(x)=r(y)$, then for every $j$,
\[
a_j(x)=\bar a_j(r(x))=\bar a_j(r(y))=a_j(y).
\]
Therefore $x\sim_{\Audit}y$ and $q_{\Audit}(x)=q_{\Audit}(y)$, proving
Equation~\eqref{arch:eq-kernel-inclusion}.
\end{proof}

\begin{sproposition}[Meaning of maximal compression]
\label{arch:prop-maximal-compression}
\Defn{} Theorem~\ref{arch:thm-compression} proves maximal compression in a
semantic, task-relative sense. It does not claim a universally shortest bit string,
a unique notation, a smallest computational cost, or a uniquely privileged drawing.
It proves only that no exact representation for the declared audit family may
discard a distinction on which an audit answer changes.
\end{sproposition}


\subsection{Finite-budget auditability and status labeling}

\begin{sdefinition}[Finite audit budget]
\label{arch:def-audit-budget}
\Defn{} A finite audit budget is a tuple
\begin{equation}
 b=(n_V,n_E,n_T,n_C,n_L),
\label{arch:eq-audit-budget}
\end{equation}
where $n_V,n_E,n_T,n_C,n_L$ respectively bound the active vertices, active edges or
path segments, active type tokens, explicit local checks, and legend/instruction
items that a declared verifier protocol may require for one audit. The tuple is a
protocol parameter. It is not a claim that all humans share one fixed capacity.
\end{sdefinition}

\begin{sdefinition}[Budgeted audit certificate]
\label{arch:def-budgeted-certificate}
\Defn{} Let $\Audit$ be a finite audit family on $X_{R,\delta}$ and let $b$ be a
finite audit budget. A \emph{budgeted audit certificate} for $\Audit$ is a selected
finite subgraph, an exact factorization of $\Audit$ through that selected record,
and a finite legend/change record whose active counts are componentwise bounded by
$b$. It is the executable decision interface between a full typed constraint record
and a finite verifier: it states which roles are active, which relations changed,
which local check validates each required relation, and which absent or mismatched
relation invalidates the conclusion. A claim is \emph{human-audited at budget $b$}
only when such a certificate is supplied and its declared checks are executed by a
verifier satisfying the relevant protocol premise.
\end{sdefinition}

\begin{sproposition}[Why finite selected certificates are required]
\label{arch:prop-budgeted-need}
\Defn{} Consider a declared audit family $\Audit$ and a finite verifier budget $b$.
A representation that either (i) identifies two records on which an audit answer
differs, or (ii) provides no budgeted audit certificate, does not establish a
reliable human audit of that question at $b$. In case (i), exactness fails by
Theorem~\ref{arch:thm-fiber}; in case (ii), no finite executable verification path
has been supplied. The required conclusion is not that the physical question is
false, but that its human-audit status remains open at the declared budget.
\end{sproposition}

\begin{sproposition}[Audit-minimality is the exact simplicity condition]
\label{arch:prop-exact-simplicity}
\Defn{} For a declared audit family, ``maximally simple'' means the coarsest exact
semantic representation, not a universal claim about aesthetic preference,
psychology, or a uniquely privileged diagram. By
Theorem~\ref{arch:thm-compression}, the audit quotient removes every distinction
that leaves all declared audit answers unchanged and forbids removal of any
distinction on which an answer changes. Thus it is the unique exact simplicity
condition up to isomorphism of the quotient and reversible re-encoding.
\end{sproposition}


\subsection{Universal finite-graph rendering in three dimensions}

\begin{sdefinition}[Topology-faithful typed rendering]
\label{arch:def-rendering}
\Defn{} Let $G=(V,E,\tau)$ be a finite typed multigraph. A topology-faithful typed
rendering in $\mathbb R^d$ is an embedding of its underlying one-dimensional
CW-complex into $\mathbb R^d$ together with a finite legend assigning every active
vertex and edge type to an unambiguous token: label, texture, line style,
orientation marker, port marker, or finite displayed value. The rendering need not
use Euclidean distance to represent every P-space metric datum; metric data may be
carried by labels, weights, slices, or controlled view changes.
\end{sdefinition}

\begin{slemma}[Three-dimensional embedding of finite typed multigraphs]
\label{arch:lem-r3-embedding}
\Thm{} Every finite typed multigraph admits a topology-faithful typed rendering in
$\mathbb R^3$.
\end{slemma}

\begin{proof}
Choose pairwise distinct vertex locations in the plane $z=0\subset\mathbb R^3$ in
general position, and around each choose a disjoint small closed ball. Reserve a
distinct port on the ball boundary for each incident edge. Enumerate the finite edge
set and assign distinct positive heights. Each non-loop edge is drawn from its
source port vertically to its assigned height, horizontally to its target port, and
vertically down. Distinct heights prevent horizontal-edge intersections; general
position and disjoint port balls prevent unwanted intersections elsewhere. Loops
use two distinct ports and a short detour at their own height; parallel edges use
distinct ports and heights. The arcs meet only at declared common endpoints. A
finite injective legend preserves the finite active type set. Thus the graph embeds
as a typed one-dimensional complex in $\mathbb R^3$.
\end{proof}

\begin{scorollary}[Every finite P-space serialization has an $\mathbb R^3$ rendering]
\label{arch:cor-r3-pspace}
\Thm{} Every finite P-space serialization admits a topology-faithful,
label-preserving rendering in $\mathbb R^3$.
\end{scorollary}

\begin{proof}
A finite P-space serialization is a finite typed multigraph by
Definition~\ref{def:pspace}. Apply Lemma~\ref{arch:lem-r3-embedding}.
\end{proof}

\begin{slemma}[Two dimensions are not universally sufficient]
\label{arch:lem-r2-insufficient}
\Thm{} No topology-faithful crossing-free rendering rule in $\mathbb R^2$ can
represent every finite graph.
\end{slemma}

\begin{proof}
The complete graph $K_5$ has $v=5$ vertices and $e=10$ edges. Every simple planar
graph with $v\geq3$ has $e\leq3v-6$ by Euler's formula; for $K_5$, $3v-6=9<10$.
Thus $K_5$ is not planar and no universal crossing-free $\mathbb R^2$ rule exists.
\end{proof}

\begin{stheorem}[Minimal universal Euclidean rendering dimension]
\label{arch:thm-minimal-r3}
\Thm{} Three is the least Euclidean dimension in which every finite typed P-space
graph can be rendered topology-faithfully without unintended edge intersections:
every such graph is renderable in $\mathbb R^3$, while no universal crossing-free
rendering rule for arbitrary finite typed graphs exists in $\mathbb R^2$.
\end{stheorem}

\begin{proof}
The positive result is Corollary~\ref{arch:cor-r3-pspace}; the negative result is
Lemma~\ref{arch:lem-r2-insufficient}.\end{proof}

\begin{sproposition}[Scope of the rendering theorem]
\label{arch:prop-rendering-scope}
\Defn{} Theorem~\ref{arch:thm-minimal-r3} proves a universal graph-rendering fact.
It does not assert that all physical geometry is Euclidean, that a P-space metric
is the displayed Euclidean metric, or that a two-dimensional display cannot render
a particular planar selected subgraph. It establishes only that a representation
architecture intended to render arbitrary finite typed graphs without topological
ambiguity needs at least three Euclidean spatial dimensions.
\end{sproposition}

\subsection{Multiscale before/after visual audit}

\begin{sdefinition}[Interaction-indexed audit selector]
\label{arch:def-selector}
\Defn{} Let $i$ denote a declared interaction class or question, such as strong,
weak, electromagnetic, gravitational, or a composed interaction. At scale
$\delta$, an interaction-indexed audit selector is a typed map
\begin{equation}
S_{i,\delta}:\Pspace_{\mathcal E;R,\delta}\longrightarrow
\Pspace^{\mathrm{sel}}_{i,\delta}
\label{arch:eq-selector}
\end{equation}
that retains exactly the vertices, edges, types, and metric/boundary data required
to answer $\Audit_{i,\delta}$. It may retain a local subgraph, an extended-support
subgraph, or a whole-region subgraph depending on the interaction and question. It
need not retain every distinction in the full P-space graph simultaneously.
\end{sdefinition}

\begin{sdefinition}[Before/after transition record and visual certificate]
\label{arch:def-before-after}
\Defn{} For an admissible operation or finite composition $\gamma:x\to y$ at
$(R,\delta)$, define
\begin{equation}
\Tr_{i,\delta}(\gamma)=
\left(S_{i,\delta}(\Pspace_x),S_{i,\delta}(\Pspace_y),
S_{i,\delta}(\gamma),\Delta_{i,\delta}\right),
\label{arch:eq-transition-record}
\end{equation}
where $\Delta_{i,\delta}$ records typed additions, deletions, or modifications of
states, edges, labels, boundaries, support signatures, or declared metric values.
A visual transition certificate is
\begin{equation}
\View^3_{i,\delta}(\gamma)=
\left(\Emb^3(S_{i,\delta}(\Pspace_x)),\Emb^3(S_{i,\delta}(\Pspace_y));
\Emb^3(S_{i,\delta}(\gamma)),\Delta_{i,\delta},\mathcal L_{i,\delta}\right),
\label{arch:eq-visual-certificate}
\end{equation}
where $\mathcal L_{i,\delta}$ is a finite legend and audit instruction set. The
before/after views may share stable placement for retained vertices.
\end{sdefinition}

\begin{stheorem}[Existence of multiscale before/after $\mathbb R^3$ certificates]
\label{arch:thm-before-after}
\Thm{} For every finite selected interaction audit $\Audit_{i,\delta}$ and every
admissible finite transition $\gamma$, a visual transition certificate exists.
\end{stheorem}

\begin{proof}
The selected before-state, after-state, and transition graphs are finite typed
multigraphs. Lemma~\ref{arch:lem-r3-embedding} supplies topology-faithful
embeddings. Retained vertices may be fixed at shared positions across before/after
views, while new vertices and edges use unused ports and heights. Attaching the
finite change record and legend yields the certificate.
\end{proof}

\begin{sproposition}[Audit is multiscale, not single-view]
\label{arch:prop-multiscale}
\Defn{} A strong-interaction audit may select local support at one scale; a weak
interaction audit may select another scale and ordered history; an electromagnetic
audit may select a local field/interface structure; and a gravitational audit may
select extended support, boundary, or global composition. The full P-space graph
remains available, but each audit uses only the finite selected structure sufficient
for its own question.
\end{sproposition}

\begin{sdefinition}[Human visual verifier model]
\label{arch:def-human-verifier}
\Cnd{} A human visual verifier at budget $b$ is a protocol-conforming observer able,
within a declared finite time and attention budget, to distinguish the finite active
visual token set in a declared legend; follow finite edge/path traces in a
navigable $\mathbb R^3$ rendering; compare matched before/after vertices, edges,
and labels; execute the finite local checks in $\mathcal L_{i,\delta}$; and report
failure when a required token, path, type, boundary, regime, or declared change is
absent or mismatched. This is a verification-capability premise, not a theorem
about every actual human under every workload.
\end{sdefinition}

\begin{stheorem}[Conditional human-audit reliability]
\label{arch:thm-human-audit}
\Cnd{} Suppose: (i) the audit family $\Audit_{i,\delta}$ factors through the
selected graph $S_{i,\delta}(\Pspace_{\mathcal E;R,\delta})$; (ii) the selected
graph is rendered by a visual transition certificate; (iii) the certificate legend
and local checks are injective for the active type set; and (iv) the auditor
satisfies Definition~\ref{arch:def-human-verifier}. Then the certificate supplies
a sound and complete human-executable audit protocol for $\Audit_{i,\delta}$.
\end{stheorem}

\begin{proof}
By (i), the selected graph determines all values of the audit family. By
Theorem~\ref{arch:thm-before-after}, it and its change record have a
topology-faithful rendering. Assumption (iii) prevents visual conflation of active
types or checks; (iv) guarantees execution and failure reporting. Hence every
positive certificate corresponds to the required graph relations and every missing
or mismatched relation triggers protocol failure. The theorem supplies no
diagnosis of the present research community; it establishes only the stated
conditional relation between a finite certificate and a declared audit protocol.
\end{proof}

\begin{sproposition}[Why direct higher-dimensional human visualization is not required]
\label{arch:prop-highdim}
\Defn{} The raw P-space schema may carry arbitrarily rich or high-dimensional
physical data through labels, metrics, refinement maps, and typed relations. The
architecture does not require a human to directly perceive an intrinsically
higher-dimensional spatial object. The selected finite graph is rendered in
$\mathbb R^3$; additional variables are represented by finite typed tokens,
layered views, and scale changes.
\end{sproposition}

\subsection{Necessity up to invariant-preserving equivalence in novel domains}

\begin{sdefinition}[Novel physical extension]
\label{arch:def-novel-extension}
\Defn{} A novel physical extension is a new bounded-region theory object or
proposed realization that introduces at least one new state type, operation type,
boundary/interface role, regime, support relation, or cross-theory bridge while
remaining subject to the declared admissibility and typing discipline. Novelty is
structural, not a predicate of whether an evaluator has previously encountered the
verbal description.
\end{sdefinition}

\begin{sdefinition}[P-space/SOP audit invariant and equivalent representation]
\label{arch:def-equivalence}
\Defn{} For a novel-domain audit family $\Audit$, the P-space/SOP audit invariant is
the equivalence class of distinctions required to prevent records with different
audit outcomes from being identified. It includes state, operation, carrier,
boundary, interface, regime, support, ordered history, and P-space/A(R)-space
bridge distinctions whenever the audit depends on them. A representation
$r:X\to Z$ is P-space/SOP-equivalent for $\Audit$ when it is sound and complete
for $\Audit$ and
\begin{equation}
\ker(r)\subseteq\ker(q_{\Audit}).
\label{arch:eq-equivalence}
\end{equation}
It may use graphs, equations, data structures, simulation states, proof objects, or
other notation. Literal use of the symbols ``P-space'' and ``SOP'' is not required.
\end{sdefinition}

\begin{stheorem}[Fiber necessity]
\label{arch:thm-fiber}
\Thm{} Let $r:X\to Z$ be a proposed representation for $\Audit$. If there are
$x,y\in X$ with $r(x)=r(y)$ and some $a_j\in\Audit$ for which $a_j(x)\neq a_j(y)$,
then no decision procedure that factors only through $r$ is both sound and complete
for $\Audit$.
\end{stheorem}

\begin{proof}
Any decision procedure factoring through $r$ gives the same result on $x$ and $y$.
A sound-and-complete procedure must differ whenever a required audit value differs.
This is a contradiction.
\end{proof}

\begin{stheorem}[Necessity of invariant preservation for reliable novel-domain audit]
\label{arch:thm-necessity}
\Cnd{} Let $\mathcal N$ be a family of novel physical extensions, each with a
finite audit family on a finite-resolution P-space record. Any representation
architecture that reliably audits every member of $\mathcal N$ must be
P-space/SOP-equivalent for every such audit family. In particular, it must preserve
every state, operation, boundary, regime, support, history, and bridge distinction
on which an audit answer changes.
\end{stheorem}

\begin{proof}
If the architecture erases a required invariant, there are records differing on an
audit answer but identified by the representation. Theorem~\ref{arch:thm-fiber}
rules out sound-and-complete audit. Theorem~\ref{arch:thm-compression} identifies
the requisite lower bound on distinctions; Definition~\ref{arch:def-equivalence}
names the resulting condition P-space/SOP equivalence.
\end{proof}

\begin{sproposition}[Necessity is architectural, not notational]
\label{arch:prop-notational}
\Defn{} A literal graph file, diagram syntax, or the names P-space and SOP are not
uniquely necessary. Reversible re-encoding is acceptable. What is necessary is
preservation of the output-relevant typed invariant architecture. An alternative
can defeat this necessity claim only by remaining sound and complete without
identifying a pair whose audit outcomes differ.
\end{sproposition}

\begin{stheorem}[Full P-space/SOP architectural theorem]
\label{arch:thm-full}
\Cnd{} Under Hypothesis~\ref{arch:hyp-generator},
Axioms~\ref{arch:ax-closure}--\ref{arch:ax-scale}, the typed P-space/A(R)-space
bridge discipline of Definitions~\ref{def:arspace} and \ref{def:sop-record-v9},
and the human visual-verifier model of Definition~\ref{arch:def-human-verifier},
the following jointly hold:
\begin{enumerate}[label=\textup{(F\arabic*)},leftmargin=1.45cm]
\item Every admissible bounded-region functional state and finite transition has an
all-scale finite-resolution P-space representation.
\item P-space paths represent physical compositions only when source, target,
operator, boundary, interface, and regime types compose; untyped substitution is
not a physical transition.
\item For every finite audit family, the audit quotient is the maximally compressed
exact semantic representation; every sound-and-complete alternative preserves an
extensionally equivalent quotient.
\item Every finite selected P-space graph and before/after transition record has a
topology-faithful, label-preserving rendering in $\mathbb R^3$.
\item $\mathbb R^3$ is the least Euclidean dimension sufficient for a universal
topology-faithful rendering rule for arbitrary finite typed P-space graphs.
\item For each declared interaction class and scale, a coherent before/after
certificate exists and need preserve only the distinctions required by that audit.
\item A claim is human-audited at a declared finite budget only when it supplies a
budgeted certificate for its stated audit family; otherwise the appropriate status
is unaudited at that budget, not settled.
\item Whenever the certificate's finite vocabulary and local checks lie within the
verifier budget, it supports a sound-and-complete human-executable audit of its
declared question.
\item A theory representation cannot reliably audit a novel physical domain unless
it preserves P-space/SOP-equivalent output-relevant invariants for every audit it
claims to decide.
\end{enumerate}
\end{stheorem}

\begin{proof}
(F1) is Theorem~\ref{arch:thm-representation}. (F2) follows from
Definition~\ref{def:pspace}, Axiom~\ref{arch:ax-closure}, and the typed bridge
record in Definition~\ref{def:arspace}. (F3) is Theorem~\ref{arch:thm-compression}.
(F4)--(F6) are Corollary~\ref{arch:cor-r3-pspace},
Theorems~\ref{arch:thm-minimal-r3} and \ref{arch:thm-before-after}. (F7) is
Proposition~\ref{arch:prop-budgeted-need}; (F8) is
Theorem~\ref{arch:thm-human-audit}; and (F9) is
Theorem~\ref{arch:thm-necessity}.
\end{proof}

\begin{sproposition}[Conditional global coherence of the architectural graph]
\label{arch:prop-global-coherence}
\Cnd{} The architectural hypothesis graph is globally coherent under the declared
assessment discipline: every cross-space transfer is mediated by a typed bridge;
every theorem has local premises; universal claims are finite-resolution or
schema-level; all-scale maps preserve the support and boundary relations required
by the declared audit; certificates compile active relations into bounded executable
checks; and the exact factorization condition identifies which distinctions must
travel through every representation change. The architecture thereby supplies the
scale and verifier stages of the constraint-transport sequence in
Equation~\eqref{eq:joint-transport}. Physical commitments not established by the
definitions and proofs remain hypotheses or open obligations. This proposition is
an in-budget coherence claim in the sense of
Definition~\ref{def:assessment-budget}, not empirical validation of a particular
physical generator or user interface.
\end{sproposition}



%=================================================================
\section{Conditional temporal trace-quotient and constraint inheritance hypothesis}
\label{supp:inheritance}

\paragraph{Claim status.}
This section is a \emph{conditional global-coherence hypothesis record}. It is
not an empirical claim that an earlier civilization existed, that any earlier
population encountered the worked-case hazard, or that an absence of trace is
positive evidence for a predecessor. Its object is narrower: to identify when the
absence of a surviving material or historical predecessor is not an adequate
observation channel for deciding whether a declared hazard class was previously
encountered. The appropriate verdict is therefore
$\mathsf{GC}_{B}(\mathfrak H)$ for a declared history hypothesis and finite
assessment budget $B$, with physical realization, trace completeness, and
historical instantiation retained as explicit residue. It is not
$\mathsf{EmpVal}(\mathfrak H)$.

The same typed distinction that separates a local observation from a composition
predicate separates a present-day trace record from a complete history of use.
In the recursive calculus, a trace quotient is a proposed $\LOp$-program and is
licensed only by its preservation-and-detection certificate. Constraint inheritance
is a $\Register\circ\Recall$ transport of a reconstructible safety case, not a
transfer of an untyped procedure, artifact, or warning.

\begin{sdefinition}[Narrow-custody capability history]
\label{def:narrow-custody-history}
\Defn{} A \emph{narrow-custody capability history} is a finite or countably
infinite record
\begin{equation}
\mathfrak H=
\left(
\mathfrak C,\mathcal A,\{u_k\}_{k\geq 1},\{K_k\}_{k\geq 1},
\{M_k\}_{k\geq 1},\mathcal T_D,\mathcal G
\right),
\label{eq:narrow-custody-history}
\end{equation}
where:
\begin{enumerate}
\item $\mathfrak C$ is a typed composition instance
(Definition~\ref{def:typed-composition-instance});

\item $\mathcal A$ is a declared physical artifact or procedure with a
persistent P-space footprint $\mathsf{Foot}_P(\mathcal A)$;

\item $u_k$ is the $k$th use attempt, including its declared support,
ordered preparation path, and terminal branch;

\item $K_k$ is the set of operators able to execute the operational procedure at
use $k$; $M_k$ is the set of operators able to reconstruct the complete typed
constraint certificate at use $k$;

\item $\mathcal T_D$ is a declared deep-time trace-and-detection channel, mapping
a history to the material, archival, geochemical, stratigraphic, or other records
that a specified later observer can in principle recover; and

\item $\mathcal G$ is the regeneration predicate: it holds only when the complete
constraint certificate can be reconstructed, corrected, and transmitted beyond a
single narrow custodian lineage or institution.
\end{enumerate}

Operational competence and safety-model competence are distinct. In particular,
$K_k\neq\varnothing$ does not imply $M_k\neq\varnothing$, and $M_k\neq\varnothing$
does not imply $\mathcal G=1$. A group can retain a procedure without retaining
the bridge conditions, regime limits, support thresholds, retention conditions,
safe witnesses, output verification rules, and stopping conditions needed to
control its hazard class.
\end{sdefinition}

\begin{sdefinition}[Low-footprint, repeated-use, narrow-custody condition]
\label{def:low-footprint-condition}
\Defn{} A narrow-custody capability history satisfies the \emph{low-footprint,
repeated-use, narrow-custody condition}, written
\begin{equation}
\mathsf{LRN}(\mathfrak H)=1,
\label{eq:lrn-condition}
\end{equation}
when all of the following hold:
\begin{enumerate}
\item $\mathsf{Foot}_P(\mathcal A)$ is sufficiently small, persistent, or
reusable that use count is not bounded by a proportionate growth in durable
material footprint;

\item the attempted-use horizon is unbounded or long relative to the survival
and transmission horizon, so that $u_1,u_2,\ldots$ can be repeatedly composed by
only a few operators;

\item the complete certificate is in narrow custody:
\begin{equation}
M_k\subseteq K_k,
\qquad
|M_k|\ll |K_k|\ \text{or}\ \mathcal G=0;
\label{eq:narrow-certificate-custody}
\end{equation}

\item the history contains no verified bridge from use, artifact, or failure to a
later detectable record under $\mathcal T_D$.
\end{enumerate}

The definition does not assign numerical thresholds to ``small'', ``long'', or
``narrow''. Those thresholds are application-specific and must be declared before
the condition is used. The structural point is that a capability's causal leverage,
its persistent physical footprint, and the propagation of its complete safety model
are independent coordinates.
\end{sdefinition}

\begin{sproposition}[Cumulative exposure under narrow custody]
\label{prop:cumulative-exposure}
\Cnd{} Let $p_k$ be the conditional probability that $u_k$ enters a terminal
unresolved branch, conditional on survival through $u_1,\ldots,u_{k-1}$ and on
execution of $u_k$. Then
\begin{equation}
\Pr(\mathrm{survival\ through\ }N)=
\prod_{k=1}^{N}(1-p_k).
\label{eq:conditional-cumulative-survival}
\end{equation}
If $0\leq p_k<1$ and
\begin{equation}
\sum_{k=1}^{\infty}p_k=\infty,
\label{eq:divergent-conditional-risk}
\end{equation}
then
\begin{equation}
\lim_{N\to\infty}
\Pr(\mathrm{survival\ through\ }N)=0.
\label{eq:zero-asymptotic-survival}
\end{equation}
\textbf{Type:} \Cnd{} on the declared conditional probabilities. The proposition
neither estimates any $p_k$ nor asserts that the worked-case path, amplification,
or terminal branch refers to a physical capability.
\end{sproposition}

\begin{proof}
The chain rule for conditional survival gives
Equation~\eqref{eq:conditional-cumulative-survival}. Since
$\log(1-p_k)\leq -p_k$ for $0\leq p_k<1$,
\begin{equation}
\log\Pr(\mathrm{survival\ through\ }N)
=
\sum_{k=1}^{N}\log(1-p_k)
\leq
-\sum_{k=1}^{N}p_k.
\end{equation}
Under Equation~\eqref{eq:divergent-conditional-risk}, the right-hand side
tends to $-\infty$, proving Equation~\eqref{eq:zero-asymptotic-survival}.
\end{proof}

The relevance of Proposition~\ref{prop:cumulative-exposure} is not that broad
diffusion is necessary for exposure. Its antecedent permits repeated use by a
small set of operators. A high capital-cost technology can sometimes be made safer
by limiting access to a few sites; the low-footprint condition removes that
inference. Narrow custody can simultaneously limit scrutiny, prevent independent
model reconstruction, and allow cumulative use count to grow.

\begin{sdefinition}[Historical encounter predicate and trace quotient]
\label{def:historical-encounter}
\Defn{} Let
\begin{equation}
\mathsf{Encounter}(\mathfrak H)=1
\label{eq:encounter-predicate}
\end{equation}
if at least one use $u_k$ in the history enters the declared hazard-relevant
corridor or terminal branch, and let
\begin{equation}
\Lambda_D(\mathfrak H)=\mathcal T_D(\mathfrak H)
\label{eq:trace-quotient}
\end{equation}
be the recoverable trace record under the specified later detection channel $D$.
The statement ``no surviving precedent'' is the observation
$\Lambda_D(\mathfrak H)=\varnothing$ or, more generally, equality to a declared
null trace record. It is not identical to
$\mathsf{Encounter}(\mathfrak H)=0$ unless a trace-completeness bridge has been
proved.
\end{sdefinition}

\begin{sproposition}[Finite historical non-observability witness]
\label{prop:historical-nonobservability}
\Thm{} Let a declared history class contain two histories
$\mathfrak H_0$ and $\mathfrak H_1$ such that:
\begin{enumerate}
\item $\mathsf{Encounter}(\mathfrak H_0)=0$;
\item $\mathsf{Encounter}(\mathfrak H_1)=1$;
\item the complete operational and safety-model record of the encounter branch
in $\mathfrak H_1$ is lost across a declared discontinuity, so $\mathcal G=0$;
and
\item the declared trace channel maps both histories to the same later record:
\begin{equation}
\Lambda_D(\mathfrak H_0)=\Lambda_D(\mathfrak H_1).
\label{eq:trace-separation-equality}
\end{equation}
\end{enumerate}
Then no decision procedure factoring only through $\Lambda_D$ is sound and
complete for $\mathsf{Encounter}$ on a class containing both histories.
\end{sproposition}

\begin{proof}
For any putative trace-only decider
$D_{\mathrm{trace}}=F\circ\Lambda_D$,
Equation~\eqref{eq:trace-separation-equality} gives
\begin{equation}
D_{\mathrm{trace}}(\mathfrak H_0)=D_{\mathrm{trace}}(\mathfrak H_1).
\end{equation}
But the encounter predicate differs by conditions (1)--(2), contradicting
soundness and completeness.
\end{proof}

Proposition~\ref{prop:historical-nonobservability} is a direct analogue of
Proposition~\ref{prop:no-local-decider}. It proves a conditional information
statement, not a historical claim: when the declared trace channel identifies an
encounter and a no-encounter history, absence of a trace cannot decide the history.
A critic can defeat the antecedent by supplying a trace-completeness theorem,
a durable forensic signature, or another verified bridge showing that every
relevant encounter must remain detectable. The proposition requires that bridge;
it does not presume its absence.

\begin{sdefinition}[Constraint certificate]
\label{def:constraint-certificate}
\Defn{} A \emph{constraint certificate} for a candidate capability is a typed
record
\begin{equation}
\mathsf{Cert}(\mathfrak C)=
\left(
\Sigma_2(\mathfrak C),
\mathsf{Corr},
\mathsf{SafeWit},
\mathsf{Amp},
\mathsf{Contain},
\mathsf{Stop},
\mathsf{Output},
\mathsf{Trace}
\right),
\label{eq:constraint-certificate}
\end{equation}
which includes the second-order admissibility signature, corridor status, safe
non-traversal witnesses, amplification and containment conditions, reversal or
stopping conditions, declared invariant output observables, and the trace and
transmission assumptions under which historical precedent would be evidence.
A certificate is \emph{regeneratively propagated} only when a successor can
reconstruct and correct the certificate, including its unresolved branches, rather
than merely repeat an operational procedure.
\end{sdefinition}

\begin{sproposition}[Constraint inheritance requirement]
\label{prop:constraint-inheritance}
\Cnd{} For any history satisfying $\mathsf{LRN}(\mathfrak H)=1$, retention of
$\mathcal A$ or of a use protocol is not an adequate substitute for regenerative
propagation of $\mathsf{Cert}(\mathfrak C)$. In the absence of a verified
trace-completeness bridge, a later operator who receives only the artifact or the
procedure cannot infer from the absence of a warning that the corridor was
previously assessed, blocked, safe, or unencountered. The appropriate safety
object to transmit is therefore the typed certificate, with explicit status labels
for open, blocked, bridge-required, regime-invalid, support-incomplete, and
hazard-coupled branches.
\end{sproposition}

\subsection{Gravity-facing constraint transport across time}
\label{supp:constraint-transport-time}

The temporal trace and inheritance components complete the time axis of the
constraint record in Definition~\ref{def:typed-constraint-record}. Their function
is not to add historical narrative to a gravity result. They determine whether a
gravitational constraint remains the same typed constraint when it crosses a
historical discontinuity.

A missing trace, unresolved bridge, unavailable safe witness, or failed quotient
factorization is not by itself a corridor verdict. It is a typed residue with a
declared status and permitted next operation: refine the representation, supply
the relevant bridge or safe witness, restrict the question, construct a defeating
witness, or retain the branch as unresolved.

\begin{sproposition}[Temporal trace as a calibrated negative-inference operator]
\label{prop:trace-calibration}
\Defn{}/\Cnd{} The trace quotient
\begin{equation}
\Lambda_D:\mathcal H\longrightarrow\mathcal O
\end{equation}
acts as a calibrated negative-inference operator. A null trace supports the
negative conclusion $\mathsf{Encounter}=0$ only through a declared bridge
\begin{equation}
\mathsf B_{\tau}:
\mathcal H\rightsquigarrow\mathcal O
\end{equation}
that specifies preservation, recovery, and recognition conditions for the relevant
encounter class. Where the bridge identifies an encounter history and a
no-encounter history, the null trace carries no exclusion force for that predicate.
The operator therefore makes the trace-preservation bridge an explicit component of
gravity-risk assessment rather than an implicit historical assumption.
\end{sproposition}

\begin{sproposition}[Constraint inheritance as safety-case transport]
\label{prop:inheritance-transport}
\Cnd{} Let
\begin{equation}
\mathsf{Prop}_{k\to k+1}:
\mathsf{Con}_{k}\longmapsto\mathsf{Con}_{k+1}
\end{equation}
be a proposed successor transfer of a gravity-facing constraint record. The
transfer preserves the safety case only when every bridge, regime, support,
history, registration, certificate, and unresolved-branch distinction on which the
declared audit family changes is reconstructible in $\mathsf{Con}_{k+1}$. Equivalently,
for the declared audit invariant $\mathsf{Inv}_{\Audit}$,
\begin{equation}
\mathsf{Inv}_{\Audit}(\mathsf{Con}_{k+1})
=
\mathsf{Inv}_{\Audit}(\mathsf{Con}_{k}).
\label{eq:inheritance-invariant}
\end{equation}
The operation preserves the scope conditions of a gravitational constraint when the
relevant theory object, simulator, apparatus, institution, or operator changes; it
distinguishes reuse of a procedure from reconstruction of its safety argument; and
it carries unresolved physical branches forward as unresolved rather than silently
reclassifying them as safe or impossible. Its structural precedent is a checkable
certificate object rather than a mere provenance statement
\cite{S_Necula1997,S_MoreauEtAl2013}.
\end{sproposition}

\begin{sproposition}[Joint transport closure]
\label{prop:joint-transport-closure}
\Cnd{} The all-scale, audit, trace, and inheritance components implement the typed
transport sequence
\begin{equation}
\mathsf{Constraint}
\xrightarrow{\ \mathsf{Bridge}\ }
\mathsf{Theory\ object}
\xrightarrow{\ \mathsf{Scale}\ }
\mathsf{Selected\ representation}
\xrightarrow{\ \mathsf{Audit}\ }
\mathsf{Finite\ verifier}
\xrightarrow{\ \mathsf{Inheritance}\ }
\mathsf{Successor\ assessor}.
\label{eq:joint-transport}
\end{equation}
For a declared audit family, this composition determines whether a conclusion
derived, simulated, or registered in one typed setting remains the same constraint
in another, must be marked bridge-required, or fails because a transfer erases an
answer-changing distinction. The temporal corollary is thus a constraint on the
legitimacy of negative gravity-relevant inference, and the inheritance condition is
the mechanism by which the complete safety case remains reconstructible,
correctable, and auditable over time.
\end{sproposition}

%=================================================================
\section{Conditional scalar T-space hypothesis: a non-operative future extension}
\label{supp:tspace}

\paragraph{Status, role, and claim boundary.}
\Cnd{} This section records a separate conditional global-coherence hypothesis
for a componentwise scalar T-space. It does not enlarge the carrier of the main
strict-separation witness, alter the corridor predicate, enter either principal
proposition, or add a premise to any Propagation Obligation. It does not posit an
external universal clock, an observer-independent global simultaneity relation, a
new fundamental field, or an empirically selected temporal ontology. Its purpose
is narrower: to test whether an auditable scalar temporal representation can be
constructed from declared material duration realizers and declared physical
comparison support, with every point at which comparison fails retained as an
explicit undefinedness condition.

The intended audit role parallels the role of P-space. A physical claim becomes
auditable when its carrier, operator, regime, boundary, and admissible composition
are inspectable. A temporal claim should likewise expose the material process that
carries the claimed duration, the scalar state advance that is counted, the
physical support that makes a comparison legal, the registration bridge from a
clock record to a physical claim, and the condition under which no common T-space
is defined. No temporal relation is inferred merely from a shared background
coordinate or a timestamp in A(R)-space.

\paragraph{Piecewise theoretical input.}
Let
\begin{equation}
\mathcal E=\bigoplus_{\alpha\in A}\mathcal E_\alpha
\label{eq:tspace-piecewise-energy-action}
\end{equation}
denote the declared piecewise composition of energy/action functionals,
effective descriptions, material-clock models, state descriptions, and boundary
conditions admitted for the P-space domain. As in the P-space construction, the
symbol $\bigoplus$ records typed branch composition, not an assertion that every
branch is simultaneously active, numerically additive, or reducible to one
fundamental equation. A T-space candidate is therefore derived from the same
piecewise physical specification that types the underlying P-space, rather than
introduced as a free-standing coordinate line.

For an admissible temporal edge $e$, let $\mathcal J(e)$ be the active set of
well-typed duration realizers. Every active realizer must either emit a scalar
candidate duration in a declared common unit or supply a declared calibration
bridge to one. The construction uses the following nonempty realization classes.

\begin{scard}[Invariant proper time and calibrated physical clocks]
\label{card:tspace-proper-clock}
\Imp{} \textbf{Imported:} for a timelike material history in a Lorentzian regime,
general relativity supplies the invariant scalar proper-time increment
\begin{equation}
d\tau_{\mathrm{GR}}
=
\frac{1}{c}\sqrt{-g_{\mu\nu}\,dx^\mu dx^\nu}.
\label{eq:tspace-proper-time}
\end{equation}
Calibrated physical clocks supply operational scalar duration records by comparing
a declared clock phase $\phi$ to a calibrated transition frequency $\omega_0$,
so that $d\tau_{\mathrm{clk}}=d\phi/\omega_0$ within the declared clock model.
\textbf{Regime:} timelike histories and the stated relativistic and clock
calibration domains. \textbf{Use here:} establishes a nonempty, physically
auditable class of scalar duration realizers. \textbf{Not taken to establish:} a
global simultaneity convention, a universal scalar coordinate, or a temporal
relation between arbitrary disconnected regions
\cite{S_Einstein1916,S_Chou2010}.
\end{scard}

\begin{scard}[Material-reference scalar clocks]
\label{card:tspace-material-reference}
\Imp{} \textbf{Imported:} the Brown--Kucha\v{r} dust construction gives a material
reference system in which dust proper time is a canonical time variable along the
declared dust congruence. \textbf{Regime:} the specified canonical gravity model
with its dust reference field. \textbf{Use here:} demonstrates that a scalar
material time reference can be a typed physical branch rather than an externally
imposed parameter. \textbf{Not taken to establish:} an absolute temporal metric;
the construction itself introduces a privileged dynamical reference frame and time
foliation, so it is admissible here only on its declared material support and with
its branch status retained \cite{S_BrownKuchar1995}.
\end{scard}

\begin{scard}[State-dependent thermal flow]
\label{card:tspace-thermal-flow}
\Imp{} \textbf{Imported:} the thermal-time hypothesis associates a one-parameter
modular flow with a statistical state $\rho$ of an algebra of observables. In a
formal representation,
\begin{equation}
\alpha_s^{\rho}(A)=\rho^{is}A\rho^{-is}.
\label{eq:tspace-thermal-flow}
\end{equation}
\textbf{Regime:} the relevant algebraic/statistical state description.
\textbf{Use here:} supplies a nonempty state-dependent local flow branch that can
parameterize material evolution. \textbf{Not taken to establish:} that $s$ is an
elapsed duration in a common physical unit. It contributes to a T-space duration
only under an explicit calibration bridge
\begin{equation}
\mathsf{Bridge}_{\mathrm{th}\to\mathrm{clk}}:
\kappa_\rho\,ds\longmapsto d\tau_{\mathrm{clk}}.
\label{eq:tspace-thermal-calibration}
\end{equation}
Without that bridge, thermal flow remains an endogenous state parametrization and
not a silently substituted metric duration \cite{S_ConnesRovelli1994}.
\end{scard}

\begin{scard}[Discrete causal-order duration candidate]
\label{card:tspace-causal-set}
\Imp{} \textbf{Imported:} causal-set theory represents spacetime by a locally
finite partial causal order, with longest-chain constructions used as candidates
for timelike distance. For a declared causal-set branch, let $L(x,y)$ be the
length of a maximal causal chain between comparable elements and write
\begin{equation}
\Delta\tau_{\mathrm{cs}}(x,y)=\ell_{\mathcal E,\delta}\,L(x,y),
\label{eq:tspace-causal-set-duration}
\end{equation}
where $\ell_{\mathcal E,\delta}$ is an explicit conversion scale.
\textbf{Regime:} the declared discrete causal-order model and any stated recovery
bridge. \textbf{Use here:} contributes a discrete branch for partial temporal
comparability. \textbf{Not taken to establish:} that all physical duration is
causal-set chain length or that this branch is active in the worked case
\cite{S_BombelliEtAl1987}.
\end{scard}

\subsection{Temporal-support composition}
\label{supp:tspace-composition}

The T-space construction does \emph{not} multiply duration values emitted by
incompatible theories. Its composition is a typed conjunction of the predicates
that license transfer of a scalar duration through an edge. For a candidate edge
$e$ between strongly typed P-space states, define
\begin{equation}
\mathsf{Adm}_{T,\mathcal E,\delta}(e)
:=
\mathsf{Dur}_{\mathcal E}(e)
\cdot
\mathsf{FieldSupp}_{\mathcal E,\delta}(e)
\cdot
\mathsf{ClockSupp}_{\mathcal E,\delta}(e)
\cdot
\mathsf{RegBridge}_{\mathcal E,\delta}(e)
\cdot
\mathsf{ModelAgree}_{\mathcal E,\delta}(e)
\in\{0,1\}.
\label{eq:tspace-admissibility-product}
\end{equation}
Here the product is Boolean conjunction, not numerical physics. Its factors are:
\begin{enumerate}
\item $\mathsf{Dur}_{\mathcal E}(e)=1$ only if a declared duration realizer is
well-typed for the carrier, causal relation, regime, and boundary of $e$;
\item $\mathsf{FieldSupp}_{\mathcal E,\delta}(e)=1$ only if a declared physical
field, causal transfer, common material reference, or other typed support relation
licenses comparison at resolution $\delta$;
\item $\mathsf{ClockSupp}_{\mathcal E,\delta}(e)=1$ only if each endpoint contains
a declared duration-bearing material process, not merely an A(R)-space timestamp;
\item $\mathsf{RegBridge}_{\mathcal E,\delta}(e)=1$ only if every clock phase,
signal, detector record, or frequency comparison used in the inference is linked
to the physical duration claim by an explicit bridge with calibration and
uncertainty; and
\item $\mathsf{ModelAgree}_{\mathcal E,\delta}(e)=1$ only if the active duration
realizers agree on their overlap within declared approximation and uncertainty
bounds.
\end{enumerate}

More explicitly, if $D_j(e)\in\mathbb R_{\geq 0}$ is the duration emitted by
$j\in\mathcal J(e)$ after any necessary declared calibration, then
\begin{equation}
\mathsf{ModelAgree}_{\mathcal E,\delta}(e)=1
\quad\Longleftrightarrow\quad
\left|D_i(e)-D_j(e)\right|\leq\varepsilon_{ij}(e)
\qquad\forall i,j\in\mathcal J(e),
\label{eq:tspace-model-agreement}
\end{equation}
where $\varepsilon_{ij}(e)$ records the joint uncertainty and approximation
allowance. Disagreement is not averaged into apparent compatibility: it leaves
the edge bridge-required, regime-conflicted, or locally defeated.

When Equation~\eqref{eq:tspace-admissibility-product} is satisfied, the edge
receives the scalar duration class
\begin{equation}
w_T(e):=[D_j(e)]_{\varepsilon}\in\mathbb R_{\geq0};
\qquad
\mathsf{Adm}_{T,\mathcal E,\delta}(e)=0
\Longrightarrow
w_T(e)=\bot.
\label{eq:tspace-edge-duration}
\end{equation}
The symbol $\bot$ means undefined under the declared physical support conditions;
it does not mean zero, infinity, or an unmeasured finite value.

\begin{sdefinition}[Temporal-support graph and partial scalar T-space]
\label{def:tspace}
\Defn{} At declared specification $(\mathcal E,\delta)$, the \emph{temporal-support
graph} induced from $\Pspace_{\mathcal E,\delta}$ is
\begin{equation}
\mathcal G_{T,\mathcal E,\delta}
=(V_T,E_T,\tau_T,w_T),
\label{eq:tspace-support-graph}
\end{equation}
where $V_T$ contains only declared clock-bearing or duration-bearing material
states and
\begin{equation}
E_T=
\left\{e\in E_P:\mathsf{Adm}_{T,\mathcal E,\delta}(e)=1\right\}.
\label{eq:tspace-edges}
\end{equation}
Every edge retains its P-space carrier, operator, regime, boundary, field support,
clock model, registration bridge, active realizer set, uncertainty, and status.
The partial scalar T-space is the componentwise object
\begin{equation}
\Tspace_{\mathcal E,\delta}
:=
\bigsqcup_{C\in\pi_0(\mathcal G_{T,\mathcal E,\delta})}
\left(C,\Theta_C,\mu_T,\chi_T\right),
\label{eq:tspace-componentwise-object}
\end{equation}
only for components satisfying the scalarity condition below.
\end{sdefinition}

\subsection{Componentwise scalarity and the emergent local rate}
\label{supp:tspace-scalarity}

Let $\sim_0$ identify states connected by a declared zero-duration
synchronization edge. For a connected quotient component
$C\subseteq V_T/\sim_0$, a scalar coordinate
\begin{equation}
\Theta_C:C\longrightarrow\mathbb R
\label{eq:tspace-component-coordinate}
\end{equation}
is admitted only if declared comparison is path-independent.  Let $\ell$ be a
closed comparison loop in the underlying comparison graph, with $\sigma_e=+1$ or
$-1$ according to traversal orientation. The scalarity condition is
\begin{equation}
\sum_{e\in\ell}\sigma_e w_T(e)=0
\qquad
\text{for every admissible comparison loop }\ell.
\label{eq:tspace-integrability}
\end{equation}

If Equation~\eqref{eq:tspace-integrability} holds, select an arbitrary reference
class $\bar\nu_0\in C$ and define
\begin{equation}
\Theta_C(\bar\nu)
=
\sum_{e\in\gamma_{\bar\nu_0\to\bar\nu}}\sigma_e w_T(e).
\label{eq:tspace-path-coordinate}
\end{equation}
The result is independent of the admissible comparison path. It has only affine
reference freedom,
\begin{equation}
\Theta_C\mapsto a\Theta_C+b,
\qquad a>0,
\label{eq:tspace-affine-freedom}
\end{equation}
until a reference material clock and scale are declared. To make the induced
distance a metric rather than merely a pseudometric, identify
\begin{equation}
\bar\nu\sim_{\Theta}\bar\nu'
\quad\Longleftrightarrow\quad
\Theta_C(\bar\nu)=\Theta_C(\bar\nu'),
\qquad
\overline C:=C/\sim_{\Theta}.
\label{eq:tspace-coordinate-quotient}
\end{equation}
Within $\overline C$, the induced scalar temporal metric is
\begin{equation}
\mu_T([\bar\nu],[\bar\nu'])
=
\left|\Theta_C(\bar\nu)-\Theta_C(\bar\nu')\right|
\in[0,\infty).
\label{eq:tspace-scalar-metric}
\end{equation}
For states not contained in a common admissible temporal-support component,
\begin{equation}
\mu_T(\bar\nu,\bar\nu')=\bot.
\label{eq:tspace-cross-component-undefined}
\end{equation}
Thus lack of common intersecting field support is not silently represented as a
very large scalar distance; it blocks construction of a common T-space component.

The local emergent scalar is not the coordinate itself. Relative to a component
reference coordinate, define the \emph{chronometric compliance field}
\begin{equation}
\chi_T(\nu\mid\Theta_C)
:=
\frac{d\tau_\nu}{d\Theta_C}
\in(0,\infty),
\qquad
d\tau_\nu=\chi_T(\nu\mid\Theta_C)\,d\Theta_C.
\label{eq:tspace-chronometric-compliance}
\end{equation}
Operationally, where a calibrated physical clock is active,
\begin{equation}
\chi_T(\nu\mid\Theta_C)
=
\frac{d\phi_\nu/\omega_{0,\nu}}{d\Theta_C}.
\label{eq:tspace-operational-compliance}
\end{equation}
The value $\chi_T=1$ means that the declared material carrier advances at the
component reference rate; $\chi_T<1$ and $\chi_T>1$ record lower and higher local
accumulated duration per component reference increment. Under
Equation~\eqref{eq:tspace-affine-freedom}, $\chi_T\mapsto a^{-1}\chi_T$, while
the physical increment $d\tau_\nu$ remains invariant. Zero and infinity are not
ordinary physical values of this field: they indicate, respectively, a failure of
the clock-bearing state condition or an asymptotic/model-boundary condition.

\paragraph{Temporal inhomogeneity, not automatic intrinsic curvature.}
A one-dimensional scalar coordinate does not by itself carry nontrivial intrinsic
Riemann curvature. The present hypothesis therefore does not call arbitrary
variation in $\chi_T$ ``T-space curvature.'' A usable diagnostic is instead the
P-space inhomogeneity of the scalar field,
\begin{equation}
K_T(\nu):=
\left\|\nabla_P\log\chi_T(\nu)\right\|_{\mu_P},
\label{eq:tspace-inhomogeneity}
\end{equation}
where the relevant P-space derivative and metric are defined. Any claim linking
this diagnostic to gravitational curvature, a new field equation, or a novel
observable is residue for future work.

\subsection{Audit record, conditional coherence, and residue}
\label{supp:tspace-audit}

\begin{sdefinition}[Temporal audit record]
\label{def:tspace-audit-record}
\Defn{} A temporal relation claim between material states $\nu_x$ and $\nu_y$ is
represented by the typed audit record
\begin{equation}
\mathsf{TAudit}(\nu_x,\nu_y)=
\left(
\nu_x,\nu_y,C,\gamma,\mathcal J,\Theta_C,\chi_T,
\mathsf{FieldSupp},\mathsf{RegBridge},\varepsilon,\mathsf{status}
\right).
\label{eq:tspace-audit-record}
\end{equation}
The record identifies the material carriers, common component if any, admissible
comparison path, duration realizers, scalar coordinate, local rate field, physical
support, A(R)-space bridge, uncertainty, and status. A reader therefore need not
re-derive every underlying temporal theory in order to audit the finite question:
what changed, what counted the change, what physically related the counters, and
where does the license for comparison end?
\end{sdefinition}

\begin{sproposition}[Conditional componentwise scalar coherence]
\label{prop:tspace-coherence}
\Cnd{} Fix $(\mathcal E,\delta)$ and a declared P-space instance. Suppose that:
\begin{enumerate}
\item every active temporal realizer is defined on its declared typed domain;
\item every active overlap satisfies the agreement condition in
Equation~\eqref{eq:tspace-model-agreement};
\item every retained edge satisfies the duration, field-support, clock-support,
and registration-bridge conditions in
Equation~\eqref{eq:tspace-admissibility-product}; and
\item every candidate temporal-support component satisfies the integrability
condition in Equation~\eqref{eq:tspace-integrability}.
\end{enumerate}
Then each such component induces a well-defined scalar metric
$\mu_T$ by Equation~\eqref{eq:tspace-scalar-metric}, and all cross-component
relations remain undefined by Equation~\eqref{eq:tspace-cross-component-undefined}.
The hypothesis graph may receive the status
$\mathsf{GC}_B(\Tspace)=\mathsf{pass}$ only in the assessment-bounded sense of
Definition~\ref{def:assessment-budget}.
\end{sproposition}

\begin{proof}
The agreement condition makes every retained edge duration single-valued up to its
declared uncertainty quotient. The admissibility product excludes edges lacking a
typed duration, physical support, clock carrier, or registration bridge. The
integrability condition makes the path sum in
Equation~\eqref{eq:tspace-path-coordinate} independent of the selected admissible
comparison path, so it defines $\Theta_C$ on the quotient component. The coordinate quotient in
Equation~\eqref{eq:tspace-coordinate-quotient} identifies exactly the
zero-distance classes. The absolute coordinate difference is therefore
nonnegative, symmetric, and zero only on a quotient class, hence it defines
$\mu_T$ within that component. No
path and therefore no licensed scalar relation exists between disconnected
components; assigning one would introduce an undeclared bridge. The
assessment-bounded verdict follows only as the status defined in
Definition~\ref{def:assessment-budget}, not as empirical validation.
\end{proof}

\paragraph{Failure predicates and future-work residue.}
The scalar construction fails locally rather than being repaired by a hidden global
time whenever
\begin{equation}
\begin{array}{rcl}
\mathsf{Dur}_{\mathcal E}(e)&=&0,\quad\text{no duration realizer is typed for the edge;}\\
\mathsf{FieldSupp}_{\mathcal E,\delta}(e)&=&0,\quad\text{no declared physical comparison support exists;}\\
\mathsf{ClockSupp}_{\mathcal E,\delta}(e)&=&0,\quad\text{no material clock-bearing state is present;}\\
\mathsf{RegBridge}_{\mathcal E,\delta}(e)&=&0,\quad\text{a record cannot support the physical duration claim;}\\
\mathsf{ModelAgree}_{\mathcal E,\delta}(e)&=&0,\quad\text{active branches disagree beyond their declared allowance;}\\
\sum_{e\in\ell}\sigma_e w_T(e)&\neq&0,\quad\text{the presumed scalar component is non-integrable.}
\end{array}
\label{eq:tspace-failure-predicates}
\end{equation}

The following remain explicit residue: the physical law, if any, determining
$\chi_T$; the exact field-overlap conditions sufficient for common-component
membership; calibration among active realization branches; the extent and
connectedness of T-space components in physical regimes; whether the construction
yields predictions beyond known relativistic clock behaviour; and whether any
inhomogeneity diagnostic in Equation~\eqref{eq:tspace-inhomogeneity} has an
independent physical interpretation. These are future-work questions. They are
not imported into the present gravity carrier by the act of defining the hypothesis.

%=================================================================
\section{Propagation obligations and their failure predicates}
\label{supp:obligations}

Each obligation below is a derived side condition on a program in the recursive
constraint-transport calculus. The four primitive assessment actions remain
$\Register$, $\Recall$, $\GOp$, and $\LOp$; the obligations specify when a program
using them becomes invalid because it erases, mis-types, or fails to preserve an
answer-changing relation. Each is stated with the failure predicate under which a
layer lacking it produces an unsafe or invalid constraint transfer.

\begin{sobligation}[Type-safe transfer]\label{ob:transfer}
\Defn{} The layer represents the carrier, operator, projection, boundary, regime,
empirical-anchor, and registration roles of each theory-object and permits a
constraint to transfer only under a bridge preserving them. \textbf{Failure
predicate:} a constraint native to one object is transferred to another with a
mismatched or absent role mapping, producing either a false barrier (a constraint
imported that does not apply) or a false license (a constraint dropped that does
apply).
\end{sobligation}

\begin{sobligation}[Registration/physical separation]\label{ob:separation}
\Defn{} The layer keeps \Pspace{} and \ARspace{} distinct and bridges them
explicitly. \textbf{Failure predicate:} a bound on a registration/encoding
structure is substituted for a bound on a physical configuration, or conversely,
without a bridge---the specific error the worked case is built to expose, since its
constraint concerns reachability of a registration structure while its candidate
barrier (an energy condition) is native to the physical configuration.
\end{sobligation}

\begin{sobligation}[Composition admissibility]\label{ob:composition}
\Defn{} The layer evaluates constraints that are properties of a composition---path
admissibility, corridor existence, regime maintenance along a path---rather than of
any single step. \textbf{Failure predicate:} the layer evaluates only per-step
admissibility and reports a path admissible because every step is admissible, while
a global property of the composition (a separating barrier, a regime exit) makes the
path inadmissible. This is the capability whose absence
Proposition~\ref{prop:insufficiency-supp} identifies in the derive-and-test
paradigm. \textbf{Note:} corridor existence for the worked case---whether an
energy-condition, constraint-satisfaction, or encoding-positivity barrier separates
ordinary matter from the target---is the designated open question, and is exactly an
instance this obligation requires the layer to be able to pose.
\end{sobligation}

\begin{sobligation}[Ordered history and state retention]\label{ob:history}
\Defn{} The layer represents the ordered path, the re-linearization basepoint of
each increment, and the retention condition by which a prepared registration state
remains available to the next increment. In a controlled physical realization, a
necessary timing condition has the schematic form
\begin{equation}
\tau_{\mathrm{sense}}(x_k)+\tau_{\mathrm{control}}(x_k)+
\tau_{\mathrm{act}}(x_k)<\tau_{\mathrm{relax}}(x_k),
\label{eq:retention-cadence}
\end{equation}
with the relevant synchronization requirement included when support is distributed.
\textbf{Failure predicate:} the layer treats an unordered ledger of individually
admissible operations as a path even though the registered distinction relaxes,
is overwritten, or is applied in an order in which the composed operation does not
exist.
\end{sobligation}

\begin{sobligation}[Support and global-invariant tracking]\label{ob:support}
\Defn{} The layer represents the declared support $S$, its incidence or topology,
and the global invariant whose attainment is required by the bridge. It
explicitly distinguishes a full-support preparation from a proper-support proxy or
locally indistinguishable marginal. \textbf{Failure predicate:} the layer treats
local agreement as evidence that the global invariant has been attained, or treats
a sub-threshold support intervention as a decisive null test of a full-support
predicate. Global support is not acausality and does not by itself establish either
reachability or hazard.
\end{sobligation}

\begin{sobligation}[Regime tracking]\label{ob:regime}
\Defn{} The layer tracks, along a simulated path, whether each step remains within
the regime in which its imported couplings and safety properties hold, marking the
path bridge-required where it does not. \textbf{Failure predicate:} the layer
applies a coupling or safety property (the linear first law, non-signalling,
wedge-local confinement) outside its regime of validity, because the path has left
the linear neighbourhood without a local indicator and the layer did not track the
exit.
\end{sobligation}

\begin{sobligation}[Wedge-locality transfer]\label{ob:wedge}
\Defn{}/\Cnd{} The layer requires an explicit bridge before importing the
wedge-local reconstruction property of Card~\ref{card:qec} from holographic models
to a realistic target. \textbf{Failure predicate:} reconstruction in the realistic
setting is not wedge-local, so acting on the encoding support produces bulk effects
outside the corresponding causal wedge, yet the layer reports causal-wedge
confinement because it imported entanglement-wedge locality without the required
P-space bridge.
\end{sobligation}

\begin{sobligation}[Safe witness and output verification]\label{ob:safewitness}
\Defn{} Before classifying a composed route as validation--hazard coupled, the layer
searches for a safe non-traversal witness: a barrier certificate, full-support
invariant, distributed measurement, bounded surrogate, or other proof object that
decides the relevant predicate without traversing the candidate hazardous branch.
It also requires a declared gauge- or coordinate-invariant output observable for
any claimed terminal P-space change. \textbf{Failure predicate:} the layer either
converts the absence of a supplied certificate into a claim that traversal is the
only decision route in reality, or treats an untyped apparatus artifact as evidence
for a terminal geometric change.
\end{sobligation}

\begin{sobligation}[Hazard-coupled flagging]\label{ob:hazard}
\Defn{} The layer identifies constraints whose empirical validation is coupled to a
hazard (Proposition~\ref{prop:coupling-supp}) and requires simulation or a safe decision object before any proposed experiment. The flag is conjunctive: it requires an open
corridor, a reachable amplifying sector, and the absence of established
containment, reversal, stopping, or safe-witness conditions. \textbf{Failure
predicate:} the layer recommends, or fails to flag, an experiment whose success
condition is also its unresolved hazard---in the worked case, a traversal whose
demonstration of corridor existence is the lighting of a potentially uncontrolled
positive feedback on a globally supported geometric variable. \textbf{Associated
cost:} the existential-risk asymmetry is recorded only for the branch in which
these conditions jointly hold; global support alone is not treated as proof of
uncontrollability.
\end{sobligation}

\begin{sobligation}[Intergenerational constraint propagation]\label{ob:inheritance}
\Defn{} For a candidate capability inside a declared temporal history hypothesis,
the layer carries a regenerable typed constraint certificate across successor
operators and institutions whenever low persistent footprint, repeated use, narrow
custody, and no verified trace-completeness bridge are jointly possible.
\textbf{Failure predicate:} a reusable artifact, procedure, archive, or null trace
record is substituted for the complete certificate, allowing a later operator to
infer that a corridor was previously assessed, blocked, safe, or unencountered
without a preservation-and-detection bridge.
\end{sobligation}


\begin{sproposition}[Necessity of the second-order layer]\label{prop:necessity}
\Cnd{}/\Thm{} For the declared witness class, a first-order local-observation
quotient cannot discharge the obligations because it is an unlicensed $\LOp$
contraction: it erases at least one relation \emph{between} theory-objects or one
property \emph{along} a composition on which corridor status changes.
Propositions~\ref{prop:strict-separation-witness},
\ref{cor:strict-witness-unlicensed-contraction}, and
\ref{prop:second-order-nonreducibility} establish the strict and semantic forms
of this claim. A sound and complete decision must either retain a registration at
least as discriminating as the bridge, regime, history, and support signature, or
obtain a certificate that changes the relevant decision branch without making the
forbidden identification. A P-space/A(R)-space representation is one such
second-order realization; the result does not claim that its notation is uniquely
privileged.
\end{sproposition}

%=================================================================
\section{Scope, status, and the deciding question}
\label{supp:scope}

\paragraph{Status architecture.} The supplement's universal constructions are
schema-level and finite-resolution: they specify how typed registrations,
quotients, audit records, and verifier certificates compose within their declared
premises. The worked-case route, its corridor conditions, its hazard branch, the
topological interaction/registration bridge, temporal-history branch, and the
scalar T-space construction each remain in the
record with their stated \Cnd{}, \Hyp{}, \Open{}, bridge, regime, and empirical
reference conditions. The associated physical values are determined by the
certificates and observations that satisfy those conditions. The finite
strict-separation theorem, by contrast, is fixed by its stated witness class and
fiber criterion.

\paragraph{Novelty status.} Novelty is assessed operationally. A statement that
the present deployment is exhausted by a familiar antecedent theorem is a proposed
$\LOp$ contraction. It becomes valid at a declared budget only when it preserves
the novelty signature of Definition~\ref{def:novelty-residue}. The remaining
generated registrations are retained as $\NovRes_{B,\Theta}$ with explicit status.
The domain-general scope of $\GOp$ is schema-level: it accepts any registrable
payload and emits the registrations required by the declared assessment program.

\paragraph{The deciding question.} Corridor existence
(Obligation~\ref{ob:composition}) is the designated physical question for the
worked case: whether an in-regime, retained, full-support admissible path connects
ordinary matter to the target, or a gravity-native barrier separates them. The
article establishes the class of constraint and the obligations on a layer that
adjudicates it. A verified safe non-traversal witness changes the validation branch
without converting the global predicate into a local one.

\paragraph{First-order and second-order readings.} At the first-order layer, a
channel of CPTP operations, measurement, and classical feedback is an engineering
composition. At the second-order layer, the article specifies which gravitational
constraints transfer across theory-objects, when a local-observation quotient is
inadequate, which distinctions an alternative-gravity simulation layer must
represent for a sound and complete decision over the declared class, and how the
novelty of that typed deployment is assessed without erasing its generated
registrations. The supplementary audit infrastructure supplies finite-resolution,
quotienting, rendering, and certificate constructions for implementations of that
calculus. Together with the source-to-composition records, the instantiated four-role
certificate family and the gravitational leverage composition compose existing
physics objects, their derivations, and their empirical registrations through
explicit bridges into an inspectable status-labeled global hypothesis. The
observability factorization derives the paired $\GPair$ closure in which that
composition is assessed. At the declared assessment budget, the composed record
emits its global-coherence verdict; its productivity, realization, and empirical
records specify the next quantitative-instantiation tasks. The matter-observability
section identifies a corresponding class of physics hard problems whose energy-scaled
matter protocols remain inside one invariant quotient.

\begin{thebibliography}{99}

\bibitem[Dray and 't~Hooft(1985)]{S_DraytHooft1985}
T.~Dray and G.~'t~Hooft, The gravitational shock wave of a massless particle,
\emph{Nuclear Physics B} \textbf{253} (1985) 173--188.
\url{https://doi.org/10.1016/0550-3213(85)90525-5}.

\bibitem[Shenker and Stanford(2014a)]{S_ShenkerStanford2014}
S.~H.~Shenker and D.~Stanford, Black holes and the butterfly effect,
\emph{Journal of High Energy Physics} \textbf{03} (2014) 067.
\url{https://doi.org/10.1007/JHEP03(2014)067}.

\bibitem[Shenker and Stanford(2014b)]{S_ShenkerStanford2014Multiple}
S.~H.~Shenker and D.~Stanford, Multiple shocks,
\emph{Journal of High Energy Physics} \textbf{12} (2014) 046.
\url{https://doi.org/10.1007/JHEP12(2014)046}.

\bibitem[Jeans(1902)]{S_Jeans1902}
J.~H.~Jeans, The stability of a spherical nebula,
\emph{Philosophical Transactions of the Royal Society of London. Series A}
\textbf{199} (1902) 1--53.
\url{https://doi.org/10.1098/rsta.1902.0012}.

\bibitem[Binney and Tremaine(2008)]{S_BinneyTremaine2008}
J.~Binney and S.~Tremaine, \emph{Galactic Dynamics}, 2nd ed.,
Princeton University Press, Princeton, NJ (2008).

\bibitem[Chavanis(2010)]{S_Chavanis2010}
P.-H.~Chavanis, Dynamical stability of infinite homogeneous
self-gravitating systems: application of the Nyquist method,
arXiv:1002.0291 (2010).



\bibitem[Moreau et al.(2013)]{S_MoreauEtAl2013}
L.~Moreau, P.~Missier, and the PROV-DM Revision Team, \emph{PROV-DM: The PROV
Data Model}, W3C Recommendation, 30 April 2013.
\url{https://www.w3.org/TR/prov-dm/}.

\bibitem[Necula(1997)]{S_Necula1997}
G.~C.~Necula, Proof-carrying code, in \emph{Proceedings of the 24th ACM
SIGPLAN--SIGACT Symposium on Principles of Programming Languages}, Paris,
France, 1997, pp.~106--119.
\url{https://doi.org/10.1145/263699.263712}.

\bibitem[Wilson(1971)]{S_Wilson1971}
K.~G.~Wilson, Renormalization group and critical phenomena. I. Renormalization
group and the Kadanoff scaling picture, \emph{Physical Review B} \textbf{4}
(1971) 3174--3183.
\url{https://doi.org/10.1103/PhysRevB.4.3174}.

\bibitem[Einstein(1916)]{S_Einstein1916}
A.~Einstein, Die Grundlage der allgemeinen Relativit\"atstheorie,
\emph{Annalen der Physik} \textbf{354} (1916) 769--822.
\url{https://doi.org/10.1002/andp.19163540702}.

\bibitem[Chou et al.(2010)]{S_Chou2010}
C.~W.~Chou, D.~B.~Hume, T.~Rosenband, and D.~J.~Wineland, Optical clocks and
relativity, \emph{Science} \textbf{329} (2010) 1630--1633.
\url{https://doi.org/10.1126/science.1192720}.

\bibitem[Brown and Kucha\v{r}(1995)]{S_BrownKuchar1995}
J.~D.~Brown and K.~V.~Kucha\v{r}, Dust as a standard of space and time in
canonical quantum gravity, \emph{Physical Review D} \textbf{51} (1995)
5600--5629.
\url{https://doi.org/10.1103/PhysRevD.51.5600}.

\bibitem[Connes and Rovelli(1994)]{S_ConnesRovelli1994}
A.~Connes and C.~Rovelli, Von Neumann algebra automorphisms and
time-thermodynamics relation in generally covariant quantum theories,
\emph{Classical and Quantum Gravity} \textbf{11} (1994) 2899--2917.
\url{https://doi.org/10.1088/0264-9381/11/12/007}.

\bibitem[Bombelli et al.(1987)]{S_BombelliEtAl1987}
L.~Bombelli, J.~Lee, D.~Meyer, and R.~D.~Sorkin, Space-time as a causal set,
\emph{Physical Review Letters} \textbf{59} (1987) 521--524.
\url{https://doi.org/10.1103/PhysRevLett.59.521}.


\bibitem[Van~Raamsdonk(2010)]{S_VanRaamsdonk2010}
M.~Van~Raamsdonk, Building up spacetime with quantum entanglement,
\emph{Gen. Relativ. Gravit.} \textbf{42} (2010) 2323--2329.

\bibitem[Maldacena and Susskind(2013)]{S_MaldacenaSusskind2013}
J.~Maldacena, L.~Susskind, Cool horizons for entangled black holes,
\emph{Fortschr. Phys.} \textbf{61} (2013) 781--811.

\bibitem[Jacobson(2016)]{S_Jacobson2016}
T.~Jacobson, Entanglement equilibrium and the Einstein equation,
\emph{Phys. Rev. Lett.} \textbf{116} (2016) 201101.

\bibitem[Faulkner et al.(2014)]{S_FaulknerEtAl2014}
T.~Faulkner, M.~Guica, T.~Hartman, R.~C.~Myers, M.~Van~Raamsdonk, Gravitation from
entanglement in holographic CFTs, \emph{J. High Energy Phys.} \textbf{2014}(03)
(2014) 051.

\bibitem[Almheiri et al.(2015)]{S_AlmheiriDongHarlow2015}
A.~Almheiri, X.~Dong, D.~Harlow, Bulk locality and quantum error correction in
AdS/CFT, \emph{J. High Energy Phys.} \textbf{2015}(04) (2015) 163.

\bibitem[Pastawski et al.(2015)]{S_PastawskiEtAl2015}
F.~Pastawski, B.~Yoshida, D.~Harlow, J.~Preskill, Holographic quantum
error-correcting codes: toy models for the bulk/boundary correspondence,
\emph{J. High Energy Phys.} \textbf{2015}(06) (2015) 149.

\bibitem[Dong et al.(2016)]{S_DongHarlowWall2016}
X.~Dong, D.~Harlow, A.~C.~Wall, Reconstruction of bulk operators within the
entanglement wedge in gauge-gravity duality, \emph{Phys. Rev. Lett.} \textbf{117}
(2016) 021601.

\bibitem[Nielsen and Chuang(2010)]{S_NielsenChuang2010}
M.~A.~Nielsen, I.~L.~Chuang, \emph{Quantum Computation and Quantum Information},
Cambridge University Press, 2010 (no-communication theorem for local operations).

\bibitem[Gisin(1990)]{S_Gisin1990}
N.~Gisin, Weinberg's non-linear quantum mechanics and supraluminal communications,
\emph{Phys. Lett. A} \textbf{143} (1990) 1--2.

\bibitem[Polchinski(1991)]{S_Polchinski1991}
J.~Polchinski, Weinberg's nonlinear quantum mechanics and the
Einstein--Podolsky--Rosen paradox, \emph{Phys. Rev. Lett.} \textbf{66} (1991)
397--400.



\bibitem[Davies and Lewis(1970)]{S_DaviesLewis1970}
E.~B.~Davies and J.~T.~Lewis, An operational approach to quantum probability,
\emph{Communications in Mathematical Physics} \textbf{17} (1970) 239--260.
\url{https://doi.org/10.1007/BF01647093}.

\bibitem[Ozawa(1984)]{S_Ozawa1984}
M.~Ozawa, Quantum measuring processes of continuous observables,
\emph{Journal of Mathematical Physics} \textbf{25} (1984) 79--87.
\url{https://doi.org/10.1063/1.526000}.

\bibitem[Particle Data Group(2024)]{S_PDG2024}
S.~Navas et al. (Particle Data Group), Review of Particle Physics,
\emph{Physical Review D} \textbf{110} (2024) 030001.
\url{https://doi.org/10.1103/PhysRevD.110.030001}.

\bibitem[Will(2014)]{S_Will2014}
C.~M.~Will, The confrontation between general relativity and experiment,
\emph{Living Reviews in Relativity} \textbf{17} (2014) 4.
\url{https://doi.org/10.12942/lrr-2014-4}.

\bibitem[Berti et al.(2015)]{S_Berti2015}
E.~Berti et al., Testing general relativity with present and future astrophysical
observations, \emph{Classical and Quantum Gravity} \textbf{32} (2015) 243001.
\url{https://doi.org/10.1088/0264-9381/32/24/243001}.

\bibitem[Shamir(1979)]{S_Shamir1979}
A.~Shamir, How to share a secret,
\emph{Communications of the ACM} \textbf{22} (1979) 612--613.
\url{https://doi.org/10.1145/359168.359176}.

\bibitem[Jurdjevic and Sussmann(1972)]{S_JurdjevicSussmann1972}
V.~Jurdjevic and H.~J.~Sussmann, Control systems on Lie groups,
\emph{Journal of Differential Equations} \textbf{12} (1972) 313--329.
\url{https://doi.org/10.1016/0022-0396(72)90035-6}.


\bibitem[Maki et al.(1962)]{S_MakiNakagawaSakata1962}
Z.~Maki, M.~Nakagawa, and S.~Sakata, Remarks on the unified model of elementary
particles, \emph{Progress of Theoretical Physics} \textbf{28} (1962) 870--880.
\url{https://doi.org/10.1143/PTP.28.870}.

\bibitem[Wolfenstein(1978)]{S_Wolfenstein1978}
L.~Wolfenstein, Neutrino oscillations in matter, \emph{Physical Review D}
\textbf{17} (1978) 2369--2374.
\url{https://doi.org/10.1103/PhysRevD.17.2369}.

\bibitem[Arnowitt et al.(1962)]{S_ArnowittDeserMisner1962}
R.~Arnowitt, S.~Deser, and C.~W.~Misner, The dynamics of general relativity,
in \emph{Gravitation: An Introduction to Current Research}, ed. L.~Witten,
Wiley, New York, 1962, pp.~227--265.
\url{https://arxiv.org/abs/gr-qc/0405109}.


\bibitem[Terhal et al.(2001)]{S_TerhalDiVincenzoLeung2001}
B.~M.~Terhal, D.~P.~DiVincenzo, and D.~W.~Leung, Hiding bits in Bell states,
\emph{Physical Review Letters} \textbf{86} (2001) 5807--5810.
\url{https://doi.org/10.1103/PhysRevLett.86.5807}.

\bibitem[Bravyi et al.(2010)]{S_BravyiHastingsMichalakis2010}
S.~Bravyi, M.~B.~Hastings, and S.~Michalakis, Topological quantum order:
stability under local perturbations, \emph{Journal of Mathematical Physics}
\textbf{51} (2010) 093512.
\url{https://doi.org/10.1063/1.3490195}.

\bibitem[Liu et al.(2007)]{S_LiuChristandlVerstraete2007}
Y.-K.~Liu, M.~Christandl, and F.~Verstraete, Quantum computational complexity of
the $N$-representability problem: QMA complete, \emph{Physical Review Letters}
\textbf{98} (2007) 110503.
\url{https://doi.org/10.1103/PhysRevLett.98.110503}.

\bibitem[Giddings(2005)]{S_Giddings2005}
S.~B.~Giddings, Observables in effective gravity, \emph{Physical Review D}
\textbf{74} (2006) 106006.
\url{https://doi.org/10.1103/PhysRevD.74.106006}.

\bibitem[Khavkine(2015)]{S_Khavkine2015}
I.~Khavkine, Local and gauge invariant observables in gravity,
\emph{Classical and Quantum Gravity} \textbf{32} (2015) 185019.
\url{https://doi.org/10.1088/0264-9381/32/18/185019}.

\bibitem[Faulkner et al.(2017)]{S_FaulknerEtAl2017}
T.~Faulkner, F.~M.~Haehl, E.~Hijano, O.~Parrikar, C.~Rabideau, and
M.~Van~Raamsdonk, Nonlinear gravity from entanglement in conformal field theories,
\emph{Journal of High Energy Physics} \textbf{2017}(08) (2017) 057.
\url{https://doi.org/10.1007/JHEP08(2017)057}.

\bibitem[Jafferis et al.(2016)]{S_JafferisEtAl2016}
D.~L.~Jafferis, A.~Lewkowycz, J.~Maldacena, and S.~J.~Suh, Relative entropy equals
bulk relative entropy, \emph{Journal of High Energy Physics} \textbf{2016}(06)
(2016) 004.
\url{https://doi.org/10.1007/JHEP06(2016)004}.

\bibitem[Tarski(1955)]{S_Tarski1955}
A.~Tarski, A lattice-theoretical fixpoint theorem and its applications,
\emph{Pacific Journal of Mathematics} \textbf{5} (1955) 285--309.
\url{https://doi.org/10.2140/pjm.1955.5.285}.

\end{thebibliography}

\end{document}
